% Initiation à la publication assistée par ordinateur — Informatique 3ème — Cours signé OnBuch+
% Programme officiel MINESEC. Compiler avec : tectonic N07-initiation-publication-assistee-par-ordinateur.tex
\newcommand{\DOCMATIERE}{Informatique}
\newcommand{\DOCNIVEAU}{3ème}
\newcommand{\DOCMODULE}{Informatique – classe de 3ème}
\newcommand{\DOCLECON}{N07}
\newcommand{\DOCTITRE}{Initiation à la publication assistée par ordinateur}
% OnBuch+ — préambule « cours signés » ENRICHI (compilé avec tectonic / XeLaTeX)
% Système visuel « Pop » d'OnBuch+ : fond crème (jamais blanc), bordures encre
% épaisses, ombres « dures » décalées, titres Archivo Black en capitales.
% Version MATHÉMATIQUES : graphes (pgfplots/tikz), boîtes pédagogiques
% (définition, propriété, méthode, attention, le savais-tu, exercice résolu,
% exercice, TP, bilan), page de garde et signature OnBuch+ ancrée en bas.
%
% Macros attendues dans main.tex AVANT \input{preamble} :
%   \DOCMATIERE  \DOCNIVEAU  \DOCMODULE  \DOCLECON (numéro)  \DOCTITRE
\documentclass[11pt]{article}

\usepackage{fontspec}
\usepackage[french]{babel}
\usepackage[a4paper,margin=2cm,top=3.4cm,bottom=2.8cm,headheight=1.4cm,footskip=1.3cm]{geometry}
\usepackage{amsmath,amssymb,amsfonts}
\usepackage{mathtools} % \xmapsto, \coloneqq... souvent utilisés par les modèles
\usepackage{cancel} % simplifications barrées (bilans de forces)
% \abs{z} (module, valeur absolue) et \norm{u} (norme), utilisés par les modèles
\providecommand{\abs}{}\let\abs\relax\DeclarePairedDelimiter{\abs}{\lvert}{\rvert}
\providecommand{\norm}{}\let\norm\relax\DeclarePairedDelimiter{\norm}{\lVert}{\rVert}
\usepackage[table]{xcolor} % [table] : fournit aussi \rowcolors (alternance de lignes)
\usepackage{pagecolor}
\usepackage{tikz}
\usetikzlibrary{arrows.meta,positioning,calc,shapes.geometric,shapes.misc,shapes.arrows,shapes.symbols,decorations.pathmorphing,decorations.pathreplacing,decorations.markings,patterns,shadows,mindmap,trees,fit,backgrounds,matrix,intersections,angles,quotes}
\usepackage{pgfplots}
\usepackage[version=4]{mhchem} % équations nucléaires \ce{^{235}_{92}U}
\usepackage[european,siunitx]{circuitikz} % schémas électriques
\pgfplotsset{compat=1.18}
\usepgfplotslibrary{fillbetween,groupplots} % aires entre courbes (calcul intégral)
\usepackage{enumitem}
\usepackage{fancyhdr}
% [most] charge toutes les bibliothèques tcolorbox (listings, minted, raster,
% documentation...) et gonfle inutilement la mémoire TeX sur un document long
% (~300 boîtes) : on ne charge que ce qui est réellement utilisé.
\usepackage[skins,breakable]{tcolorbox}
\usepackage{graphicx}
\usepackage{colortbl}
\usepackage{booktabs}
\usepackage{tabularx}
\usepackage{diagbox} % case d'en-tête diagonale des tableaux à double entrée
% Colonnes extensibles centrée / gauche / droite (C, L, R), souvent utilisées par les modèles
\newcolumntype{C}{>{\centering\arraybackslash}X}
\newcolumntype{L}{>{\raggedright\arraybackslash}X}
\newcolumntype{R}{>{\raggedleft\arraybackslash}X}
\usepackage{array}
\usepackage{multirow}
\usepackage{multicol}
\usepackage{siunitx}
% parse-numbers=false : accepte aussi \SI{2,5 \times 10^{-3}}{...}
\sisetup{locale=FR,output-decimal-marker={,},group-minimum-digits=4,parse-numbers=false}
\DeclareSIUnit\ohm{\ensuremath{\symup{\Omega}}} % U+2126 absent de la police
\DeclareSIUnit\division{div} % réglages d'oscilloscope (ms/div, V/div)
\DeclareSIUnit\tr{tr} % tours (vitesse de rotation, ex. \SI{1500}{\tr\per\minute})
\DeclareSIUnit\bit{bit}
\DeclareSIUnit\byte{o} % octet
\DeclareSIUnit\inch{po} % pouce
\DeclareSIUnit\ppp{ppp} % points par pouce (résolution d'image)
\usepackage{qrcode}
% Blocs de code source (PHP, C, HTML, JavaScript, SQL) — spécifique à la
% matière Informatique : listings + habillage aux couleurs OnBuch+ (voir le
% style \lstset ci-dessous, à côté des autres boîtes pédagogiques).
\usepackage{listings}
% \degree : commande fréquemment utilisée par les modèles pour un angle ou une
% température en dehors de \SI{}{} (non fournie par les paquets chargés).
\providecommand{\degree}{^{\circ}}
% \qed (fin de démonstration), utilisé par les modèles sans amsthm
\providecommand{\qed}{\hfill$\square$}
% \barre : double barre des valeurs interdites dans un tableau de variations (tkz-tab)
\providecommand{\barre}{\ensuremath{\|}}
% environnement proof (amsthm non chargé) utilisé par les modèles
\ifdefined\proof\else\newenvironment{proof}[1][Démonstration]{\par\noindent\textit{#1.}\ }{\qed\par}\fi
\usepackage{lastpage}
\usepackage{hyperref}
\hypersetup{hidelinks}

% ============================================================
% Palette « Pop » OnBuch+ — mêmes hex que la classe P (plus_ui.dart)
% ============================================================
\definecolor{poporange}{HTML}{F59321}
\definecolor{popdark}{HTML}{E07A0C}
\definecolor{popcream}{HTML}{FFF6E8}
\definecolor{popink}{HTML}{1C1714}
\definecolor{poppurple}{HTML}{7A5AE0}
\definecolor{popgreen}{HTML}{2E9E6A}
\definecolor{poppink}{HTML}{E0457B}
\definecolor{popblue}{HTML}{2F7DE1}
\definecolor{popgold}{HTML}{C98A12}
\definecolor{popmuted}{HTML}{8A7F76}
\definecolor{popline}{HTML}{EADFD2}
\definecolor{popgray}{HTML}{8A7F76}
\colorlet{popgrey}{popgray}
% Alias tolérants (le modèle nomme parfois les couleurs différemment)
\colorlet{popwhite}{white}
\colorlet{popblack}{popink}
\colorlet{popred}{poppink}
\colorlet{popyellow}{popgold}
% Teintes claires pour fonds de tableaux / zones
\colorlet{popblueL}{popblue!10!white}
\colorlet{poporangeL}{poporange!14!white}
\colorlet{popgreenL}{popgreen!12!white}
\colorlet{poppurpleL}{poppurple!12!white}
\colorlet{poppinkL}{poppink!12!white}
\colorlet{popgoldL}{popgold!14!white}

\pagecolor{popcream}

% ============================================================
% Typographie
% ============================================================
\defaultfontfeatures{Ligatures=TeX}
\setmainfont{PlusJakartaSans-Medium.ttf}[Path=../../../fonts/,BoldFont=PlusJakartaSans-Bold.ttf]
% Maths en sans-serif (Fira Math) avec chiffres et lettres droites de Plus
% Jakarta, pour que les formules se fondent dans le texte.
\usepackage[math-style=ISO,bold-style=ISO]{unicode-math}
\setmathfont{FiraMath-Regular.otf}[Path=../../../fonts/]
\setmathfont{PlusJakartaSans-Medium.ttf}[Path=../../../fonts/,range={up/num,up/{Latin,latin},\mathup}]
\usepackage{icomma} % virgule décimale sans espace en mode math
% Caractères mathématiques Unicode absents des polices de texte (Plus Jakarta,
% Archivo Black) : rendus par leur équivalent mathématique, en texte comme en
% formule — sinon ils s'affichent en carré vide (ex. « Arithmétique dans ℤ »).
\usepackage{newunicodechar}
\newunicodechar{ℝ}{\ensuremath{\mathbb{R}}}
\newunicodechar{ℤ}{\ensuremath{\mathbb{Z}}}
\newunicodechar{ℂ}{\ensuremath{\mathbb{C}}}
\newunicodechar{ℕ}{\ensuremath{\mathbb{N}}}
\newunicodechar{ℚ}{\ensuremath{\mathbb{Q}}}
\newunicodechar{𝔻}{\ensuremath{\mathbb{D}}}
\newunicodechar{α}{\ensuremath{\alpha}}
\newunicodechar{β}{\ensuremath{\beta}}
\newunicodechar{γ}{\ensuremath{\gamma}}
\newunicodechar{δ}{\ensuremath{\delta}}
\newunicodechar{ε}{\ensuremath{\varepsilon}}
\newunicodechar{θ}{\ensuremath{\theta}}
\newunicodechar{λ}{\ensuremath{\lambda}}
\newunicodechar{μ}{\ensuremath{\mu}}
\newunicodechar{σ}{\ensuremath{\sigma}}
\newunicodechar{τ}{\ensuremath{\tau}}
\newunicodechar{φ}{\ensuremath{\varphi}}
\newunicodechar{ψ}{\ensuremath{\psi}}
\newunicodechar{ω}{\ensuremath{\omega}}
\newunicodechar{Σ}{\ensuremath{\Sigma}}
% majuscules produites par \MakeUppercase dans les titres (α-aminés -> Α)
\newunicodechar{Α}{\ensuremath{\alpha}}
\newunicodechar{Β}{\ensuremath{\beta}}
\newunicodechar{Γ}{\ensuremath{\Gamma}}
\newunicodechar{Δ}{\ensuremath{\Delta}}
\newunicodechar{Ε}{\ensuremath{\varepsilon}}
\newunicodechar{Θ}{\ensuremath{\Theta}}
\newunicodechar{Λ}{\ensuremath{\Lambda}}
\newunicodechar{Μ}{\ensuremath{\mu}}
\newunicodechar{Τ}{\ensuremath{\tau}}
\newunicodechar{Φ}{\ensuremath{\Phi}}
\newunicodechar{Ψ}{\ensuremath{\Psi}}
\newunicodechar{Ω}{\ensuremath{\Omega}}
\newunicodechar{↦}{\ensuremath{\mapsto}}
\newunicodechar{∈}{\ensuremath{\in}}
\newunicodechar{∉}{\ensuremath{\notin}}
\newunicodechar{∧}{\ensuremath{\wedge}}
\newunicodechar{⇔}{\ensuremath{\Leftrightarrow}}
\newunicodechar{⟺}{\ensuremath{\Longleftrightarrow}}
\newunicodechar{⇒}{\ensuremath{\Rightarrow}}
\newunicodechar{⟹}{\ensuremath{\Longrightarrow}}
\newunicodechar{∘}{\ensuremath{\circ}}
\newunicodechar{∪}{\ensuremath{\cup}}
\newunicodechar{∩}{\ensuremath{\cap}}
\newunicodechar{≡}{\ensuremath{\equiv}}
\newunicodechar{⊂}{\ensuremath{\subset}}
\newunicodechar{⊥}{\ensuremath{\perp}}
\newunicodechar{∃}{\ensuremath{\exists}}
\newunicodechar{∀}{\ensuremath{\forall}}
\newunicodechar{∛}{\ensuremath{\sqrt[3]{\,}}}
\newunicodechar{∜}{\ensuremath{\sqrt[4]{\,}}}
\newunicodechar{⃗}{\ensuremath{{}^{\to}}}
\newunicodechar{ⁿ}{\ifmmode^{n}\else\textsuperscript{\ensuremath{n}}\fi}
\newunicodechar{ˣ}{\ifmmode^{x}\else\textsuperscript{\ensuremath{x}}\fi}
\newunicodechar{ᵇ}{\ifmmode^{b}\else\textsuperscript{\ensuremath{b}}\fi}
\newunicodechar{ʳ}{\ifmmode^{r}\else\textsuperscript{\ensuremath{r}}\fi}
\newunicodechar{ᵘ}{\ifmmode^{u}\else\textsuperscript{\ensuremath{u}}\fi}
\newunicodechar{ʸ}{\ifmmode^{y}\else\textsuperscript{\ensuremath{y}}\fi}
\newunicodechar{ᵃ}{\ifmmode^{a}\else\textsuperscript{\ensuremath{a}}\fi}
\newunicodechar{ᵉ}{\ifmmode^{e}\else\textsuperscript{\ensuremath{e}}\fi}
\newunicodechar{ᵏ}{\ifmmode^{k}\else\textsuperscript{\ensuremath{k}}\fi}
\newunicodechar{ᵖ}{\ifmmode^{p}\else\textsuperscript{\ensuremath{p}}\fi}
\newunicodechar{ᴺ}{\ifmmode^{N}\else\textsuperscript{\ensuremath{N}}\fi}
\newunicodechar{ᶜ}{\ifmmode^{c}\else\textsuperscript{\ensuremath{c}}\fi}
\newunicodechar{ʰ}{\ifmmode^{h}\else\textsuperscript{\ensuremath{h}}\fi}
\newunicodechar{ᵠ}{\ifmmode^{q}\else\textsuperscript{\ensuremath{q}}\fi}
\newunicodechar{ₙ}{\ifmmode_{n}\else\textsubscript{\ensuremath{n}}\fi}
\newunicodechar{ₐ}{\ifmmode_{a}\else\textsubscript{\ensuremath{a}}\fi}
\newunicodechar{ᵢ}{\ifmmode_{i}\else\textsubscript{\ensuremath{i}}\fi}
\newunicodechar{ᵣ}{\ifmmode_{r}\else\textsubscript{\ensuremath{r}}\fi}
\newunicodechar{ₖ}{\ifmmode_{k}\else\textsubscript{\ensuremath{k}}\fi}
\newunicodechar{ᵦ}{\ifmmode_{\beta}\else\textsubscript{\ensuremath{\beta}}\fi}
\newunicodechar{ₓ}{\ifmmode_{x}\else\textsubscript{\ensuremath{x}}\fi}
\newfontfamily\popdisplay{ArchivoBlack-Regular.ttf}[Path=../../../fonts/]
\newfontfamily\poplogo{SpaceGrotesk-Bold.ttf}[Path=../../../fonts/]
\newfontfamily\popsemi{PlusJakartaSans-SemiBold.ttf}[Path=../../../fonts/]
\newfontfamily\popxbold{PlusJakartaSans-ExtraBold.ttf}[Path=../../../fonts/]
\newfontfamily\popxboldls{PlusJakartaSans-ExtraBold.ttf}[Path=../../../fonts/,LetterSpace=3.0]

\setlist[enumerate]{leftmargin=1.7em,itemsep=5pt,topsep=4pt}
\setlist[itemize]{leftmargin=1.7em,itemsep=4pt,topsep=4pt,label={\color{poporange}\textbullet}}
\setlength{\parindent}{0pt}
\setlength{\parskip}{5pt}
\renewcommand{\arraystretch}{1.25}

% pgfplots : style Pop par défaut
\pgfplotsset{
  every axis/.append style={axis line style={popink,line width=1pt},tick style={popink},
    grid style={popline,line width=0.5pt},label style={font=\small},tick label style={font=\footnotesize}},
  popcurve/.style={poporange,line width=1.8pt,no markers,smooth},
}
% Même style disponible dans un \draw TikZ simple (hors pgfplots)
\tikzset{popcurve/.style={poporange,line width=1.8pt}}
\tikzset{popgrid/.style={popline,very thin}}
\colorlet{popcurve}{poporange} % le nom du style est parfois utilisé comme couleur

% ============================================================
% Pill / squiggle
% ============================================================
\newcommand{\poppill}[3]{%
  \tikz[baseline=-0.55ex]\node[draw=popink,line width=1.2pt,rounded corners=2.6pt,%
    fill=#2,inner xsep=6.5pt,inner ysep=3pt]{%
    \popxboldls\fontsize{7}{7}\selectfont\color{#3}\MakeUppercase{#1}};%
}
\newcommand{\popsquiggle}[1]{%
  \tikz[baseline=-0.4ex]{
    \draw[#1,line width=2.6pt,line cap=round]
      plot [smooth] coordinates {(0,0.09) (0.34,0.01) (0.68,0.21) (1.02,0.01) (1.36,0.10)};
  }%
}
% Mot-clé surligné (marqueur orange sous le texte)
\newcommand{\cle}[1]{{\popxbold\color{popdark}#1}}

% ============================================================
% En-tête / pied
% ============================================================
\pagestyle{fancy}
\fancyhf{}
\renewcommand{\headrulewidth}{0pt}
\renewcommand{\footrulewidth}{0pt}
\lhead{\raisebox{-0.15ex}{\poplogo\fontsize{13}{13}\selectfont\color{popink}OnBuch\color{poporange}+}}
\rhead{\poppill{\DOCMATIERE\ \textbullet\ \DOCNIVEAU\ \textbullet\ Leçon \DOCLECON}{white}{popink}}
\lfoot{\popxbold\fontsize{7.6}{7.6}\selectfont\color{popmuted}\MakeUppercase{Cours signé OnBuch+}\ \textbullet\ {\popsemi\DOCTITRE}}
\rfoot{\tikz[baseline=-0.6ex]\node[rounded corners=8.5pt,fill=popink,minimum height=17pt,minimum width=17pt,inner xsep=5pt,inner sep=0]{\popxbold\fontsize{8}{8}\selectfont\color{popcream}\thepage\,/\,\pageref*{LastPage}};}

% ============================================================
% Titres
% ============================================================
\newcommand{\chaptitre}[2]{%
  \begin{tcolorbox}[enhanced,colback=poporange,colframe=popink,boxrule=2.6pt,arc=11pt,
      drop shadow=popink,left=15pt,right=15pt,top=12pt,bottom=11pt,before skip=3pt,after skip=18pt]
    \poppill{#2}{popink}{popcream}\par\vspace{9pt}
    {\raggedright\hyphenpenalty=10000\exhyphenpenalty=10000\popdisplay\fontsize{22}{24}\selectfont\color{popcream}\MakeUppercase{#1}\par}
  \end{tcolorbox}
}
% Section numérotée : pastille orange avec le numéro + titre Archivo
\newcounter{popsec}
\newcommand{\coursec}[1]{%
  \stepcounter{popsec}\setcounter{popsub}{0}%
  \par\vspace{16pt}\noindent
  \begin{minipage}{\linewidth}
  \tikz[baseline=-0.7ex]\node[draw=popink,line width=1.6pt,fill=poporange,rounded corners=4pt,minimum size=22pt,inner sep=0]{\popdisplay\fontsize{12}{12}\selectfont\color{popcream}\thepopsec};%
  \hspace{8pt}{\popdisplay\fontsize{15}{17}\selectfont\color{popink}\MakeUppercase{#1}}\par
  \vspace{4pt}\noindent{\color{poporange}\rule{36pt}{3.6pt}}
  \end{minipage}\par\nopagebreak\vspace{8pt}
}
\newcounter{popsub}
\newcommand{\courssub}[1]{%
  \stepcounter{popsub}%
  \par\vspace{9pt}\noindent{\popxbold\fontsize{11.5}{13}\selectfont\color{popink}{\color{poporange}\thepopsec.\thepopsub}\hspace{6pt}#1}\par\nopagebreak\vspace{4pt}
}

% ============================================================
% Boîtes pédagogiques — même recette : fond blanc, bordure épaisse colorée,
% ombre dure encre, badge plein. Titre optionnel [#1].
% ============================================================
\tcbset{popbase/.style={breakable,enhanced,colback=white,boxrule=2.3pt,arc=9pt,
  shadow={3pt}{-3pt}{0pt}{popink},left=14pt,right=14pt,top=11pt,bottom=11pt,before skip=11pt,after skip=15pt,
  fontupper=\popsemi\fontsize{10.6}{14.5}\selectfont\color{popink},
  fontlower=\popsemi\fontsize{10.6}{14.5}\selectfont\color{popink}}}
% Coche dessinée (le glyphe ✓ n'existe pas dans les polices du document)
\newcommand{\popcheck}[1]{\tikz[baseline=-0.5ex]\draw[#1,line width=1.6pt,line cap=round,line join=round](-0.13,0.0)--(-0.04,-0.09)--(0.13,0.1);}
\providecommand{\coche}{\popcheck{popgreen}}
\providecommand{\resultat}[1]{\par\smallskip\noindent\fcolorbox{popgreen}{popgreenL}{\parbox{\dimexpr\linewidth-2\fboxsep-2\fboxrule\relax}{\textbf{Résultat.} #1}}\par\smallskip}
\newcommand{\popboxhead}[3]{\poppill{#1}{#2}{white}\ifx\relax#3\relax\else\hspace{6pt}{\popxbold\fontsize{10.6}{12}\selectfont\color{popink}#3}\fi\par\vspace{7pt}}

\newtcolorbox{aretenir}[1][]{popbase,colframe=popgreen,before upper={\popboxhead{À retenir}{popgreen}{#1}}}
\newtcolorbox{exemplebox}[1][]{popbase,colframe=poporange,before upper={\popboxhead{Exemple}{poporange}{#1}}}
\newtcolorbox{definition}[1][]{popbase,colframe=popblue,before upper={\popboxhead{Définition}{popblue}{#1}}}
\newtcolorbox{propriete}[1][]{popbase,colframe=poppurple,before upper={\popboxhead{Propriété}{poppurple}{#1}}}
\newtcolorbox{methode}[1][]{popbase,colframe=popgold,before upper={\popboxhead{Méthode}{popgold}{#1}}}
\newtcolorbox{attention}[1][]{popbase,colframe=poppink,before upper={\popboxhead{Attention}{poppink}{#1}}}
\newtcolorbox{experience}[1][]{popbase,colframe=popblue,colback=popblueL,before upper={\popboxhead{Expérience}{popblue}{#1}}}
% « Le savais-tu ? » : carte violette pleine (contexte, culture, Cameroun)
\newtcolorbox{savaistu}[1][]{popbase,colback=poppurple,colframe=popink,
  fontupper=\popsemi\fontsize{10.4}{14.2}\selectfont\color{white},
  before upper={\poppill{Le savais-tu\,?}{white}{poppurple}\ifx\relax#1\relax\else\hspace{6pt}{\popxbold\color{white}#1}\fi\par\vspace{7pt}}}
% Exercice résolu : énoncé en haut, \tcblower puis la solution sur fond crème
\newcounter{popexo}\newcounter{popexr}
\newtcolorbox{exoresolu}[1][]{popbase,colframe=poporange,colbacklower=popcream,
  segmentation style={popink,line width=1.2pt,dashed},
  before upper={\stepcounter{popexr}\popboxhead{Exercice résolu \thepopexr}{poporange}{#1}},
  before lower={\poppill{Solution}{popink}{popcream}\par\vspace{6pt}}}
% Exercice (non résolu en place) — la correction est dans la partie Corrigés
\newtcolorbox{exercice}[1][]{popbase,colframe=popink,
  before upper={\stepcounter{popexo}\popboxhead{Exercice \thepopexo}{popink}{#1}}}
% Corrigé
\newtcolorbox{corrige}[1][]{popbase,colframe=popgreen,colback=popgreenL,
  before upper={\popboxhead{Corrigé}{popgreen}{#1}}}
% Objectifs / prérequis / bilan
\newtcolorbox{objectifs}{popbase,colframe=popink,colback=poporangeL,
  before upper={\popboxhead{Objectifs de la leçon}{popink}{}}}
\newtcolorbox{prerequis}{popbase,colframe=popink,colback=popblueL,
  before upper={\popboxhead{Prérequis}{popblue}{}}}
\newtcolorbox{bilan}{popbase,colframe=popink,colback=white,boxrule=2.8pt,
  before upper={\popboxhead{Fiche bilan}{poporange}{L'essentiel à connaître pour l'examen}}}
% Figure encadrée avec légende
\newtcolorbox{popfigure}[1][]{enhanced,breakable=false,colback=white,colframe=popline,boxrule=1.4pt,arc=8pt,
  left=10pt,right=10pt,top=10pt,bottom=8pt,before skip=10pt,after skip=12pt,halign=center,
  lower separated=false,halign lower=center,
  fontlower=\popsemi\fontsize{9}{11.5}\selectfont\color{popmuted},
  #1}
\newcommand{\legende}[1]{\par\vspace{6pt}{\popsemi\fontsize{9}{11.5}\selectfont\color{popmuted}#1}}

% ============================================================
% Bloc de code source — \begin{lstlisting}[language=Php]...\end{lstlisting}
% (ou language=C, HTML, JavaScript, SQL ; option omise = pseudo-code brut).
% Même recette visuelle que les figures (\popfigure) : cadre encre épais,
% fond clair (popcream), légende possible juste après via \legende{...}
% comme pour une figure (listings ne sait pas arrondir ses coins : on garde
% un cadre rectangulaire, cohérent avec les autres tableaux du document).
% CONTRAT : les modèles ne doivent utiliser QUE \begin{lstlisting}[...] tel
% quel (pas de \begin{tcblisting}, pas de \verb pour un bloc multi-lignes).
% ============================================================
\lstdefinestyle{poplst}{
  backgroundcolor=\color{popcream},
  basicstyle=\popsemi\fontsize{8.6}{11}\selectfont\color{popink},
  keywordstyle=\bfseries\color{popblue},
  commentstyle=\itshape\color{popmuted},
  stringstyle=\color{popgreen},
  numberstyle=\tiny\color{popmuted},
  identifierstyle=\color{popink},
  frame=l,
  rulecolor=\color{popink},
  framerule=1.6pt,
  framesep=7pt,
  framexleftmargin=7pt,
  framexrightmargin=7pt,
  xleftmargin=2pt,
  xrightmargin=2pt,
  breaklines=true,
  breakatwhitespace=false,
  showstringspaces=false,
  showspaces=false,
  showtabs=false,
  tabsize=2,
  columns=flexible,
  keepspaces=true,
  numbers=none,
  aboveskip=10pt,
  belowskip=10pt,
  captionpos=b,
}
\lstset{style=poplst, language=PHP}
% La distribution TeX embarquée ne fournit pas le fichier de langue
% JavaScript de listings (lstlang*.dat) : on la redéfinit nous-mêmes
% pour que \begin{lstlisting}[language=JavaScript] fonctionne.
\lstdefinelanguage{JavaScript}{
  keywords={break,case,catch,class,const,continue,debugger,default,delete,do,
    else,export,extends,finally,for,function,if,import,in,instanceof,let,new,
    return,static,super,switch,this,throw,try,typeof,var,void,while,with,yield,
    async,await,of,get,set},
  keywordstyle=\bfseries\color{popblue},
  ndkeywords={true,false,null,undefined,NaN,Infinity},
  ndkeywordstyle=\bfseries\color{poporange},
  identifierstyle=\color{popink},
  sensitive=true,
  comment=[l]{//},
  morecomment=[s]{/*}{*/},
  string=[b]",
  morestring=[b]',
  morestring=[b]`,
}
% Même limitation pour CSS et les fichiers de configuration (apacheconf,
% nginx, ini...) : ils ne sont pas fournis. CSS et un style générique de
% type "config" (commentaires # ou ;, pas de mots-clés) couvrent les besoins.
\lstdefinelanguage{CSS}{
  keywords={color,background,background-color,margin,padding,border,
    border-radius,font,font-size,font-weight,font-family,width,height,
    display,position,top,left,right,bottom,float,clear,text-align,
    line-height,list-style,cursor,overflow,box-shadow,transition,
    flex,grid,align-items,justify-content,z-index},
  keywordstyle=\bfseries\color{popblue},
  identifierstyle=\color{popink},
  sensitive=false,
  comment=[s]{/*}{*/},
  string=[b]",
  morestring=[b]',
}
\lstdefinelanguage{apacheconf}{
  sensitive=false,
  morecomment=[l]{\#},
}
\lstdefinelanguage{Apache}{
  sensitive=false,
  morecomment=[l]{\#},
}
\lstdefinelanguage{text}{sensitive=false}
\lstdefinelanguage{powershell}{
  keywords={New-Object,New-SmbShare,Get-Acl,Set-Acl,Get-ChildItem,Set-Content,
    Get-Content,Write-Host,ForEach-Object,Where-Object,Select-Object,
    Test-Path,Remove-Item,Copy-Item,Move-Item,Start-Process,Stop-Process,
    If,Else,ElseIf,ForEach,While,Function,Return,Try,Catch},
  keywordstyle=\bfseries\color{popblue},
  identifierstyle=\color{popink},
  sensitive=false,
  comment=[l]{\#},
  string=[b]",
  morestring=[b]',
}
\lstdefinelanguage{ini}{
  sensitive=false,
  morecomment=[l]{;},
  morecomment=[l]{\#},
}

% ============================================================
% Signature OnBuch+
% ============================================================
\newcommand{\poplogoinline}[1]{{\poplogo\fontsize{#1}{#1}\selectfont\color{popink}OnBuch\color{poporange}+}}
\newcommand{\signatureonbuch}{%
  \begin{tcolorbox}[enhanced,colback=popink,colframe=popink,boxrule=2.2pt,arc=10pt,
      drop shadow=poporange,left=14pt,right=14pt,top=10pt,bottom=10pt,before skip=0pt,after skip=0pt]
    \tikz[baseline=-0.7ex]\node[circle,fill=popgreen,draw=popcream,line width=1.3pt,minimum size=22pt,inner sep=0]{\popcheck{white}};%
    \hspace{9pt}\begin{minipage}[c]{0.62\linewidth}
      {\popxbold\fontsize{12}{13}\selectfont\color{popcream}Signé\ \textbullet\ {\poplogo OnBuch\color{poporange}+}}\par\vspace{2pt}
      {\raggedright\popsemi\fontsize{8}{10.5}\selectfont\color{popline}Ludovic A.\\\MakeUppercase{Conforme au programme officiel MINESEC}\par}
    \end{minipage}\hfill
    {\poplogo\fontsize{22}{22}\selectfont\color{popcream}OnBuch\color{poporange}+}
  \end{tcolorbox}%
}

% ============================================================
% Page de garde
% #1 titre · #2 matière · #3 niveau · #4 module · #5 n° leçon · #6 durée ·
% #7 description / objectifs (texte ou itemize)
% ============================================================
\newcommand{\pagegarde}[7]{%
  \thispagestyle{empty}
  \newgeometry{a4paper,margin=1.7cm,top=1.6cm,bottom=1.4cm}%
  \begin{tikzpicture}[remember picture,overlay]
    \fill[poporange!18] ([xshift=-3.2cm,yshift=-3.6cm]current page.north east) circle (4.6cm);
    \fill[poppurple!14] ([xshift=2.4cm,yshift=7.6cm]current page.south west) circle (3.8cm);
  \end{tikzpicture}%
  {\poplogo\fontsize{30}{30}\selectfont\color{popink}OnBuch\color{poporange}+}\hfill
  \raisebox{0.6ex}{\poppill{Cours signé \textbullet\ Premium}{popink}{popcream}}\par
  \vspace{1pt}\popsquiggle{poppurple}\par
  \vspace{8pt}
  \begin{tcolorbox}[enhanced,colback=poporange,colframe=popink,boxrule=2.8pt,arc=13pt,
      drop shadow=popink,left=18pt,right=18pt,top=12pt,bottom=13pt,after skip=10pt]
    \poppill{#2 \textbullet\ #3}{popink}{popcream}\hspace{5pt}\poppill{#4}{white}{popink}\par\vspace{8pt}
    {\popdisplay\fontsize{14}{14}\selectfont\color{popink}LEÇON #5}\par\vspace{4pt}
    {\raggedright\hyphenpenalty=10000\exhyphenpenalty=10000\popdisplay\fontsize{25}{27}\selectfont\color{popcream}\MakeUppercase{#1}\par}
    \vspace{8pt}
    {\popxbold\fontsize{9.5}{9.5}\selectfont\color{popink}\MakeUppercase{Durée conseillée : #6}}
  \end{tcolorbox}
  \begin{tcolorbox}[enhanced,colback=white,colframe=popink,boxrule=2.2pt,arc=10pt,
      drop shadow=popink,left=16pt,right=16pt,top=10pt,bottom=10pt,after skip=8pt]
    \poppill{Dans cette leçon}{poppurple}{white}\par\vspace{5pt}
    \popsemi\fontsize{9.3}{12.6}\selectfont\color{popink}#7
  \end{tcolorbox}
  \noindent\begin{minipage}[t]{0.487\linewidth}
    \begin{tcolorbox}[enhanced,colback=popgreen,colframe=popink,boxrule=2.2pt,arc=10pt,
        drop shadow=popink,left=11pt,right=11pt,top=10pt,bottom=10pt,height=3.1cm,valign=center]
      \tcbox[colback=white,colframe=popink,boxrule=1.3pt,arc=5pt,boxsep=0pt,nobeforeafter,
        left=3pt,right=3pt,top=3pt,bottom=3pt,valign=center]{\qrcode[height=1.9cm]{https://play.google.com/store/apps/details?id=com.luvvix.onbuch&hl=fr}}%
      \hspace{8pt}\begin{minipage}[c]{0.5\linewidth}
        {\popdisplay\fontsize{11}{12.5}\selectfont\color{white}\MakeUppercase{Télécharge OnBuch}}\par\vspace{4pt}
        {\raggedright\popsemi\fontsize{8.4}{11}\selectfont\color{white}Gratuit \textbullet\ Google Play\\Tuteur IA, annales, concours, résultats\par}
      \end{minipage}
    \end{tcolorbox}
  \end{minipage}\hfill
  \begin{minipage}[t]{0.487\linewidth}
    \begin{tcolorbox}[enhanced,colback=poppurple,colframe=popink,boxrule=2.2pt,arc=10pt,
        drop shadow=popink,left=11pt,right=11pt,top=10pt,bottom=10pt,height=3.1cm,valign=center]
      \tikz[baseline=-0.7ex]\node[circle,fill=white,draw=popink,line width=1.3pt,minimum size=1.6cm,inner sep=0]{\popdisplay\fontsize{24}{24}\selectfont\color{poppurple}+};%
      \hspace{8pt}\begin{minipage}[c]{0.6\linewidth}
        {\popdisplay\fontsize{11}{12.5}\selectfont\color{white}\MakeUppercase{Passe à OnBuch+}}\par\vspace{4pt}
        {\raggedright\popsemi\fontsize{8.4}{11}\selectfont\color{white}Tous les cours signés, Tuteur IA illimité, fiches PDF — onglet OnBuch+ de l'app\par}
      \end{minipage}
    \end{tcolorbox}
  \end{minipage}
  \vspace{8pt}
  \popcontenu
  \vfill
  \signatureonbuch
  \restoregeometry
  \newpage
}

% Rangée « Ce cours contient » (page de garde)
\newcommand{\poptile}[3]{% #1 couleur #2 gros texte #3 légende
  \begin{tcolorbox}[enhanced,colback=#1,colframe=popink,boxrule=2pt,arc=9pt,shadow={2.5pt}{-2.5pt}{0pt}{popink},
      left=6pt,right=6pt,top=9pt,bottom=9pt,halign=center,valign=center,height=1.75cm,width=\linewidth,nobeforeafter]
    {\popdisplay\fontsize{12.5}{13}\selectfont\color{white}\MakeUppercase{#2}}\par\vspace{3pt}
    {\centering\popsemi\fontsize{7.8}{9.5}\selectfont\color{white}#3\par}
  \end{tcolorbox}}
\newcommand{\popcontenu}{%
  \par\noindent{\popxboldls\fontsize{7.5}{7.5}\selectfont\color{popmuted}CE COURS SIGNÉ CONTIENT}\par\vspace{7pt}
  \noindent\begin{minipage}[t]{0.235\linewidth}\poptile{popblue}{Cours}{complet, illustré, pas à pas}\end{minipage}\hfill
  \begin{minipage}[t]{0.235\linewidth}\poptile{poporange}{Méthodes}{et exercices résolus}\end{minipage}\hfill
  \begin{minipage}[t]{0.235\linewidth}\poptile{poppink}{Exercices}{type examen + corrigés}\end{minipage}\hfill
  \begin{minipage}[t]{0.235\linewidth}\poptile{popgreen}{Fiche bilan}{l'essentiel à retenir}\end{minipage}\par}

% Sommaire
\newtcolorbox{sommaire}{enhanced,colback=white,colframe=popink,boxrule=2.2pt,arc=10pt,
  drop shadow=popink,left=18pt,right=18pt,top=13pt,bottom=13pt,before skip=6pt,after skip=20pt,
  before upper={\poppill{Sommaire}{poppurple}{white}\par\vspace{9pt}\popsemi\fontsize{10.6}{14.5}\selectfont\color{popink}}}

% Signature de fin de cours — poussée tout en bas de la dernière page
\newcommand{\findecours}{%
  \par\vfill\nopagebreak
  \begin{center}\popsquiggle{poporange}\end{center}\vspace{4pt}
  \signatureonbuch
}
\AtBeginDocument{\def\llbracket{\mathopen{[\mkern-3.5mu[}}\def\rrbracket{\mathclose{]\mkern-3.5mu]}}} % intervalles d'entiers (glyphe ⟦ absent de la police)
% Glyphes absents de Fira Math : redéfinis à partir de glyphes disponibles
\AtBeginDocument{%
  \def\setminus{\mathbin{\backslash}}\let\smallsetminus\setminus
  \def\vdots{\mathord{\vbox{\baselineskip3.2pt\lineskiplimit0pt\kern4pt\hbox{.}\hbox{.}\hbox{.}}}}%
  \def\ddots{\mathinner{\mkern1mu\raise6.5pt\hbox{.}\mkern2mu\raise3.6pt\hbox{.}\mkern2mu\raise0.7pt\hbox{.}\mkern1mu}}%
  \def\triangle{\mathord{\scalebox{0.9}{$\symup{\Delta}$}}}%
  \def\nabla{\mathord{\raisebox{\depth}{\scalebox{1}[-1]{$\symup{\Delta}$}}}}%
  \def\checkmark{\popcheck{popdark}}%
  \def\star{\mathbin{\ast}}\def\bigstar{\mathord{\ast}}%
  \def\diamond{\mathbin{\scalebox{0.6}{$\Diamond$}}}%
  \def\bigcup{\mathop{\mathchoice{\raisebox{-0.3em}{\scalebox{1.7}{$\cup$}}}{\scalebox{1.25}{$\cup$}}{\cup}{\cup}}\displaylimits}%
  \def\bigcap{\mathop{\mathchoice{\raisebox{-0.3em}{\scalebox{1.7}{$\cap$}}}{\scalebox{1.25}{$\cap$}}{\cap}{\cap}}\displaylimits}%
  \def\lesseqqgtr{\mathrel{\lessgtr}}\def\bigoplus{\mathop{\oplus}\displaylimits}%
}
\providecommand{\ssi}{si et seulement si} % abréviation fréquente
\AtBeginDocument{\providecommand{\wideparen}{\overparen}} % arc de cercle

%% FIGFIX-BEGIN v5 — correctifs globaux des figures (docs/onbuch-generation/tools/figfix)
\usepackage{adjustbox}\usepackage{etoolbox}\tcbuselibrary{raster}
% toute figure TikZ/pgfplots plus large que la ligne est réduite à la largeur disponible
\BeforeBeginEnvironment{tikzpicture}{\begin{adjustbox}{max width=\linewidth}}
\AfterEndEnvironment{tikzpicture}{\end{adjustbox}}
% légendes pgfplots lisibles par-dessus les courbes
\pgfplotsset{every axis/.append style={legend style={fill=white,fill opacity=0.92,text opacity=1,draw=black!15,inner sep=3pt}}}
% étiquettes d'arêtes (syntaxe edge["..."]) avec fond blanc
\tikzset{every edge quotes/.append style={fill=white,inner sep=1.2pt}}
%% FIGFIX-END
\begin{document}

\pagegarde{Initiation à la publication assistée par ordinateur}{Informatique}{3ème}{Informatique – classe de 3ème}{N07}{1 séance (fiche de progression harmonisée)}{Cette leçon initie les élèves à la Publication Assistée par Ordinateur (PAO) : découverte des notions de base de l'infographie, des types d'images numériques et des logiciels de mise en page. À travers des situations concrètes de la vie courante (affiches, bulletins, cartes de visite), les élèves apprennent à concevoir et justifier une démarche de création de documents imprimés de qualité.
\vspace{4pt}
\begin{multicols}{2}\small
\begin{itemize}
\item À la fin de cette leçon, je sais définir l'infographie et la Publication Assistée par Ordinateur (PAO)
\item À la fin de cette leçon, je sais distinguer une image matricielle (bitmap) d'une image vectorielle et citer leurs usages
\item À la fin de cette leçon, je sais citer au moins trois logiciels de PAO et leurs domaines d'application
\item À la fin de cette leçon, je sais identifier les éléments d'une mise en page (marges, colonnes, orientation, en-tête, pied de page)
\item À la fin de cette leçon, je sais décrire les étapes de conception d'un document de PAO (analyse du besoin, maquette, réalisation, impression)
\item À la fin de cette leçon, je sais appliquer les notions de la leçon à des situations concrètes de la vie courante (affiche, bulletin, carte d'invitation)
\end{itemize}\end{multicols}}

\begin{sommaire}
\begin{multicols}{2}
\begin{enumerate}
\item Infographie et PAO : définitions
\item Les images numériques : matricielles et vectorielles
\item Logiciels d'infographie et de PAO
\item Les éléments d'une mise en page
\item Démarche de conception d'un document de PAO
\item Applications concrètes et rédaction de la démarche
\item Activité d'intégration
\item Méthodes et exercices résolus
\item Exercices
\item Corrigés des exercices
\item Fiche bilan
\end{enumerate}
\end{multicols}
\end{sommaire}

\begin{objectifs}
\begin{itemize}
\item À la fin de cette leçon, je sais définir l'infographie et la Publication Assistée par Ordinateur (PAO)
\item À la fin de cette leçon, je sais distinguer une image matricielle (bitmap) d'une image vectorielle et citer leurs usages
\item À la fin de cette leçon, je sais citer au moins trois logiciels de PAO et leurs domaines d'application
\item À la fin de cette leçon, je sais identifier les éléments d'une mise en page (marges, colonnes, orientation, en-tête, pied de page)
\item À la fin de cette leçon, je sais décrire les étapes de conception d'un document de PAO (analyse du besoin, maquette, réalisation, impression)
\item À la fin de cette leçon, je sais appliquer les notions de la leçon à des situations concrètes de la vie courante (affiche, bulletin, carte d'invitation)
\item À la fin de cette leçon, je sais rédiger et justifier une démarche de création d'un document, une réponse ou une conclusion
\item À la fin de cette leçon, je sais choisir le logiciel adapté à un besoin précis et justifier mon choix
\end{itemize}
\end{objectifs}

\begin{prerequis}
\begin{itemize}
\item Notions générales sur l'ordinateur et ses périphériques (classe de 4ème)
\item Utilisation de base d'un logiciel de traitement de texte (saisie, mise en forme simple)
\item Notion de fichier, dossier et formats de fichiers courants
\item Manipulation de la souris et du clavier, environnement graphique (fenêtres, menus)
\end{itemize}
\end{prerequis}

\newpage

\coursec{Infographie et PAO : définitions}

L'association des jeunes du quartier Nkol-Eton à Yaoundé organise son grand tournoi de football annuel. Pour faire venir du public, il faut une belle affiche à coller au carrefour, et aussi un petit bulletin d'information de deux pages qui raconte la vie du quartier. Le président de l'association se rend au cybercafé du coin : le patron annonce \num{5 000} FCFA par page de conception. Pour l'affiche et les deux pages du bulletin, cela fait \num{15 000} FCFA, une somme importante pour la caisse de l'association !

Armelle, élève de 3ème, assiste à la discussion. Elle se dit : « Mais nous avons des ordinateurs à l'école, et je sais déjà taper un texte. Ne pourrais-je pas concevoir moi-même ces documents, avec des textes bien présentés, des images et des couleurs ? » Sa question en soulève d'autres : quels logiciels utiliser ? Quelles règles respecter pour obtenir un document qui ait l'air \emph{professionnel}, et non un simple texte tapé à la va-vite ?

C'est exactement l'objet de cette leçon : découvrir l'\cle{infographie} et la \cle{Publication Assistée par Ordinateur} (PAO), comprendre ce qui les distingue du traitement de texte que tu connais déjà, et voir à quels documents de la vie courante elles s'appliquent.

\textbf{Problématique :} comment l'ordinateur permet-il de concevoir des documents destinés à l'impression ou à la diffusion, et en quoi cette activité dépasse-t-elle la simple saisie d'un texte ?

\courssub{Qu'est-ce que l'infographie ?}

\courssub{L'intuition : l'ordinateur comme atelier d'images}

Avant l'ordinateur, créer une image destinée à l'impression (un logo, une affiche, une illustration de journal) demandait des outils coûteux : pinceaux, encres, tables à dessin, machines de reprographie. La moindre erreur obligeait à tout recommencer. Avec l'ordinateur, on peut dessiner, retoucher une photo, déplacer un élément, changer une couleur, puis annuler et recommencer autant de fois qu'on veut. L'ordinateur est devenu un \emph{atelier complet de création d'images}. Cet ensemble de techniques porte un nom : l'infographie.

\begin{definition}[Infographie]
L'\cle{infographie} est l'ensemble des techniques de \textbf{création} et de \textbf{manipulation} d'\textbf{images numériques} à l'aide d'un ordinateur.
\end{definition}

Décortiquons cette définition :

\begin{itemize}
\item \textbf{Création} : produire une image à partir de rien, par exemple dessiner le logo du tournoi de Nkol-Eton (un ballon stylisé et le nom du quartier).
\item \textbf{Manipulation} : modifier une image existante, par exemple recadrer la photo de l'équipe, éclaircir une photo trop sombre, ou détourer un joueur pour le placer sur un fond coloré.
\item \textbf{Images numériques} : des images stockées sous forme de fichiers informatiques (formats JPEG, PNG, etc.), que l'ordinateur peut enregistrer, copier et transmettre.
\item \textbf{À l'aide d'un ordinateur} : c'est l'outil central, avec des logiciels spécialisés que nous découvrirons dans la suite de la leçon.
\end{itemize}

\begin{exemplebox}[L'infographie autour de toi]
L'infographie est partout dans la vie quotidienne camerounaise :
\begin{itemize}
\item la retouche des photos d'identité dans les studios photo de Yaoundé ou de Douala ;
\item la conception des logos des entreprises (brasseries, compagnies de transport, boutiques) ;
\item les pochettes de musique et les affiches de concerts ;
\item les illustrations des manuels scolaires que tu utilises en classe.
\end{itemize}
\end{exemplebox}

\courssub{Qu'est-ce que la PAO ?}

\courssub{L'intuition : assembler texte, images et couleurs sur une page}

Créer une image, c'est une chose. Mais l'affiche du tournoi n'est pas qu'une image : elle contient aussi un titre (« Tournoi de football de Nkol-Eton »), des informations (dates, lieu, frais d'inscription), une photo, un logo, des couleurs. Il faut \emph{disposer tous ces éléments sur la page} de manière harmonieuse et lisible, puis imprimer le résultat. Cette activité d'assemblage et de mise en page porte un nom précis : la PAO.

\begin{definition}[Publication Assistée par Ordinateur (PAO)]
La \cle{Publication Assistée par Ordinateur} (\cle{PAO}) est l'utilisation de l'ordinateur et de \textbf{logiciels spécialisés} pour \textbf{concevoir des documents destinés à l'impression ou à la diffusion} : affiches, journaux, brochures, cartes de visite, faire-part, etc.
\end{definition}

Trois idées sont essentielles dans cette définition :

\begin{enumerate}
\item \textbf{Des logiciels spécialisés} : la PAO utilise des logiciels conçus pour la mise en page (nous les étudierons en section 3), qui permettent de placer chaque élément avec précision.
\item \textbf{Concevoir} : la PAO n'est pas seulement de la saisie ; c'est un travail de \emph{conception} où l'on choisit les polices, les couleurs, la position des images, l'organisation en colonnes.
\item \textbf{Destinés à l'impression ou à la diffusion} : le document final est pensé pour être imprimé (affiche, journal papier) ou diffusé (bulletin envoyé en PDF par WhatsApp, par exemple).
\end{enumerate}

\begin{aretenir}[Infographie et PAO : deux activités complémentaires]
L'\textbf{infographie} s'occupe des \textbf{images} (les créer, les retoucher). La \textbf{PAO} s'occupe de la \textbf{page entière} : elle assemble le texte, les images (souvent produites par infographie) et les couleurs pour fabriquer un document prêt à imprimer ou à diffuser. Les deux se complètent : on retouche la photo de l'équipe (infographie), puis on la place dans l'affiche avec le titre et les dates (PAO).
\end{aretenir}

La production d'un document de PAO suit une chaîne logique, de l'idée jusqu'à l'impression :

\begin{popfigure}
\begin{center}
\begin{tikzpicture}[
  etape/.style={rectangle, rounded corners=4pt, draw=popblue, fill=popblueL, thick, minimum width=2.6cm, minimum height=1.1cm, align=center, font=\small\bfseries},
  descr/.style={font=\scriptsize, text=popdark, align=center, below=2pt},
  fleche/.style={-{Stealth[length=3mm]}, very thick, poporange}
]
\node[etape] (idee) at (0,0) {Idée};
\node[etape] (maquette) at (3.6,0) {Maquette};
\node[etape, align=center] (saisie) at (7.2,0){Saisie et\\mise en page};
\node[etape] (impression) at (10.8,0) {Impression};
\node[descr, align=center] at (idee.south){que veut-on\\communiquer ?};
\node[descr, align=center] at (maquette.south){croquis de la page\\sur papier};
\node[descr, align=center] at (saisie.south){travail sur ordinateur\\avec le logiciel de PAO};
\node[descr, align=center] at (impression.south){document final\\papier ou PDF};
\draw[fleche] (idee) -- (maquette);
\draw[fleche] (maquette) -- (saisie);
\draw[fleche] (saisie) -- (impression);
\end{tikzpicture}
\end{center}
\legende{La chaîne de production d'un document de PAO : de l'idée à l'impression.}
\end{popfigure}

\courssub{PAO et traitement de texte : quelle différence ?}

Tu connais déjà le traitement de texte (par exemple Microsoft Word) : tu y as tapé des lettres et des exposés. Alors, en quoi la PAO est-elle différente ?

\begin{propriete}[Différence entre traitement de texte et PAO]
\begin{itemize}
\item Le \textbf{traitement de texte} gère surtout le \textbf{texte} : saisie, correction, orthographe, paragraphes, styles simples. Le texte s'écoule de haut en bas, page après page.
\item La \textbf{PAO} gère la \textbf{mise en page complète} : texte \textbf{+} images \textbf{+} couleurs. Chaque élément est placé dans un \emph{bloc} (ou \emph{cadre}) que l'on positionne librement et précisément sur la page : titre en haut, texte en deux colonnes, photo à droite, logo dans un coin.
\end{itemize}
\end{propriete}

\begin{popfigure}
\begin{center}
\begin{tikzpicture}[font=\scriptsize]
% Cadre traitement de texte
\draw[thick, popdark, fill=popcream] (0,0) rectangle (5.6,4.6);
\node[font=\small\bfseries, text=popdark] at (2.8,4.25) {Traitement de texte};
\draw[popline] (0.3,3.7) -- (5.3,3.7);
\draw[popline] (0.3,3.4) -- (5.3,3.4);
\draw[popline] (0.3,3.1) -- (4.6,3.1);
\draw[popline] (0.3,2.7) -- (5.3,2.7);
\draw[popline] (0.3,2.4) -- (5.3,2.4);
\draw[popline] (0.3,2.1) -- (5.3,2.1);
\draw[popline] (0.3,1.8) -- (3.9,1.8);
\draw[popline] (0.3,1.4) -- (5.3,1.4);
\draw[popline] (0.3,1.1) -- (5.3,1.1);
\draw[popline] (0.3,0.8) -- (5.3,0.8);
\draw[popline] (0.3,0.5) -- (4.2,0.5);
\node[text=popmuted, align=center] at (2.8,-0.35) {le texte s'écoule\\de haut en bas};
% Cadre PAO
\begin{scope}[xshift=7cm]
\draw[thick, popblue, fill=popblueL] (0,0) rectangle (5.6,4.6);
\node[font=\small\bfseries, text=popblue] at (2.8,4.25) {Mise en page PAO};
\draw[fill=poporange, draw=poporange] (0.3,3.5) rectangle (5.3,3.95);
\node[font=\scriptsize\bfseries, text=white] at (2.8,3.72) {GRAND TITRE};
\draw[fill=popgreenL, draw=popgreen] (0.3,0.4) rectangle (2.5,3.3);
\draw[popgreen] (0.45,2.9) -- (2.35,2.9);
\draw[popgreen] (0.45,2.65) -- (2.35,2.65);
\draw[popgreen] (0.45,2.4) -- (2.35,2.4);
\draw[popgreen] (0.45,2.15) -- (2.35,2.15);
\draw[popgreen] (0.45,1.9) -- (2.35,1.9);
\draw[popgreen] (0.45,1.65) -- (2.35,1.65);
\draw[popgreen] (0.45,1.4) -- (2.35,1.4);
\draw[popgreen] (0.45,1.15) -- (2.35,1.15);
\draw[popgreen] (0.45,0.9) -- (2.35,0.9);
\draw[popgreen] (0.45,0.65) -- (2.35,0.65);
\node[text=popgreen] at (1.4,3.12) {colonne 1};
\draw[fill=popgreenL, draw=popgreen] (2.7,0.4) rectangle (3.9,3.3);
\draw[popgreen] (2.85,2.9) -- (3.75,2.9);
\draw[popgreen] (2.85,2.65) -- (3.75,2.65);
\draw[popgreen] (2.85,2.4) -- (3.75,2.4);
\draw[popgreen] (2.85,2.15) -- (3.75,2.15);
\draw[popgreen] (2.85,1.9) -- (3.75,1.9);
\draw[popgreen] (2.85,1.65) -- (3.75,1.65);
\draw[popgreen] (2.85,1.4) -- (3.75,1.4);
\draw[popgreen] (2.85,1.15) -- (3.75,1.15);
\draw[popgreen] (2.85,0.9) -- (3.75,0.9);
\draw[popgreen] (2.85,0.65) -- (3.75,0.65);
\node[text=popgreen] at (3.3,3.12) {colonne 2};
\draw[fill=poppinkL, draw=poppink, thick] (4.1,1.5) rectangle (5.3,3.3);
\node[text=poppink, align=center] at (4.7,2.4) {photo};
\draw[fill=popgoldL, draw=popgold, thick] (4.1,0.4) rectangle (5.3,1.3);
\node[text=popgold] at (4.7,0.85) {logo};
\node[text=popmuted, align=center] at (2.8,-0.35) {blocs placés librement :\\colonnes, image, logo};
\end{scope}
\end{tikzpicture}
\end{center}
\legende{Le même contenu présenté en traitement de texte (à gauche) et en mise en page PAO avec colonnes, photo et logo (à droite).}
\end{popfigure}

\begin{attention}[Erreur fréquente : confondre PAO et traitement de texte]
Beaucoup d'élèves (et d'adultes !) croient que « faire de la PAO » signifie simplement taper un texte dans Word avec quelques images insérées. C'est une erreur. Certes, Word permet \emph{un peu} de mise en page (colonnes, images, zones de texte), mais ce n'est \textbf{pas} un logiciel de PAO professionnel : le placement précis des blocs y est difficile, et les éléments bougent facilement quand on modifie le texte. Un vrai logiciel de PAO est conçu dès le départ pour contrôler la position de chaque élément au millimètre près. Retiens : \textbf{le traitement de texte sert à rédiger, la PAO sert à mettre en page}.
\end{attention}

\courssub{Domaines d'application de la PAO}

La PAO est utilisée partout où l'on produit des documents imprimés ou diffusés avec une présentation soignée :

\begin{center}
\begin{tabularx}{\linewidth}{|l|X|}
\hline
\rowcolor{popblueL} \textbf{Domaine} & \textbf{Exemples de documents produits} \\
\hline
Presse & journaux, magazines, pages publicitaires \\
\hline
Publicité & affiches, flyers, banderoles, cartes de visite \\
\hline
Édition scolaire & manuels scolaires, bulletins de notes, journaux de l'établissement \\
\hline
Communication d'entreprise & brochures, catalogues, rapports annuels, en-têtes de lettres \\
\hline
Affichage administratif & avis, communiqués, programmes de cérémonies officielles \\
\hline
\end{tabularx}
\end{center}

\begin{exemplebox}[Des documents de PAO dans la vie camerounaise]
Sans le savoir, tu croises chaque jour des documents conçus en PAO :
\begin{itemize}
\item les \textbf{affiches de campagne} électorale ou les affiches de matchs et de concerts collées dans les rues de Yaoundé et de Douala ;
\item les \textbf{bulletins scolaires} et les journaux des établissements ;
\item les \textbf{menus de restaurant} plastifiés, avec photos des plats et prix en FCFA ;
\item les \textbf{faire-part} de mariage, de baptême ou de deuil, très codifiés dans nos traditions ;
\item les \textbf{tracts et bulletins d'église} distribués le dimanche.
\end{itemize}
\end{exemplebox}

\begin{exemplebox}[L'affiche du tournoi de Nkol-Eton]
Revenons à Armelle. L'affiche du tournoi doit contenir : un grand titre accrocheur, les dates et le lieu, les frais d'inscription, la photo de l'équipe championne de l'an dernier et le logo de l'association, le tout disposé \emph{librement} sur la page (le titre en haut en grandes lettres colorées, la photo au centre, les informations pratiques en bas dans un encadré). Ce placement libre et précis de texte, photo et logo est un \textbf{travail typique de PAO} : impossible à réaliser proprement avec une simple saisie de texte qui s'écoule ligne après ligne.
\end{exemplebox}

\begin{savaistu}[Les journaux camerounais et la PAO]
Les grands journaux camerounais, comme \emph{Cameroon Tribune} ou \emph{Le Messager}, sont entièrement mis en page par PAO avant impression. Chaque matin, les maquettistes placent les articles, les titres, les photos et les publicités dans des colonnes précises, puis le fichier final part à l'imprimerie. Sans la PAO, impossible de produire chaque jour des milliers d'exemplaires en quelques heures !
\end{savaistu}

\begin{exoresolu}[Traitement de texte ou PAO ?]
\textbf{Énoncé.} Pour chacune des tâches suivantes, indique si elle relève plutôt du traitement de texte ou de la PAO, en justifiant ta réponse :
\begin{enumerate}
\item rédiger une lettre de demande de stage adressée à une entreprise ;
\item concevoir l'affiche du tournoi de football de Nkol-Eton ;
\item taper un exposé de cinq pages sur les fleuves du Cameroun ;
\item réaliser le bulletin d'information de deux pages de l'association, avec des articles en colonnes et des photos.
\end{enumerate}
\tcblower
\textbf{Solution détaillée.}
\begin{enumerate}
\item \textbf{Traitement de texte.} Une lettre est essentiellement du texte organisé en paragraphes, de haut en bas, sans placement libre d'images ni de blocs.
\item \textbf{PAO.} L'affiche exige de placer librement un titre, une photo, un logo et des encadrés d'informations, avec des couleurs : c'est une mise en page complète.
\item \textbf{Traitement de texte.} Un exposé est un long texte continu, avec éventuellement des titres de paragraphes ; le texte s'écoule page après page.
\item \textbf{PAO.} Le bulletin combine articles en colonnes, photos et titres positionnés précisément sur chaque page : c'est le travail classique d'un logiciel de PAO.
\end{enumerate}
\textbf{Conclusion.} Le critère de choix est simple : si le document est dominé par le texte continu, le traitement de texte suffit ; s'il faut organiser librement texte, images et couleurs sur la page, il faut la PAO.
\end{exoresolu}

\begin{exercice}[À toi de jouer]
\begin{enumerate}
\item Définis l'infographie et la PAO avec tes propres mots, puis cite deux différences entre ces deux activités.
\item Armelle hésite : pour taper le texte des articles du bulletin, vaut-il mieux utiliser un traitement de texte ou un logiciel de PAO ? Justifie ta réponse.
\item Cite trois documents de ton quartier ou de ton établissement qui ont probablement été conçus en PAO, et explique pourquoi pour l'un d'entre eux.
\end{enumerate}
\end{exercice}

\begin{corrige}[Exercice]
\begin{enumerate}
\item L'infographie est l'ensemble des techniques de création et de manipulation d'images numériques à l'aide d'un ordinateur. La PAO est l'utilisation de l'ordinateur et de logiciels spécialisés pour concevoir des documents destinés à l'impression ou à la diffusion. Différences possibles : l'infographie porte sur les images, la PAO sur la page entière ; l'infographie produit des éléments (logo, photo retouchée) que la PAO assemble ensuite avec le texte.
\item Pour \emph{taper} les articles, le traitement de texte est adapté (saisie, correction orthographique). Mais pour \emph{mettre en page} le bulletin complet (colonnes, photos, titres), il faudra un logiciel de PAO. Les deux outils sont donc complémentaires.
\item Exemples attendus : affiche d'un match ou d'une campagne, menu d'un restaurant, faire-part de mariage, tract d'église, journal de l'établissement. Justification type : le menu du restaurant combine photos des plats, prix et cadres colorés disposés librement, ce qui exige une mise en page complète, donc la PAO.
\end{enumerate}
\end{corrige}

\begin{aretenir}[L'essentiel de la section]
\begin{itemize}
\item L'\textbf{infographie} = créer et manipuler des \textbf{images numériques} avec un ordinateur.
\item La \textbf{PAO} = concevoir avec des logiciels spécialisés des \textbf{documents destinés à l'impression ou à la diffusion} (affiches, journaux, brochures, cartes).
\item Le \textbf{traitement de texte} gère surtout le texte ; la \textbf{PAO} gère la mise en page complète : texte + images + couleurs, placés librement en blocs.
\item La PAO s'applique à la presse, la publicité, l'édition scolaire, la communication d'entreprise et l'affichage administratif.
\end{itemize}
\end{aretenir}

Armelle a maintenant sa réponse : oui, elle peut concevoir elle-même l'affiche et le bulletin, à condition de comprendre les images numériques, de choisir le bon logiciel et de respecter les règles de mise en page. C'est précisément ce que nous allons découvrir dans les sections suivantes.

\coursec{Les images numériques : matricielles et vectorielles}

Dans la section précédente, nous avons découvert que l'infographie est l'ensemble des techniques de création et de manipulation d'images par ordinateur, et que la PAO (publication assistée par ordinateur) permet de composer des documents imprimés ou diffusés : affiches, journaux scolaires, cartes de visite. Mais avant de composer un beau document, il faut comprendre ce qu'est exactement une \emph{image numérique}. Pourquoi une photo agrandie devient-elle floue ? Pourquoi le logo d'une entreprise reste-t-il net, même imprimé sur un immense panneau publicitaire ? Toute cette section répond à ces questions. Tu vas découvrir qu'il existe deux grandes familles d'images : les images \cle{matricielles} et les images \cle{vectorielles}.

\courssub{Le pixel : la brique de base de l'image numérique}

\textbf{Intuition.} Observe de très près l'écran de ton téléphone ou de ton ordinateur, par exemple en y déposant une goutte d'eau qui fait loupe : tu verras que l'image est formée d'une multitude de petits points colorés. C'est exactement comme une mosaïque ou comme les tableaux faits de petits carreaux de couleur : de loin, on voit un visage ; de très près, on ne voit que des carrés.

\begin{definition}[Pixel]
Un \cle{pixel} (contraction de l'anglais \emph{picture element}, « élément d'image ») est le \textbf{plus petit élément d'une image numérique}. C'est un petit point carré auquel on attribue une seule couleur. Une image numérique est décrite par ses dimensions en pixels : par exemple une image de $3000 \times 2000$ pixels contient $3000$ colonnes et $2000$ lignes de pixels, soit $6\,000\,000$ de pixels (6 mégapixels).
\end{definition}

\begin{exemplebox}[Compter les pixels]
L'appareil photo d'un téléphone annonce « 12 mégapixels ». Cela signifie que chaque photo contient environ 12 millions de pixels, par exemple $4000 \times 3000 = 12\,000\,000$ pixels. Plus il y a de pixels, plus l'image contient de détails... et plus le fichier est lourd !
\end{exemplebox}

\courssub{L'image matricielle (bitmap)}

\begin{definition}[Image matricielle]
Une \cle{image matricielle} (ou image \emph{bitmap}, ou image « en mode point ») est une image formée d'une \textbf{grille de pixels}, comme un tableau de nombres où chaque case contient une couleur. L'ordinateur mémorise la couleur de chaque pixel, un par un. Les formats les plus courants sont : \textbf{JPEG} (photos compressées), \textbf{PNG} (images avec transparence possible), \textbf{BMP} (non compressé, fichiers lourds) et \textbf{GIF} (petites animations, 256 couleurs maximum).
\end{definition}

\textbf{Avantage.} L'image matricielle est parfaite pour les \textbf{photographies} : comme chaque pixel a sa propre couleur, on peut reproduire fidèlement les milliers de nuances d'un visage, d'un paysage ou d'un coucher de soleil sur le mont Cameroun.

\textbf{Inconvénient.} L'image matricielle a une taille fixe en pixels. Si on l'\textbf{agrandit trop}, l'ordinateur doit « inventer » des pixels supplémentaires : l'image devient floue et on finit par voir les petits carrés. C'est l'effet de \cle{pixelisation}.

\begin{popfigure}
\begin{center}
\begin{tikzpicture}[scale=0.55]
  % Grille de pixels : un smiley simplifié
  \fill[popgoldL] (0,0) rectangle (1,1);
  \fill[popgoldL] (1,0) rectangle (2,1);
  \fill[popgoldL] (2,0) rectangle (3,1);
  \fill[popgoldL] (3,0) rectangle (4,1);
  \fill[popgoldL] (0,1) rectangle (1,2);
  \fill[popgoldL] (4,1) rectangle (5,2);
  \fill[popgoldL] (0,2) rectangle (1,3);
  \fill[popgoldL] (4,2) rectangle (5,3);
  \fill[popgoldL] (1,3) rectangle (2,4);
  \fill[popgoldL] (2,3) rectangle (3,4);
  \fill[popgoldL] (3,3) rectangle (4,4);
  % yeux
  \fill[popdark] (1,1) rectangle (2,2);
  \fill[popdark] (3,1) rectangle (4,2);
  % bouche
  \fill[poporange] (1,0) rectangle (2,1);
  \fill[poporange] (2,0) rectangle (3,1);
  \fill[poporange] (3,0) rectangle (4,1);
  \draw[popline, step=1] (0,0) grid (5,4);
  \draw[popdark, thick] (0,0) rectangle (5,4);
  \node[below, popdark] at (2.5,-0.2) {\small Chaque carré = un pixel};
  % flèche zoom
  \draw[-{Stealth}, very thick, popblue] (5.6,2) -- (7.4,2) node[midway, above]{\small zoom};
  % zoom sur 4 pixels
  \begin{scope}[shift={(8,0)}]
    \fill[popgoldL] (0,0) rectangle (2,2);
    \fill[popgoldL] (2,0) rectangle (4,2);
    \fill[popdark] (0,2) rectangle (2,4);
    \fill[popgoldL] (2,2) rectangle (4,4);
    \draw[popline, step=2] (0,0) grid (4,4);
    \draw[popdark, thick] (0,0) rectangle (4,4);
    \node[below, popdark] at (2,-0.2) {\small En zoomant, on voit les carrés};
  \end{scope}
\end{tikzpicture}
\end{center}
\legende{Une image matricielle fortement agrandie : les pixels deviennent visibles, c'est la pixelisation.}
\end{popfigure}

\courssub{L'image vectorielle}

\textbf{Intuition.} Imagine que, au lieu de mémoriser chaque point d'un dessin, on mémorise la \emph{recette} du dessin : « trace un cercle de centre $(3, 2)$ et de rayon $5$, en rouge ; trace une ligne de $(0, 0)$ à $(10, 4)$ ». Avec cette recette, on peut redessiner l'image à n'importe quelle taille, sur une carte de visite ou sur un panneau de 10 mètres : elle sera toujours parfaitement nette, car l'ordinateur recalcule le dessin à chaque fois.

\begin{definition}[Image vectorielle]
Une \cle{image vectorielle} est une image décrite par des \textbf{formules mathématiques} : coordonnées de points, équations de lignes, de courbes et de formes (cercles, rectangles, polygones), ainsi que leurs couleurs. L'ordinateur ne stocke pas des pixels, mais des instructions de dessin. Les formats courants sont : \textbf{SVG} (très utilisé sur le Web), \textbf{EPS} et \textbf{AI} (logiciel Adobe Illustrator).
\end{definition}

\textbf{Avantage.} Une image vectorielle peut être agrandie ou réduite \textbf{sans aucune perte de qualité}, puisque l'ordinateur recalcule les formules à chaque taille d'affichage. C'est le choix idéal pour les \textbf{logos}, les \textbf{schémas}, les \textbf{plans} et les \textbf{illustrations}.

\textbf{Inconvénient.} L'image vectorielle ne convient pas aux photographies : il serait impossible (et absurde) de décrire les millions de nuances d'une photo par des formules.

\begin{propriete}[Agrandissement sans perte (admise)]
Une image \textbf{vectorielle} peut être agrandie indéfiniment sans perte de qualité : l'ordinateur recalcule les formules mathématiques à chaque nouvelle taille, donc les contours restent lisses. Au contraire, une image \textbf{matricielle} perd en qualité dès qu'on l'agrandit au-delà de sa taille d'origine : les pixels deviennent visibles et les contours apparaissent crénelés (en « escalier »). C'est ce que montre la figure ci-dessous.
\end{propriete}

\begin{popfigure}
\begin{center}
\begin{tikzpicture}
  % Cercle matriciel crénelé
  \begin{scope}[shift={(0,0)}]
    \draw[fill=popblueL, draw=popdark, thick]
      (1.0,0) -- (1.4,0) -- (1.4,0.3) -- (1.7,0.3) -- (1.7,0.7)
      -- (1.9,0.7) -- (1.9,1.1) -- (1.7,1.1) -- (1.7,1.5)
      -- (1.4,1.5) -- (1.4,1.8) -- (1.0,1.8) -- (1.0,2.0)
      -- (0.6,2.0) -- (0.6,1.8) -- (0.2,1.8) -- (0.2,1.5)
      -- (-0.1,1.5) -- (-0.1,1.1) -- (-0.3,1.1) -- (-0.3,0.7)
      -- (-0.1,0.7) -- (-0.1,0.3) -- (0.2,0.3) -- (0.2,0)
      -- (0.6,0) -- (0.6,-0.2) -- (1.0,-0.2) -- cycle;
    \node[below, popdark, align=center] at (0.8,-0.5) {\small Cercle \textbf{matriciel} agrandi :\\ \small contours crénelés};
  \end{scope}
  % Cercle vectoriel lisse
  \begin{scope}[shift={(5.5,0.9)}]
    \draw[fill=popgreenL, draw=popdark, thick] (0,0) circle (1.1);
    \node[below, popdark, align=center] at (0,-1.5) {\small Cercle \textbf{vectoriel} agrandi :\\ \small contour parfaitement lisse};
  \end{scope}
\end{tikzpicture}
\end{center}
\legende{Le même cercle fortement agrandi : en matriciel les bords sont en escalier, en vectoriel ils restent lisses.}
\end{popfigure}

\courssub{La résolution d'une image}

\textbf{Intuition.} Deux photos peuvent avoir le même nombre de pixels, mais l'une s'imprimera en grand et l'autre en petit. Tout dépend de la \emph{densité} des pixels : combien de pixels place-t-on sur un centimètre de papier ? C'est la résolution.

\begin{definition}[Résolution]
La \cle{résolution} d'une image est le \textbf{nombre de pixels par unité de longueur}. Elle s'exprime en \textbf{ppp} (pixels par pouce) ou en \textbf{dpi} (\emph{dots per inch}, points par pouce). Un pouce vaut environ $2,54$ cm. Plus la résolution est élevée, plus l'image imprimée est fine et détaillée. Pour une \textbf{impression de qualité}, on recommande une résolution de \textbf{300 dpi}. Pour un simple affichage à l'écran, 72 à 96 ppp suffisent.
\end{definition}

\begin{exemplebox}[Exemple chiffré : quelle taille d'impression ?]
Une photo de $3000 \times 2000$ pixels est imprimée à $300$ dpi. Calculons la taille du tirage papier :
\begin{align*}
\text{largeur} &= \frac{3000 \text{ pixels}}{300 \text{ pixels/pouce}} = 10 \text{ pouces} = 10 \times 2,54 \approx 25,4 \text{ cm} \\
\text{hauteur} &= \frac{2000 \text{ pixels}}{300 \text{ pixels/pouce}} \approx 6,7 \text{ pouces} \approx 6,7 \times 2,54 \approx 16,9 \text{ cm}
\end{align*}
Cette photo s'imprime donc en \textbf{bonne qualité jusqu'à environ $25 \times 17$ cm}. Si on l'imprime sur une affiche de $1$ mètre de large, chaque pixel sera étiré et l'image paraîtra floue.
\end{exemplebox}

\courssub{La profondeur de couleur}

Chaque pixel mémorise sa couleur sous forme de \textbf{bits} (0 ou 1), la plus petite unité d'information en informatique. Le nombre de bits utilisés pour coder la couleur d'un pixel s'appelle la \cle{profondeur de couleur}.

\begin{center}
\begin{tabularx}{\linewidth}{|l|c|X|}
\hline
\rowcolor{popblueL} \textbf{Type d'image} & \textbf{Bits par pixel} & \textbf{Nombre de couleurs possibles} \\
\hline
Noir et blanc & 1 bit & 2 (noir ou blanc) \\
\hline
Niveaux de gris & 8 bits & $2^8 = 256$ nuances de gris \\
\hline
Couleurs vraies & 24 bits & $2^{24} = 16\,777\,216$ couleurs \\
\hline
\end{tabularx}
\end{center}

En couleurs 24 bits, chaque pixel est codé par 3 octets : un pour le rouge, un pour le vert et un pour le bleu (codage \textbf{RVB}). En mélangeant ces trois couleurs de base avec des intensités de 0 à 255, on obtient plus de 16 millions de teintes, largement de quoi reproduire toutes les couleurs qu'un œil humain peut distinguer.

\begin{exoresolu}[Taille d'une image en mémoire]
Énoncé : une image matricielle non compressée de $800 \times 600$ pixels est codée en couleurs 24 bits. Calcule sa taille en mémoire, en octets puis en kilo-octets (Ko).
\tcblower
Solution détaillée :
\begin{enumerate}
  \item Nombre de pixels : $800 \times 600 = 480\,000$ pixels.
  \item Chaque pixel occupe 24 bits, soit $24 \div 8 = 3$ octets.
  \item Taille totale : $480\,000 \times 3 = 1\,440\,000$ octets.
  \item En kilo-octets : $1\,440\,000 \div 1024 \approx 1406$ Ko, soit environ $1,4$ Mo.
\end{enumerate}
Conclusion : les images matricielles non compressées sont lourdes ; c'est pourquoi on utilise des formats compressés comme JPEG pour les photos.
\end{exoresolu}

\courssub{Bien choisir : matricielle ou vectorielle ?}

\begin{methode}[Choisir le bon type d'image]
Pour choisir entre image matricielle et image vectorielle, pose-toi la question : « quelle est la nature de mon image et que vais-je en faire ? »
\begin{enumerate}
  \item Si c'est une \textbf{photographie} (paysage, portrait, photo de classe) : choisis une image \textbf{matricielle} (JPEG ou PNG), avec suffisamment de pixels pour l'usage prévu.
  \item Si c'est un \textbf{logo}, un \textbf{schéma}, un \textbf{plan} ou une \textbf{illustration} qui devra être redimensionné (petit sur une carte de visite, grand sur une affiche) : choisis une image \textbf{vectorielle} (SVG, EPS, AI).
  \item Vérifie la \textbf{résolution} avant l'impression : vise 300 dpi pour un document imprimé de qualité.
  \item En cas de doute, crée le dessin en vectoriel : on peut toujours le convertir ensuite en image matricielle de la taille voulue, mais l'inverse donne un mauvais résultat.
\end{enumerate}
\end{methode}

\begin{attention}[Erreur fréquente : le logo qui devient flou]
L'erreur classique des débutants en PAO : on récupère un petit logo matriciel (par exemple un JPEG de $200 \times 200$ pixels trouvé sur Internet), on l'insère dans une affiche et on l'\textbf{agrandit} pour qu'il remplisse le haut de la page. Résultat : à l'impression, le logo est \textbf{flou et pixelisé}, et le document paraît peu professionnel. La bonne démarche : utiliser une version \textbf{vectorielle} du logo (SVG), ou au minimum une image matricielle de \textbf{grande taille} dès le départ. On peut réduire une grande image sans problème, mais on ne peut pas agrandir une petite image sans perte.
\end{attention}

\begin{aretenir}[L'essentiel de la section]
\begin{itemize}
  \item Le \cle{pixel} est le plus petit élément d'une image numérique : un point coloré.
  \item Une image \cle{matricielle} est une grille de pixels (JPEG, PNG, BMP, GIF) : idéale pour les photos, mais elle se \textbf{pixelise} quand on l'agrandit.
  \item Une image \cle{vectorielle} est décrite par des formules mathématiques (SVG, EPS, AI) : elle s'agrandit \textbf{sans perte de qualité}, parfaite pour les logos et schémas.
  \item La \cle{résolution} (en ppp ou dpi) indique le nombre de pixels par pouce ; \textbf{300 dpi} est recommandé pour l'impression.
  \item La \cle{profondeur de couleur} est le nombre de bits par pixel : 1 bit (noir et blanc), 8 bits (256 gris), 24 bits (16,7 millions de couleurs).
\end{itemize}
\end{aretenir}

\begin{exercice}[À toi de jouer]
\begin{enumerate}
  \item Le lycée veut imprimer une bannière de 2 mètres de large avec le logo de l'établissement. Quel type d'image faut-il utiliser pour le logo ? Justifie ta réponse.
  \item Une photo de $2400 \times 1800$ pixels est imprimée à 300 dpi. Calcule les dimensions du tirage en centimètres (rappel : 1 pouce $= 2,54$ cm).
  \item Combien de couleurs peut représenter une image codée sur 8 bits par pixel ? Cite un usage de ce type d'image.
\end{enumerate}
\end{exercice}

\begin{corrige}[Exercice : À toi de jouer]
\begin{enumerate}
  \item Il faut utiliser une image \textbf{vectorielle} (format SVG, EPS ou AI) : elle est décrite par des formules mathématiques, donc elle peut être agrandie à 2 mètres sans aucune perte de qualité. Un logo matriciel deviendrait flou et pixelisé.
  \item Largeur : $2400 \div 300 = 8$ pouces, soit $8 \times 2,54 = 20,32$ cm. Hauteur : $1800 \div 300 = 6$ pouces, soit $6 \times 2,54 = 15,24$ cm. Le tirage de bonne qualité mesure environ $20 \times 15$ cm.
  \item Avec 8 bits par pixel, on peut représenter $2^8 = 256$ couleurs (ou 256 niveaux de gris). Usage : les images en niveaux de gris (photocopies, documents scannés en noir et blanc) ou les petites images GIF.
\end{enumerate}
\end{corrige}

\coursec{Logiciels d'infographie et de PAO}

Maintenant que nous savons distinguer une image \cle{matricielle} (pixels) d'une image \cle{vectorielle} (formules géométriques), une question pratique se pose~: \emph{avec quel outil vais-je créer, modifier ou assembler ces images pour produire mon affiche, mon bulletin ou ma carte de visite~?} C'est tout l'objet de cette section~: présenter les familles de logiciels, leurs spécialités, et vous donner une méthode pour choisir le bon outil sans vous disperser.

\courssub{Les trois grandes familles de logiciels}

En infographie et en PAO, on ne travaille pas avec un seul «~couteau suisse~». On utilise généralement une \textbf{chaîne de production}~: un logiciel pour créer les éléments graphiques (logos, illustrations, retouches photos), un autre pour les assembler en mise en page finale. Voici les trois catégories principales~:

\begin{definition}[Logiciel de PAO (Publication Assistée par Ordinateur)]
Un \cle{logiciel de PAO} est une application conçue pour \textbf{composer des pages} en assemblant du texte, des images (matricielles ou vectorielles) et des formes graphiques, dans le but de produire un document prêt à l'impression ou à la diffusion numérique (PDF). Contrairement à un traitement de texte (Word, Writer), la PAO gère le positionnement \emph{absolu} des objets (au millimètre près) et la gestion des couleurs d'imprimerie (CMJN).
\end{definition}

\begin{popfigure}
\begin{center}
\begin{tikzpicture}[
    every node/.style={font=\small},
    header/.style={fill=popblue, text=white, font=\bfseries\small, align=center},
    cat/.style={fill=popblueL, font=\bfseries\small, align=center},
    soft/.style={fill=white, align=center},
    alt/.style={fill=popcream, align=center},
    border/.style={draw=popline, thin}
]
% En-têtes
\node[header, border, minimum width=3.5cm, minimum height=0.9cm] (h1) at (0,0) {Logiciel};
\node[header, border, minimum width=3cm, minimum height=0.9cm] (h2) at (3.5,0) {Catégorie};
\node[header, border, minimum width=2.5cm, minimum height=0.9cm] (h3) at (6.5,0) {Coût / Licence};
\node[header, border, minimum width=4.5cm, minimum height=0.9cm] (h4) at (9,0) {Usage principal (Exemples 3ème)};

% Ligne 1 : Paint
\node[soft, border, minimum width=3.5cm, minimum height=0.8cm] (r1c1) at (0,-0.9) {Paint (Windows)};
\node[cat, border, minimum width=3cm, minimum height=0.8cm] (r1c2) at (3.5,-0.9) {Dessin matriciel (Bitmap)};
\node[soft, border, minimum width=2.5cm, minimum height=0.8cm] (r1c3) at (6.5,-0.9) {Gratuit (inclus)};
\node[soft, border, minimum width=4.5cm, minimum height=0.8cm] (r1c4) at (9,-0.9) {Retouche ultra-simple, recadrage, gribouillage};

% Ligne 2 : GIMP
\node[alt, border, minimum width=3.5cm, minimum height=0.8cm] (r2c1) at (0,-1.7) {GIMP};
\node[cat, border, minimum width=3cm, minimum height=0.8cm] (r2c2) at (3.5,-1.7) {Dessin matriciel (Bitmap)};
\node[alt, border, minimum width=2.5cm, minimum height=0.8cm] (r2c3) at (6.5,-1.7) {Libre / Gratuit (GPL)};
\node[alt, border, minimum width=4.5cm, minimum height=0.8cm] (r2c4) at (9,-1.7) {Retouche photo pro, détourage, filtres, calques};

% Ligne 3 : Photoshop
\node[soft, border, minimum width=3.5cm, minimum height=0.8cm] (r3c1) at (0,-2.5) {Adobe Photoshop};
\node[cat, border, minimum width=3cm, minimum height=0.8cm] (r3c2) at (3.5,-2.5) {Dessin matriciel (Bitmap)};
\node[soft, border, minimum width=2.5cm, minimum height=0.8cm] (r3c3) at (6.5,-2.5) {Payant (Abonnement)};
\node[soft, border, minimum width=4.5cm, minimum height=0.8cm] (r3c4) at (9,-2.5) {Standard pro photo, design web, mockups};

% Ligne 4 : Inkscape
\node[alt, border, minimum width=3.5cm, minimum height=0.8cm] (r4c1) at (0,-3.3) {Inkscape};
\node[cat, border, minimum width=3cm, minimum height=0.8cm] (r4c2) at (3.5,-3.3) {Dessin vectoriel};
\node[alt, border, minimum width=2.5cm, minimum height=0.8cm] (r4c3) at (6.5,-3.3) {Libre / Gratuit (GPL)};
\node[alt, border, minimum width=4.5cm, minimum height=0.8cm] (r4c4) at (9,-3.3) {Logos, icônes, schémas, affiches vectorielles};

% Ligne 5 : CorelDRAW
\node[soft, border, minimum width=3.5cm, minimum height=0.8cm] (r5c1) at (0,-4.1) {CorelDRAW};
\node[cat, border, minimum width=3cm, minimum height=0.8cm] (r5c2) at (3.5,-4.1) {Dessin vectoriel (+ mise en page)};
\node[soft, border, minimum width=2.5cm, minimum height=0.8cm] (r5c3) at (6.5,-4.1) {Payant (Licence/Ab.)};
\node[soft, border, minimum width=4.5cm, minimum height=0.8cm] (r5c4) at (9,-4.1) {Design graphique pro, signalétique, imprimerie};

% Ligne 6 : Illustrator
\node[alt, border, minimum width=3.5cm, minimum height=0.8cm] (r6c1) at (0,-4.9) {Adobe Illustrator};
\node[cat, border, minimum width=3cm, minimum height=0.8cm] (r6c2) at (3.5,-4.9) {Dessin vectoriel};
\node[alt, border, minimum width=2.5cm, minimum height=0.8cm] (r6c3) at (6.5,-4.9) {Payant (Abonnement)};
\node[alt, border, minimum width=4.5cm, minimum height=0.8cm] (r6c4) at (9,-4.9) {Logos pro, typographie, illustration complexe};

% Ligne 7 : Publisher
\node[soft, border, minimum width=3.5cm, minimum height=0.8cm] (r7c1) at (0,-5.7) {Microsoft Publisher};
\node[cat, border, minimum width=3cm, minimum height=0.8cm] (r7c2) at (3.5,-5.7) {PAO (Mise en page)};
\node[soft, border, minimum width=2.5cm, minimum height=0.8cm] (r7c3) at (6.5,-5.7) {Payant (Suite Office)};
\node[soft, border, minimum width=4.5cm, minimum height=0.8cm] (r7c4) at (9,-5.7) {Bulletins, cartes visite, affiches, modèles prêts à l'emploi};

% Ligne 8 : Scribus
\node[alt, border, minimum width=3.5cm, minimum height=0.8cm] (r8c1) at (0,-6.5) {Scribus};
\node[cat, border, minimum width=3cm, minimum height=0.8cm] (r8c2) at (3.5,-6.5) {PAO (Mise en page)};
\node[alt, border, minimum width=2.5cm, minimum height=0.8cm] (r8c3) at (6.5,-6.5) {Libre / Gratuit (GPL)};
\node[alt, border, minimum width=4.5cm, minimum height=0.8cm] (r8c4) at (9,-6.5) {Journaux scolaires, livres, PDF imprimerie (CMJN)};

% Ligne 9 : InDesign
\node[soft, border, minimum width=3.5cm, minimum height=0.8cm] (r9c1) at (0,-7.3) {Adobe InDesign};
\node[cat, border, minimum width=3cm, minimum height=0.8cm] (r9c2) at (3.5,-7.3) {PAO (Mise en page)};
\node[soft, border, minimum width=2.5cm, minimum height=0.8cm] (r9c3) at (6.5,-7.3) {Payant (Abonnement)};
\node[soft, border, minimum width=4.5cm, minimum height=0.8cm] (r9c4) at (9,-7.3) {Magazines, livres, documents longs pro};

% Ligne 10 : Canva
\node[alt, border, minimum width=3.5cm, minimum height=0.8cm] (r10c1) at (0,-8.1) {Canva (Web / App)};
\node[cat, border, minimum width=3cm, minimum height=0.8cm] (r10c2) at (3.5,-8.1) {PAO en ligne (Modèles)};
\node[alt, border, minimum width=2.5cm, minimum height=0.8cm] (r10c3) at (6.5,-8.1) {Freemium (Gratuit + Pro)};
\node[alt, border, minimum width=4.5cm, minimum height=0.8cm] (r10c4) at (9,-8.1) {Réseaux sociaux, affiches rapides, collab. classe};
\end{tikzpicture}
\end{center}
\legende{Tableau comparatif des principaux logiciels d'infographie et de PAO accessibles en classe de 3ème. Les lignes sur fond \textcolor{popblueL}{bleu clair} / \textcolor{popcream}{crème} indiquent des solutions \textbf{gratuites ou libres}, idéales pour l'école et la maison.}
\end{popfigure}

\courssub{Zoom sur trois outils accessibles au collège}

Parmi cette liste, trois logiciels retiennent particulièrement notre attention en classe de 3ème car ils sont soit déjà installés sur les postes du CDI (Publisher), soit installables légalement et gratuitement (Scribus), soit accessibles sans installation depuis n'importe quel navigateur, même sur téléphone (Canva).

\courssub{Microsoft Publisher : la PAO «~clé en main~»}

\begin{popfigure}
\begin{center}
\begin{tikzpicture}[
    scale=0.85,
    every node/.style={font=\small},
    ribbon/.style={fill=popblue!15, draw=popblue, rounded corners=3pt},
    tab/.style={fill=popblue, text=white, font=\bfseries\small, rounded corners=2pt},
    tabinact/.style={fill=popblue!40, text=white, font=\small, rounded corners=2pt},
    page/.style={fill=white, draw=popmuted, thick},
    guide/.style={draw=poporange, dashed, thick},
    gallery/.style={fill=popgreenL, draw=popgreen, rounded corners=3pt}
]
% Barre de titre
\draw[fill=popdark!90, draw=popdark] (0,7.5) rectangle (16,8.2);
\node[text=white, font=\bfseries\small] at (8,7.85) {Microsoft Publisher - [Bulletin\_Association.pub] - Enregistré};

% Ruban - Onglets
\draw[ribbon] (0,6.5) rectangle (16,7.4);
\node[tab] at (1.2,6.95) {Fichier};
\node[tabinact] at (3.2,6.95) {Accueil};
\node[tabinact] at (5.0,6.95) {Insertion};
\node[tabinact] at (6.8,6.95) {Mise en page};
\node[tabinact] at (8.8,6.95) {Mise en forme du texte};
\node[tabinact] at (11.2,6.95) {Outils d'image};
\node[tabinact] at (13.2,6.95) {Affichage};

% Zone de travail - Gauche : Galerie modèles / Pages
\draw[gallery] (0.2,0.5) rectangle (3.8,6.3);
\node[font=\bfseries\small, text=popgreen!80!black] at (2,6.0) {Galerie de modèles};
\node[draw=popline, fill=white, rounded corners=2pt, minimum width=3.2cm, minimum height=1.8cm] (mod1) at (2,4.8) {};
\node[align=center, font=\tiny] at (2,4.8) {Affiche\\Événement};
\node[draw=popline, fill=white, rounded corners=2pt, minimum width=3.2cm, minimum height=1.8cm] (mod2) at (2,3.0) {};
\node[align=center, font=\tiny] at (2,3.0) {Brochure\\3 volets};
\node[draw=popline, fill=white, rounded corners=2pt, minimum width=3.2cm, minimum height=1.8cm] (mod3) at (2,1.2) {};
\node[align=center, font=\tiny] at (2,1.2) {Carte de\\visite};

% Zone centrale : La Page
\draw[page] (4.2,0.5) rectangle (13.8,6.3);
% Guides de marge
\draw[guide] (4.8,0.5) -- (4.8,6.3);
\draw[guide] (13.2,0.5) -- (13.2,6.3);
\draw[guide] (4.2,5.7) -- (13.8,5.7);
\draw[guide] (4.2,1.1) -- (13.8,1.1);
% Contenu exemple
\node[draw=popblue, fill=popblueL, rounded corners=2pt, minimum width=8cm, minimum height=1.2cm] (title) at (9,5.8) {};
\node[font=\bfseries\small, text=popblue] at (9,5.8) {Titre : Bulletin de l'Association des Élèves};
\node[draw=popmuted, fill=popcream, rounded corners=2pt, minimum width=4cm, minimum height=2.5cm] (img) at (6.5,3.8) {};
\node[font=\tiny, align=center] at (6.5,3.8) {Zone image\\ (Photo club)};
\node[draw=popmuted, fill=popcream, rounded corners=2pt, minimum width=4cm, minimum height=2.5cm] (txt) at (11.2,3.8) {};
\node[font=\tiny, align=center] at (11.2,3.8) {Zone de texte\\ (Articles)};
\node[draw=poporange, fill=poporangeL, rounded corners=2pt, minimum width=8cm, minimum height=0.8cm] (foot) at (9,1.4) {};
\node[font=\tiny, text=poporange] at (9,1.4) {Pied de page : Contact / Mentions légales};

% Barre d'état
\draw[fill=popmuted!20, draw=popmuted] (0,0) rectangle (16,0.4);
\node[font=\tiny] at (2,0.2) {Page 1 sur 2};
\node[font=\tiny] at (8,0.2) {100\%};
\node[font=\tiny] at (14,0.2) {Français (Cameroun)};
\end{tikzpicture}
\end{center}
\legende{Schéma simplifié de l'interface de Microsoft Publisher : le \textbf{Ruban} (haut), la \textbf{Galerie de modèles} (gauche), la \textbf{Zone de page} avec guides de marge (centre).}
\end{popfigure}

\textbf{Microsoft Publisher} fait partie de la suite Office. Son atout majeur pour un élève de 3ème~: les \textbf{modèles prédéfinis} (\cle{templates}). Dès l'ouverture, une galerie propose des structures toutes faites pour «~Affiche~», «~Brochure~», «~Carte de visite~», «~Calendrier~», «~Bulletin d'information~». Vous choisissez un modèle, vous remplacez les textes et images par les vôtres, et le document est prêt.

\begin{aretenir}[Points forts de Publisher pour débuter]
\begin{itemize}
    \item Interface familière (ruban style Word/Excel).
    \item \textbf{Modèles intégrés} très variés (plus de 100 catégories).
    \item Gestion aisée des \textbf{blocs de texte liés} (le texte coule d'un cadre à l'autre automatiquement).
    \item Export natif en PDF haute qualité pour l'impression.
\end{itemize}
\end{aretenir}

\begin{attention}[Limites de Publisher]
\begin{itemize}
    \item \textbf{Payant} (licence Microsoft 365 ou Office perpétuelle). Pas installable légalement à la maison sans licence.
    \item Fichiers \texttt{.pub} \textbf{propriétaires} : illisibles sur Linux, Mac, ou sans Publisher. Pensez à exporter en \texttt{.pdf} pour partager.
    \item Pas d'outils de dessin vectoriel avancé (pas de plume de Bézier, pas de fusion de formes complexes). On \emph{importe} les logos faits ailleurs (Inkscape).
\end{itemize}
\end{attention}

\courssub{Scribus : la PAO libre et gratuite pour l'école}

\begin{savaistu}[Scribus et GIMP : la liberté légale]
Scribus (PAO) et GIMP (retouche photo) sont publiés sous licence \textbf{GPL (General Public License)}. Cela signifie qu'une école, un cybercafé, une association ou un élève a le \textbf{droit légal} de les télécharger, les installer sur \emph{tous} les ordinateurs, les copier sur clé USB, les modifier, sans payer un centime de licence. C'est la solution idéale pour équiper un CDI au Cameroun sans budget logiciel.
\end{savaistu}

\textbf{Scribus} est un logiciel de PAO \textbf{libre, multiplateforme} (Windows, Linux, macOS). Il est plus «~pro~» que Publisher~: il gère nativement les couleurs \textbf{CMJN} (quadrichromie imprimerie), les \textbf{fonds perdus} (bleed), les \textbf{repères de coupe}, les styles de paragraphes/caractères rigoureux, et génère des PDF conformes aux normes d'imprimerie (PDF/X-1a, PDF/X-3).

\begin{exemplebox}[Pour le journal du collège : Scribus s'impose]
L'équipe du journal scolaire veut imprimer 200 exemplaires chez un imprimeur local (Yaoundé, Douala, Bafoussam). L'imprimeur exige un PDF «~prêt à graver~» avec fonds perdus de 3~mm et repères de coupe. Publisher ne gère pas les fonds perdus nativement. Scribus, si. L'équipe installe Scribus sur les postes du CDI (Linux ou Windows), met en page le journal, exporte en PDF/X-1a avec repères : l'imprimeur valide du premier coup.
\end{exemplebox}

\begin{attention}[Courbe d'apprentissage Scribus]
Scribus est \textbf{moins intuitif} que Publisher. Pas de galerie de modèles «~clé en main~» au démarrage. Il faut créer ses gabarits (pages maîtresses), importer ses images, lier ses cadres de texte manuellement. C'est formateur (on comprend \emph{comment} marche la PAO), mais plus lent pour un premier document. Prévoyez une séance de prise en main.
\end{attention}

\courssub{Canva : la PAO collaborative dans le navigateur}

\textbf{Canva} (\texttt{canva.com}) change la donne~: plus besoin d'installer un logiciel. On se connecte avec un compte Google, Microsoft ou email, et on travaille \textbf{en ligne}, sur ordinateur \emph{ou} sur smartphone (application Android/iOS). C'est un atout majeur quand le CDI est complet ou pour finir le travail à la maison sur son téléphone.

\begin{aretenir}[Pourquoi Canva séduit les élèves]
\begin{itemize}
    \item \textbf{Zéro installation}, synchronisation cloud automatique.
    \item \textbf{Bibliothèque géante de modèles} (recherche par mots-clés : «~affiche rentrée scolaire~», «~carte anniversaire~»).
    \item \textbf{Collaboration temps réel} : plusieurs élèves éditent le même design en même temps (lien de partage).
    \item Éléments graphiques (photos, icônes, illustrations) intégrés, souvent gratuits.
    \item Export PDF (standard ou impression), PNG, JPG, MP4 (vidéo), GIF.
\end{itemize}
\end{aretenir}

\begin{attention}[Vigilance Canva : Freemium et Données]
\begin{itemize}
    \item Modèle \textbf{Freemium} : beaucoup d'éléments «~Pro~» (couronne) sont payants. Filtrez sur «~Gratuit~» dans la recherche.
    \item \textbf{Données personnelles} : création de compte obligatoire. Vérifiez la politique de confidentialité. Évitez d'y mettre des données sensibles (notes, numéros CNI).
    \item \textbf{Hors-ligne impossible} : pas de connexion = pas de Canva. Prévoyez une alternative locale (Publisher/Scribus) pour les zones sans réseau stable.
    \item \textbf{Précision limitée} : pas de contrôle millimétrique des fonds perdus, pas de CMJN natif (RVB seulement). Pour une impression pro chez l'imprimeur, Scribus reste supérieur.
\end{itemize}
\end{attention}

\courssub{Comment choisir~? Méthode et tableau de décision}

Ne cherchez pas «~le meilleur logiciel~». Cherchez \textbf{le bon outil pour VOTRE besoin aujourd'hui}. Voici une méthode en 4~questions.

\begin{methode}[Choisir son logiciel d'infographie/PAO en 4 questions]
\begin{enumerate}
    \item \textbf{Quel est le produit final~?}
    \begin{itemize}
        \item Logo, icône, schéma technique, affiche vectorielle $\rightarrow$ \textbf{Dessin vectoriel} (Inkscape).
        \item Retouche photo, recadrage, filtres, montage pixel $\rightarrow$ \textbf{Dessin matriciel} (GIMP, Paint).
        \item Document multipage, texte + images, prêt à imprimer (bulletin, journal, rapport, carte visite) $\rightarrow$ \textbf{PAO} (Publisher, Scribus, Canva).
    \end{itemize}
    \item \textbf{Quel environnement matériel/logiciel~?}
    \begin{itemize}
        \item Poste CDI Windows + Office installé $\rightarrow$ Publisher dispo immédiatement.
        \item Poste Linux / Windows sans licence / À la maison sans payer $\rightarrow$ Scribus (installable), GIMP, Inkscape.
        \item Téléphone / Tablette / Navigateur seulement / Travail en groupe à distance $\rightarrow$ Canva.
    \end{itemize}
    \item \textbf{Quel niveau de précision / contrainte d'impression~?}
    \begin{itemize}
        \item Imprimeur pro (fonds perdus, CMJN, repères, PDF/X) $\rightarrow$ \textbf{Scribus} (ou InDesign pro).
        \item Impression scolaire (photocopieuse, PDF standard) $\rightarrow$ Publisher, Canva (PDF Impression).
    \end{itemize}
    \item \textbf{Travaillerez-vous seul ou à plusieurs~?}
    \begin{itemize}
        \item Seul $\rightarrow$ Tous conviennent.
        \item Groupe (collaboration simultanée) $\rightarrow$ \textbf{Canva} (partage lien édition), ou fichiers partagés sur clé/Drive pour Publisher/Scribus (mais pas d'édition simultanée).
    \end{itemize}
\end{enumerate}
\end{methode}

\begin{popfigure}
\begin{center}
\begin{tikzpicture}[
    every node/.style={font=\small, align=center},
    header/.style={fill=poppurple, text=white, font=\bfseries\small},
    yes/.style={fill=popgreenL, text=popgreen!80!black, font=\bfseries\small},
    no/.style={fill=poppinkL, text=poppink!80!black, font=\bfseries\small},
    info/.style={fill=popblueL, text=popblue!80!black, font=\small},
    cell/.style={draw=popline, minimum height=0.75cm, anchor=center},
    headercell/.style={cell, header, minimum height=0.9cm}
]
% En-têtes
\node[headercell, minimum width=3.2cm] (h0) at (0,0) {Critère de décision};
\node[headercell, minimum width=2.8cm] (h1) at (3.2,0) {Inkscape (Vectoriel)};
\node[headercell, minimum width=2.8cm] (h2) at (6.0,0) {GIMP (Matriciel)};
\node[headercell, minimum width=2.8cm] (h3) at (8.8,0) {Publisher (PAO)};
\node[headercell, minimum width=2.8cm] (h4) at (11.6,0) {Scribus (PAO)};
\node[headercell, minimum width=2.8cm] (h5) at (14.4,0) {Canva (PAO Web)};

% Lignes
% 1 Logo / Illustration vectorielle
\node[cell, info, minimum width=3.2cm] at (0,-0.85) {Logo, icône, schéma, affiche vectorielle};
\node[cell, yes, minimum width=2.8cm] at (3.2,-0.85) {\textbf{Idéal}};
\node[cell, no, minimum width=2.8cm] at (6.0,-0.85) {Inadapté};
\node[cell, no, minimum width=2.8cm] at (8.8,-0.85) {Import seulement};
\node[cell, no, minimum width=2.8cm] at (11.6,-0.85) {Import seulement};
\node[cell, no, minimum width=2.8cm] at (14.4,-0.85) {Limité (formes simples)};

% 2 Retouche Photo
\node[cell, info, minimum width=3.2cm] at (0,-1.7) {Retouche photo, filtres, détourage};
\node[cell, no, minimum width=2.8cm] at (3.2,-1.7) {Inadapté};
\node[cell, yes, minimum width=2.8cm] at (6.0,-1.7) {\textbf{Idéal}};
\node[cell, no, minimum width=2.8cm] at (8.8,-1.7) {Import seulement};
\node[cell, no, minimum width=2.8cm] at (11.6,-1.7) {Import seulement};
\node[cell, no, minimum width=2.8cm] at (14.4,-1.7) {Basique (filtres)};

% 3 Mise en page multipage (bulletin, journal)
\node[cell, info, minimum width=3.2cm] at (0,-2.55) {Bulletin, journal, rapport, carte visite (texte + images)};
\node[cell, no, minimum width=2.8cm] at (3.2,-2.55) {Difficile (1 page)};
\node[cell, no, minimum width=2.8cm] at (6.0,-2.55) {Non};
\node[cell, yes, minimum width=2.8cm] at (8.8,-2.55) {\textbf{Idéal (Modèles)}};
\node[cell, yes, minimum width=2.8cm] at (11.6,-2.55) {\textbf{Idéal (Pro)}};
\node[cell, yes, minimum width=2.8cm] at (14.4,-2.55) {\textbf{Idéal (Modèles)}};

% 4 Imprimeur pro (CMJN, Fonds perdus, PDF/X)
\node[cell, info, minimum width=3.2cm] at (0,-3.4) {Impression pro (CMJN, fonds perdus 3mm, PDF/X)};
\node[cell, no, minimum width=2.8cm] at (3.2,-3.4) {Non (RVB seulement)};
\node[cell, no, minimum width=2.8cm] at (6.0,-3.4) {Non (RVB seulement)};
\node[cell, no, minimum width=2.8cm] at (8.8,-3.4) {Non natif};
\node[cell, yes, minimum width=2.8cm] at (11.6,-3.4) {\textbf{Natif}};
\node[cell, no, minimum width=2.8cm] at (14.4,-3.4) {RVB seulement};

% 5 Coût / Licence École
\node[cell, info, minimum width=3.2cm] at (0,-4.25) {Gratuit, légal, installable partout (CDI, maison)};
\node[cell, yes, minimum width=2.8cm] at (3.2,-4.25) {\textbf{Oui (Libre)}};
\node[cell, yes, minimum width=2.8cm] at (6.0,-4.25) {\textbf{Oui (Libre)}};
\node[cell, no, minimum width=2.8cm] at (8.8,-4.25) {Payant (Office)};
\node[cell, yes, minimum width=2.8cm] at (11.6,-4.25) {\textbf{Oui (Libre)}};
\node[cell, yes, minimum width=2.8cm] at (14.4,-4.25) {Freemium (Compte requis)};

% 6 Collaboration temps réel
\node[cell, info, minimum width=3.2cm] at (0,-5.1) {Travail en groupe simultané (même doc)};
\node[cell, no, minimum width=2.8cm] at (3.2,-5.1) {Non (fichier local)};
\node[cell, no, minimum width=2.8cm] at (6.0,-5.1) {Non (fichier local)};
\node[cell, no, minimum width=2.8cm] at (8.8,-5.1) {Non (fichier local)};
\node[cell, no, minimum width=2.8cm] at (11.6,-5.1) {Non (fichier local)};
\node[cell, yes, minimum width=2.8cm] at (14.4,-5.1) {\textbf{Oui (Cloud)}};

% 7 Prise en main (Facilité 3ème)
\node[cell, info, minimum width=3.2cm] at (0,-5.95) {Facilité prise en main (élève 3ème, 1ère fois)};
\node[cell, info, minimum width=2.8cm] at (3.2,-5.95) {Moyenne (Outils techniques)};
\node[cell, info, minimum width=2.8cm] at (6.0,-5.95) {Moyenne (Interface dense)};
\node[cell, yes, minimum width=2.8cm] at (8.8,-5.95) {\textbf{Facile (Ruban connu)}};
\node[cell, no, minimum width=2.8cm] at (11.6,-5.95) {Difficile (Concepts pro)};
\node[cell, yes, minimum width=2.8cm] at (14.4,-5.95) {\textbf{Très facile (Glisser-déposer)}};
\end{tikzpicture}
\end{center}
\legende{Tableau de décision : à quelle famille de logiciel correspond VOTRE besoin~? Vert = Choix recommandé. Rose = Déconseillé pour cet usage. Bleu = Critère.}
\end{popfigure}

\courssub{La notion de Modèle (Template) : ne réinventez pas la roue}

\begin{definition}[Modèle / Template (PAO)]
Un \cle{modèle} (ou \cle{template}, fichier \texttt{.pub}, \texttt{.sla}, format Canva) est un \textbf{document préconçu} contenant la structure de mise en page~: gabarits de pages (marges, colonnes, en-têtes/pieds de page), styles de texte (titres, corps, citations), cadres vides pour images, palette de couleurs, polices associées. L'utilisateur \emph{personnalise} le modèle (remplace textes/images, adapte couleurs) au lieu de partir d'une page blanche.
\end{definition}

\begin{propriete}[Avantages du modèle pour un élève]
\begin{itemize}
    \item \textbf{Gain de temps} : la structure est valide (marges, alignements, hiérarchie visuelle).
    \item \textbf{Cohérence} : styles de paragraphes prédéfinis $\rightarrow$ titres uniformes automatiquement.
    \item \textbf{Apprentissage} : en déverrouillant les calques, on \emph{voit} comment un pro a construit la page (grille, alignements, espace blanc).
    \item \textbf{Réutilisabilité} : on enregistre \emph{son} document personnalisé comme \emph{nouveau modèle} pour l'année prochaine.
\end{itemize}
\end{propriete}

\begin{exemplebox}[Armelle et le bulletin de l'association]
\textbf{Situation~:} Armelle (3ème) est secrétaire de l'Association des Élèves Délégués. Elle doit produire le bulletin trimestriel (4 pages A4, plié en 2 = 8 pages A5) pour vendredi. L'imprimeur du quartier (M.~Fouda, rue de l'École) imprime sur photocopieuse couleur A3 plié/agrafé.

\textbf{Analyse d'Armelle (méthode 4 questions)~:}
\begin{enumerate}
    \item \textit{Produit final~:} Bulletin multipage, texte + photos. $\rightarrow$ \textbf{PAO}.
    \item \textit{Environnement~:} CDI a Publisher (Office 365). Elle a un téléphone Android. Pas de PC à la maison.
    \item \textit{Précision impression~:} Photocopieuse école/quartier $\rightarrow$ PDF standard suffit. Pas de CMJN/fonds perdus stricts.
    \item \textit{Collaboration~:} Elle travaille seule sur la mise en page, mais récupère les articles (Word) et photos (WhatsApp) des autres.
\end{enumerate}

\textbf{Choix~:} \textbf{Microsoft Publisher} au CDI.
\textbf{Justification~:} Modèle «~Bulletin d'information 4 pages~» dispo dans la galerie. Elle glisse-dépose les textes Word dans les cadres liés (le texte coule tout seul), insère les photos reçues, adapte les couleurs aux codes de l'asso (vert/jaune). Exporte en \texttt{PDF Impression haute qualité}. Dépose le PDF sur la clé USB de M.~Fouda. Fait en 1h30 au lieu de 4h depuis page blanche.

\textbf{Alternative si CDI fermé~:} \textbf{Canva} sur son téléphone (modèle «~Newsletter scolaire~»), même logique, export PDF. Moins précis pour le pliage, mais dépannage parfait.
\end{exemplebox}

\begin{experience}[TP Guidé : «~Ma première carte de visite~» (Publisher ou Canva)]
\textbf{Objectif~:} Créer une carte de visite personnelle (format standard 85$\times$55~mm) en utilisant un modèle, en 20~min.

\textbf{Matériel~:} Poste Windows (Publisher) \textbf{OU} Navigateur (Canva).

\textbf{Consignes~:}
\begin{enumerate}
    \item \textbf{Ouvrir le logiciel} et chercher «~Carte de visite~» dans les modèles (Publisher : \texttt{Fichier > Nouveau > Cartes de visite} ; Canva : barre recherche «~Business card~» + filtre «~Gratuit~»).
    \item \textbf{Choisir un modèle} sobre (nom, téléphone, email, classe, une icône loisir). Éviter les modèles surchargés.
    \item \textbf{Remplacer les textes~:} Votre Prénom NOM, «~Élève 3ème~», \texttt{prenom.nom@email.com}, \texttt{+237 6XX XX XX XX}. Police lisible (Montserrat, Roboto, Open Sans). Taille mini 9~pt pour l'info, 14~pt pour le nom.
    \item \textbf{Insérer une image~:} Icône vectorielle (Inkscape $\rightarrow$ export PNG transparent) ou photo carrée recadrée (GIMP/Paint). Placer dans le cadre prévu.
    \item \textbf{Vérifier les marges~:} Aucun texte/logo à moins de 3~mm du bord (zone de coupe).
    \item \textbf{Exporter~:} \texttt{Fichier > Exporter > PDF (Impression)}. Vérifier l'aperçu : recto/verso correct ? Texte net ?
    \item \textbf{Enregistrer le fichier source~:} \texttt{.pub} (Publisher) ou sauvegarde auto (Canva) pour modification future.
\end{enumerate}

\textbf{Observations attendues~:} Les élèves constatent que le modèle gère l'alignement (grille invisible), les styles (couleurs cohérentes), le format physique. Ils comprennent l'intérêt du \texttt{.pdf} comme format de livraison universel.

\textbf{Interprétation~:} La PAO, c'est \textbf{assembler} dans un gabarit rigoureux, pas «~dessiner~» au pixel près. Le modèle est l'ami du débutant.
\end{experience}

\begin{exoresolu}[Choix argumenté pour l'affiche de la kermesse]
\textbf{Énoncé~:} L'association des parents d'élèves (APE) organise la kermesse de fin d'année. Ils ont besoin d'une affiche A3 (297$\times$420~mm) à afficher sur les grilles de l'école et chez les commerçants du quartier. L'affiche contient : le logo de l'école (fichier \texttt{.ai} vectoriel fourni par l'admin), une photo de la kermesse précédente (fichier \texttt{.jpg} 4~Mo), les infos pratiques (date, heure, lieu, contact), et un QR code vers le formulaire d'inscription en ligne. L'impression se fera sur la photocopieuse A3 de l'école (papier 160~g). Quel(s) logiciel(s) choisir et dans quel ordre~?

\tcblower

\textbf{Solution détaillée~:}
\begin{enumerate}
    \item \textbf{Analyse des éléments sources~:}
    \begin{itemize}
        \item Logo école \texttt{.ai} (vectoriel) $\rightarrow$ doit rester vectoriel pour netteté à n'importe quelle taille.
        \item Photo kermesse \texttt{.jpg} (matricielle) $\rightarrow$ vérifier résolution : 4~Mo pour A3 $\approx$ 150~DPI (acceptable pour photocopieuse, idéal 300~DPI). Peut-être retoucher luminosité/contraste.
        \item QR code $\rightarrow$ générer en vectoriel (SVG) si possible, sinon PNG haute résolution.
        \item Texte infos $\rightarrow$ mise en page typographique.
    \end{itemize}
    \item \textbf{Chaîne de production recommandée~:}
    \begin{enumerate}
        \item \textbf{Inkscape} (ou Illustrator si dispo) : Ouvrir le \texttt{.ai}, vérifier/convertir en \texttt{.svg} ou \texttt{.pdf} vectoriel. Générer le QR code en SVG (site \texttt{qr-code-generator.com} $\rightarrow$ download SVG). \textbf{Sortie~:} \texttt{logo\_ecole.svg}, \texttt{qrcode\_inscription.svg}.
        \item \textbf{GIMP} : Ouvrir \texttt{photo\_kermesse.jpg}. \texttt{Couleurs > Niveaux} (auto ou manuel) pour éclaircir si sombre. \texttt{Image > Échelle et taille de l'image} : vérifier 300~DPI pour 297~mm largeur $\approx$ 3500~px. Si moins, ne pas upscaler artificiellement (perte qualité), accepter 150~DPI. Exporter \texttt{photo\_kermesse\_prete.jpg} (JPEG qualité 90\%).
        \item \textbf{PAO (Choix selon dispo)~:}
        \begin{itemize}
            \item \textbf{CDI (Publisher)~:} Nouveau document A3. Insérer \texttt{logo\_ecole.svg} (Publisher importe SVG), \texttt{photo\_kermesse\_prete.jpg}, \texttt{qrcode\_inscription.svg}. Créer cadres texte pour infos. Aligner sur guides. Exporter \texttt{PDF Impression}.
            \item \textbf{Maison (Canva)~:} Modèle «~Affiche événement A3~». Upload SVG (Canva gère SVG). Même assemblage. Télécharger \texttt{PDF Impression}.
            \item \textbf{CDI Linux / Sans licence (Scribus)~:} Nouveau document A3, marges 10~mm, fonds perdus 3~mm (même si photocopieuse ne coupe pas, bonne habitude). Importer SVG (Fichier > Importer > Vectoriel), images, cadres texte. Styles de paragraphe pour titre/sous-titre/texte. Exporter \texttt{PDF/X-1a} (robuste).
        \end{itemize}
    \end{enumerate}
    \item \textbf{Justification finale~:} On a utilisé \textbf{le bon outil pour chaque type d'image} (Vectoriel pour logo/QR, Matriciel pour photo) et \textbf{un outil de PAO pour l'assemblage final}. C'est le flux de travail professionnel standard, adapté au niveau 3ème.
\end{enumerate}
\end{exoresolu}

\begin{attention}[Erreur classique : «~Je fais tout dans Paint / Word / PowerPoint~»]
\begin{itemize}
    \item \textbf{Paint} : pas de gestion du texte multiligne, pas de calques, résolution figée, pas d'import SVG. \textbf{Interdit} pour une affiche A3.
    \item \textbf{Word / Writer} : Traitement de texte, flux de texte continu. Positionner une image «~au millimètre~» en haut à gauche puis un texte juste à côté est un cauchemar (ancres, repérage paragraphe). Pas de fonds perdus, RVB seulement.
    \item \textbf{PowerPoint / Impress} : Mieux que Word (positionnement libre), mais format diapo (16/9 ou 4/3), pas A3/A4 natif sans bidouillage. Gestion des couleurs RVB seulement. Export PDF basique. \textbf{Acceptable pour un exposé oral projeté, pas pour une affiche imprimée.}
\end{itemize}
\textbf{Règle d'or~:} \textbf{Dessin vectoriel (logo) + Retouche matricielle (photo) + PAO (assemblage final).} Trois outils, trois métiers.
\end{attention}

\begin{aretenir}[Bilan : Votre boîte à outils 3ème]
\begin{center}
\begin{tabularx}{\linewidth}{|l|X|X|}\hline
\textbf{Besoin} & \textbf{Outil recommandé (Gratuit/Libre)} & \textbf{Outil alternatif (CDI/Office)} \\\hline
Créer un logo, icône, schéma propre & \textbf{Inkscape} (Installez-le !) & --- \\\hline
Retoucher une photo (luminosité, recadrage, filtre) & \textbf{GIMP} (Installez-le !) & Paint (basique) \\\hline
Monter l'affiche / bulletin / carte (Mise en page) & \textbf{Scribus} (Pro) ou \textbf{Canva} (Web) & \textbf{Publisher} (Modèles faciles) \\\hline
Format de livraison final pour l'imprimeur / prof & \textbf{PDF (Impression / PDF/X-1a si Scribus)} & \textbf{PDF} \\\hline
\end{tabularx}
\end{center}
\end{aretenir}

\courssub{Exercices d'application}

\begin{exercice}[Choix d'outils et justifications]
Pour chacun des besoins suivants, indiquez le(s) logiciel(s) le(s) plus adapté(s) parmi \{Paint, GIMP, Inkscape, Publisher, Scribus, Canva, PowerPoint\} et justifiez en une phrase (type d'image, contrainte environnement, contrainte sortie).
\begin{enumerate}
    \item Le logo du club «~Robotique~» (forme géométrique, texte, doit être net sur badge 3~cm et banderole 3~m).
    \item La photo de classe (floue, trop sombre) pour l'annuaire de fin d'année (impression offset pro).
    \item L'invitation pliée (A5 plié, 4 pages) pour la remise des prix, à imprimer demain sur la photocopieuse de l'école. Vous êtes au CDI.
    \item La même invitation, mais vous êtes chez vous dimanche soir, sur votre téléphone, pas de PC.
    \item Un schéma électrique (composants, fils) pour le rapport de TP Physique, à insérer dans un document Writer.
\end{enumerate}
\end{exercice}

\begin{corrige}[Exercice n°1]
\begin{enumerate}
    \item \textbf{Inkscape}. Logo = vectoriel (net à toute taille). Inkscape est l'outil vectoriel libre de référence.
    \item \textbf{GIMP} (retouche) $\rightarrow$ \textbf{Scribus} (mise en page annuaire). Photo = matricielle. GIMP pour corriger flou/niveaux. Scribus pour gestion CMJN/fonds perdus exigés par l'imprimeur offset.
    \item \textbf{Publisher}. Invitation multipage + CDI + Office + urgence + photocopieuse = Modèle Publisher «~Carte pliée~» le plus rapide.
    \item \textbf{Canva}. Téléphone + Web + Urgence + Modèles «~Invitation~» + Export PDF = Seule solution viable sur mobile sans installation.
    \item \textbf{Inkscape}. Schéma technique = vectoriel (lignes nettes, alignement grille, export SVG/PDF vectoriel pour insertion propre dans Writer). PowerPoint possible mais moins précis pour les connexions fils/composants.
\end{enumerate}
\end{corrige}

\coursec{Les éléments d'une mise en page}

Après avoir découvert les logiciels qui nous permettent de faire de la PAO (\emph{Publication Assistée par Ordinateur}), nous devons maintenant comprendre \emph{comment} organiser une page. Avoir un outil puissant comme Scribus, Publisher ou Canva ne suffit pas : si l'on place les éléments n'importe comment, le document sera illisible, laid et inefficace. La mise en page, c'est l'art d'organiser l'espace pour guider l'œil du lecteur, hiérarchiser l'information et produire un document professionnel, qu'il s'agisse d'un bulletin de notes, d'une affiche pour la fête de l'école ou d'un rapport de stage.

\courssub{Format de page et orientation}

Avant même d'écrire le premier mot, il faut choisir le \textbf{format} (la taille physique de la feuille) et son \textbf{orientation}.

\begin{definition}[Format et orientation de page]
Le \cle{format} définit les dimensions normalisées de la feuille (série A ISO 216). Les plus courants au collège et en administration sont :
\begin{itemize}
    \item \textbf{A4} (210 × 297 mm) : standard pour les lettres, rapports, bulletins, examens.
    \item \textbf{A5} (148 × 210 mm) : demi-A4, \textbf{utilisé pour} les livrets, programmes de fête, flyers.
    \item \textbf{A3} (297 × 420 mm) : double A4, \textbf{utilisé pour} les affiches, grands tableaux, plans.
\end{itemize}
L'\cle{orientation} détermine le sens de lecture :
\begin{itemize}
    \item \textbf{Portrait} (hauteur > largeur) : lecture naturelle pour les textes longs (rapports, courriers).
    \item \textbf{Paysage} (largeur > hauteur) : adapté aux tableaux larges, frises chronologiques, affiches panoramiques.
\end{itemize}
\end{definition}

\begin{methode}[Choisir le format et l'orientation]
\begin{enumerate}
    \item \textbf{Analyser le contenu :} Beaucoup de texte continu $\rightarrow$ A4 Portrait. Tableau large ou affiche visuelle $\rightarrow$ A3 Paysage. Petit livret $\rightarrow$ A5.
    \item \textbf{Vérifier l'imprimante :} La plupart des imprimantes scolaires (cybercafés, secrétariats) n'acceptent que l'A4. L'A3 nécessite souvent une imprimante spécifique ou un pliage.
    \item \textbf{Fixer le format \emph{avant} de commencer :} Changer le format en cours de route décale tous les blocs et casse la mise en page.
\end{enumerate}
\end{methode}

\begin{popfigure}
\begin{tikzpicture}[scale=0.38, every node/.style={font=\small}]
    % Page A4 frame
    \draw[popdark, line width=2pt] (0,0) rectangle (21,29.7);
    % Marges (15mm)
    \draw[popblue!50, dashed, line width=1pt] (1.5,1.5) rectangle (19.5,28.2);
    % Marge labels
    \node[popblue, rotate=90] at (0.7,15) {Marge gauche (gouttière)};
    \node[popblue] at (10.5,0.7) {Marge haute};
    \node[popblue, rotate=-90] at (20.3,15) {Marge droite};
    \node[popblue] at (10.5,29) {Marge basse};
    
    % En-tête zone
    \draw[poporange, line width=1.5pt] (1.5,26.5) rectangle (19.5,28.2);
    \node[poporange, align=center] at (10.5,27.35) {\textbf{EN-TÊTE}\newline Logo | Titre doc | Date};
    
    % Pied de page zone
    \draw[poporange, line width=1.5pt] (1.5,1.5) rectangle (19.5,3.2);
    \node[poporange, align=center] at (10.5,2.35) {\textbf{PIED DE PAGE}\newline N° page | Mentions légales | Contact};
    
    % Corps : 2 colonnes
    \draw[popgreen, line width=1.5pt] (1.5,3.5) rectangle (10.2,26.2);
    \node[popgreen, align=center] at (5.85,14.85) {\textbf{COLONNE 1}\newline Bloc texte principal};
    
    \draw[popgreen, line width=1.5pt] (10.8,3.5) rectangle (19.5,26.2);
    \node[popgreen, align=center] at (15.15,14.85) {\textbf{COLONNE 2}\newline Bloc texte / Image};
    
    % Gouttière entre colonnes
    \draw[popmuted, <->] (10.2,25) -- (10.8,25) node[midway, above, popmuted] {Gouttière (espace)};
    
    % Dimensions
    \draw[popdark, <->] (-0.5,0) -- (-0.5,29.7) node[midway, left] {297 mm};
    \draw[popdark, <->] (0,-0.5) -- (21,-0.5) node[midway, below] {210 mm};
\end{tikzpicture}
\legende{Structure type d'une page A4 Portrait en 2 colonnes : marges, en-tête, pied de page, gouttière.}
\end{popfigure}

\courssub{Marges et zones de sécurité}

Les \cle{marges} sont les espaces blancs non imprimés entre le bord du papier et le contenu. Elles ne sont pas du « vide perdu » : elles assurent la lisibilité, laissent de la place pour la reliure (gouttière) et absorbent les défauts de coupe du massicot (la machine qui coupe les ramettes).

\begin{definition}[Marges et zone utile]
On distingue :
\begin{itemize}
    \item \textbf{Marges haute et basse :} Espace pour en-tête/pied de page et respiration visuelle. Minimum conseillé : 15 à 20 mm.
    \item \textbf{Marge gauche (gouttière) :} Souvent plus large (20–25 mm) si le document est relié (agrafé, spirale, dos carré collé) pour ne pas « manger » le texte dans la pliure.
    \item \textbf{Marge droite :} Symétrique ou légèrement plus étroite (15 mm).
    \item \textbf{Zone utile (zone de sécurité) :} Rectangle intérieur aux marges où tout contenu \textbf{garanti} visible après coupe. Les imprimeurs exigent souvent 5 mm de marge technique supplémentaire à l'intérieur des marges visibles.
\end{itemize}
\end{definition}

\begin{savaistu}[Le « manger » du massicot]
Dans une imprimerie professionnelle (ou même au secrétariat du lycée avec le grand massicot), la lame coupe des centaines de feuilles à la fois. Le papier glisse légèrement : on perd facilement \textbf{2 à 5 mm} sur chaque bord. C'est pourquoi on ne place \textbf{jamais} un cadre texte ou une image \emph{exactement} au bord de la marge visible. On laisse une « zone de sécurité » de 5 mm à l'intérieur des marges. Pour un fond perdu (image qui va jusqu'au bord du papier), on étend l'image de 3 à 5 mm \emph{au-delà} du format final (fond perdu), puis on coupe.
\end{savaistu}

\courssub{Structure en colonnes}

Diviser la zone utile en \cle{colonnes} est la technique de base pour structurer un journal, un bulletin, une newsletter. Cela réduit la longueur des lignes (confort de lecture : 45 à 75 caractères par ligne), dynamise la page et permet de mixer texte et images.

\begin{propriete}[Règles de mise en colonnes]
\begin{itemize}
    \item \textbf{Nombre :} 1 colonne (rapport, lettre), 2 colonnes (bulletin, article), 3 colonnes (journal, magazine dense). Au-delà, c'est illisible en A4.
    \item \textbf{Largeur :} Égales ou inégales (ex: 2/3 - 1/3 pour une barre latérale).
    \item \textbf{Gouttière (espace inter-colonnes) :} Minimum 4 à 6 mm pour que l'œil ne saute pas d'une colonne à l'autre involontairement.
    \item \textbf{Alignement :} Les bases de lignes (baselines) des colonnes voisines doivent être \emph{alignées horizontalement} (grille de base) pour un aspect professionnel.
\end{itemize}
\end{propriete}

\courssub{En-tête et pied de page : repères et navigation}

L'\cle{en-tête} (header) et le \cle{pied de page} (footer) sont des zones répétées sur chaque page (maîtres-pages). Ils donnent l'identité du document et aident à la navigation.

\begin{exemplebox}[Bulletin de l'Association des Anciens Élèves - Lycée de Ngoa-Ekellé]
\textbf{En-tête (identique page 1 à 4) :}
\begin{itemize}
    \item \textbf{Gauche :} Logo de l'association (15 mm haut).
    \item \textbf{Centre :} « Bulletin de Liaison N° 12 — Juin 2025 » (Police Serif, 14 pt, Gras).
    \item \textbf{Droite :} « Association des Anciens du Lycée de Ngoa-Ekellé » (Police Sans Serif, 10 pt).
\end{itemize}
\textbf{Pied de page (identique page 1 à 4) :}
\begin{itemize}
    \item \textbf{Gauche :} « Siège : BP 123 Yaoundé — Tel: 2XX XX XX XX ».
    \item \textbf{Centre :} Numérotation : « Page 1 / 4 » (champ auto \texttt{PageNumber} / \texttt{TotalPages}).
    \item \textbf{Droite :} « Email : alumni@ngoa-ekelle.cm ».
\end{itemize}
\textit{Astuce :} Sur la page de garde (couverture), on masque souvent l'en-tête/pied de page via un « Maître-page Première page » différent.
\end{exemplebox}

\courssub{Blocs (cadres) de texte et d'images : le principe fondamental de la PAO}

C'est ici que la PAO diffère radicalement du traitement de texte (Word, Writer). Dans Word, le texte « coule » de haut en bas, de page en page. En PAO, \textbf{tout est objet} : chaque titre, paragraphe, photo, logo, tableau, ligne décorative est enfermé dans un \cle{cadre} (ou \cle{bloc}, \cle{frame}) indépendant, positionné librement par coordonnées $(x, y)$ sur la page.

\begin{definition}[Cadre (Frame / Text Frame / Image Frame)]
Un cadre est un conteneur rectangulaire (ou de forme libre) qui possède :
\begin{itemize}
    \item Une \textbf{position} (coin haut-gauche $x, y$ par rapport aux marges ou au bord page).
    \item Des \textbf{dimensions} (largeur $L$, hauteur $H$).
    \item Un \textbf{contenu} (texte lié/chaîné, image, couleur de fond, rien).
    \item Des \textbf{propriétés} : bordure (traits), fond, arrondis, ombre, texte contour (text wrap).
\end{itemize}
\textbf{Chaînage de cadres :} Pour un article long sur plusieurs colonnes/pages, on crée plusieurs cadres texte et on les \emph{chaîne} (lien « suivant »). Le texte excédentaire coule automatiquement du cadre 1 vers le cadre 2.
\end{definition}

\begin{popfigure}
\begin{tikzpicture}[scale=0.40, every node/.style={font=\small}]
    % Page A3 Paysage (42 x 29.7 cm) -> scale 0.4 => ~16.8cm width
    \draw[popdark, line width=2pt] (0,0) rectangle (42,29.7);
    \draw[popblue!50, dashed] (2,2) rectangle (40,27.7);
    
    % Zone Titre (Haut, pleine largeur)
    \draw[poppink, line width=2pt, rounded corners=5pt] (3,22) rectangle (39,27);
    \node[poppink, align=center, font=\bfseries] at (21,24.5) {ZONE TITRE PRINCIPAL\normalsize (Police de titre, 48-72 pt, Centré)};
    
    % Zone Image Centrale (Dominante)
    \draw[popblue, line width=2pt, rounded corners=5pt] (3,8) rectangle (25,21);
    \node[popblue, align=center] at (14,14.5) {ZONE IMAGE PRINCIPALE\normalsize (Photo HD, fond perdu possible)};
    
    % Zone Infos (Droite de l'image)
    \draw[poporange, line width=2pt, rounded corners=5pt] (27,15) rectangle (39,21);
    \node[poporange, align=center] at (33,18) {ZONE INFOS CLÉS\normalsize Date, Lieu, Heure, Tarif};
    
    % Zone Description (Sous image)
    \draw[popgreen, line width=2pt, rounded corners=5pt] (3,8) rectangle (25,7);
    \node[popgreen, align=center] at (14,7.5) {ZONE DESCRIPTION / DÉTAILS};
    
    % Zone Contact / Bas
    \draw[poppurple, line width=2pt, rounded corners=5pt] (27,3) rectangle (39,7);
    \node[poppurple, align=center] at (33,5) {ZONE CONTACT / RÉSEAUX\normalsize Tel, Email, QR Code, Logo Partenaires};
    
    % Zone Logo bas gauche
    \draw[popgold, line width=2pt, rounded corners=5pt] (3,3) rectangle (12,7);
    \node[popgold, align=center] at (7.5,5) {LOGO\normalsize Structure};
    
    % Flèches flux visuel
    \draw[popdark, ->, thick, decorate, decoration={snake, amplitude=1pt, segment length=5pt}] (21,22) -- (14,21) node[midway, right, popdark] {Œil};
    \draw[popdark, ->, thick] (14,8) -- (14,7.5);
    \draw[popdark, ->, thick] (33,15) -- (33,7);
\end{tikzpicture}
\legende{Zonage type d'une affiche A3 Paysage : chaque zone colorée est un cadre indépendant déplaçable.}
\end{popfigure}

\courssub{Typographie : polices, tailles, hiérarchie visuelle}

La typographie porte le sens. Un mauvais choix de police rend le document amateur, voire illisible.

\begin{propriete}[Règles d'or de la typographie en PAO (niveau 3ème)]
\begin{enumerate}
    \item \textbf{Limite : 2 à 3 polices maximum} par document (ex: 1 pour titres, 1 pour corps, 1 optionnelle pour citations/code).
    \item \textbf{Contraste :} Associer une \textbf{Serif} (avec empattements : Times, Garamond, Libre Baskerville) pour le corps de texte (lisibilité en pavé) et une \textbf{Sans Serif} (sans empattements : Arial, Helvetica, Roboto, Open Sans) pour les titres (modernité, impact). \textit{Éviter :} 2 Serif ensemble ou 2 Sans Serif trop proches.
    \item \textbf{Hiérarchie par la taille et le gras :}
    \begin{itemize}
        \item Titre principal : 24–36 pt (Gras).
        \item Sous-titre (intertitre) : 16–18 pt (Gras ou Semi-gras).
        \item Corps de texte : 10–12 pt (Régulier). Interligne (leading) : 1,2 à 1,5 × taille police.
        \item Notes de bas de page / Légendes : 8–9 pt (Italique ou Gris).
    \end{itemize}
    \item \textbf{Alignement :} \textbf{Justifié} pour les colonnes de journal/bulletin (bords nets), \textbf{Gauche} (drapeau) pour les lettres, rapports, web (meilleure lecture, pas de « trous » dans les mots). \textit{Jamais centré} pour un pavé de texte > 3 lignes.
\end{enumerate}
\end{propriete}

\begin{attention}[Piège classique : La « fête des polices »]
\begin{itemize}
    \item \textbf{Erreur :} Utiliser \texttt{Comic Sans} pour le titre, \texttt{Papyrus} pour le sous-titre, \texttt{Brush Script} pour une citation, \texttt{Times} pour le corps, \texttt{Impact} pour le pied de page. Résultat : document illisible, visuellement agressif, aspect « kermesse ».
    \item \textbf{Solution :} Choisir \textbf{une} famille de police complète (ex: \texttt{Roboto} : Thin, Light, Regular, Medium, Bold, Black + Italiques). On varie les \textbf{graisses} (poids) et les \textbf{tailles}, pas les familles.
    \item \textbf{Rappel légal :} Vérifier la licence de la police (Google Fonts = libre, polices Windows/Mac = souvent restreintes pour PDF commercial).
    \item \textbf{Veuves et orphelines :} Éviter qu'une ligne isolée (veuve) commence une page/colonne, ou qu'une ligne isolée (orpheline) termine un paragraphe en bas de page/colonne. Activer l'option « Empêcher veuves/orphelines » dans les styles de paragraphe.
\end{itemize}
\end{attention}

\courssub{Couleurs : harmonie, contraste et coût}

La couleur attire l'œil, code l'information (ex: titre rouge = important), mais coûte cher à l'impression.

\begin{definition}[Modèles de couleurs et impression]
\begin{itemize}
    \item \textbf{RVB (RGB) :} Rouge, Vert, Bleu. \textbf{Écran uniquement} (lumière additive). Ne pas envoyer un PDF RVB à l'imprimeur sans conversion.
    \item \textbf{CMJN (CMYK) :} Cyan, Magenta, Jaune, Noir (Key). \textbf{Impression standard} (encres soustractives). Les couleurs sont moins vives qu'à l'écran (gamut réduit).
    \item \textbf{Pantone (Couleurs directes) :} Encre pré-mélangée unique (ex: « Bleu Pantone 286 C »). Coût : 1 passage machine par couleur. Réservé aux logos, chartes graphiques strictes.
\end{itemize}
\end{definition}

\begin{methode}[Gérer la couleur pour un document scolaire]
\begin{enumerate}
    \item \textbf{Définir une palette limitée :} 1 couleur principale (identité), 1 secondaire (accent), noir, blanc, 1 gris neutre. Code hexa / CMJN noté dans un coin du gabarit.
    \item \textbf{Contraste texte/fond :} Ratio minimum 4,5:1 (accessibilité). \textbf{Noir sur blanc/jaune clair} = OK. \textbf{Gris clair sur blanc} = INTERDIT (illisible photocopié). \textbf{Rouge sur bleu} = INTERDIT (vibration optique, daltoniens).
    \item \textbf{Mode « Noir seul » (Économie) :} Si budget serré (photocopieuse école), concevoir en \textbf{niveaux de gris} + 1 couleur spot (ex: bleu pour titres) $\rightarrow$ 2 passages machine (Noir + Bleu) au lieu de 4 (CMJN).
    \item \textbf{Export PDF :} Cocher « Convertir en CMJN » ou « Simuler l'impression » dans Scribus/Publisher avant d'envoyer au cybercafé.
\end{enumerate}
\end{methode}

\courssub{Alignements, interlignes et aération : le « blanc » qui respire}

L'espace blanc (\emph{white space} ou \emph{espace négatif}) n'est pas du vide, c'est un \textbf{élément de design actif}. Il structure, sépare, repose l'œil.

\begin{propriete}[Paramètres d'aération essentiels]
\begin{itemize}
    \item \textbf{Interligne (Leading) :} Espace entre deux lignes de base. Défaut souvent 1,0 (trop serré). \textbf{Règle :} Taille police × 1,2 à 1,5. Ex: Police 11 pt $\rightarrow$ Interligne 13,2 à 16,5 pt.
    \item \textbf{Espace avant/après paragraphe :} Ne \textbf{jamais} faire « Entrée » deux fois (crée un paragraphe vide ingérable). Utiliser les styles : \texttt{Espace avant = 6 pt, Espace après = 0 pt} (ou 3 pt).
    \item \textbf{Retrait (Indentation) :} Première ligne du paragraphe rentrante (3 à 5 mm) \textbf{SAUF} premier paragraphe après un titre.
    \item \textbf{Alignement vertical dans les cadres :} Haut (défaut), Centre (cartes de visite), Bas (rares), Justifié vertical (remplit le cadre — attention aux grands écarts interlignes).
    \item \textbf{Grille de base (Baseline Grid) :} Forcer les lignes de texte de \emph{tous} les cadres à s'aligner sur une grille invisible (pas = interligne corps). Indispensable pour aligner les colonnes côte à côte.
\end{itemize}
\end{propriete}

\begin{exoresolu}[Structurer la page d'un bulletin scolaire A4 2 colonnes]
\textbf{Énoncé :} Vous devez mettre en page la « Gazette du Collège », bulletin de 4 pages A4, 2 colonnes. Contenu : Édito (1/2 page), Articles (3 pages), Agenda (1/4 page bas page 4), Photos (2). Contraintes : Impression noir/bleu (2 couleurs), agrafé 2 points (marge gauche 20 mm).
\\ \tcblower
\textbf{Solution détaillée (Démarche) :}
\begin{enumerate}
    \item \textbf{Format/Orientation :} A4 Portrait (standard, imprimante OK).
    \item \textbf{Marges :} Haute 15, Basse 15, Droite 15, Gauche \textbf{20 mm} (gouttière agrafage). Zone utile : 175 × 267 mm.
    \item \textbf{Colonnes :} 2 colonnes, largeur 83 mm chacune, gouttière 6 mm. Grille de base activée (pas 14 pt pour corps 11 pt).
    \item \textbf{Maîtres-pages (Gabarits) :}
        \begin{itemize}
            \item \texttt{Maître A (Pages intérieures 2-3)} : En-tête (Nom journal + N° page auto), Pied (Mentions légales). 2 cadres textes chainés couvrant les 2 colonnes.
            \item \texttt{Maître B (Page 1 Couverture)} : Pas d'en-tête/pied. Grands cadres libres pour Titre, Édito, Photo une.
            \item \texttt{Maître C (Page 4)} : Comme Maître A + Cadre fixe bas pour Agenda (hauteur 40 mm, verrouillé).
        \end{itemize}
    \item \textbf{Styles de paragraphes (Feuille de styles) :}
    \begin{center}
    \begin{tabularx}{\linewidth}{|l|X|X|}\hline
    \rowcolor{popblueL}
    \textbf{Style} & \textbf{Police / Taille / Couleur} & \textbf{Usage} \\\hline
    \texttt{TitreUne} & Roboto Bold 36 pt / Bleu Pantone & Une seule fois (Page 1) \\\hline
    \texttt{TitreArt} & Roboto Bold 18 pt / Noir & Début articles \\\hline
    \texttt{CorpsTexte} & Libre Baskerville 11 pt / Noir / Interl 14 pt / Justifié & Pavés articles \\\hline
    \texttt{Legende} & Roboto Italic 9 pt / Gris 60\% / Centré & Sous photos \\\hline
    \texttt{Agenda} & Roboto Medium 10 pt / Bleu / Gauche & Bas page 4 \\\hline
    \end{tabularx}
    \end{center}
    \item \textbf{Placement contenu :} Appliquer styles. Chaîner cadres texte. Insérer images (cadres image), appliquer « Texte contour » (carré 3 mm). Vérifier veuves/orphelines (dernière ligne seule en haut/page, première seule en bas).
    \item \textbf{Export :} PDF/X-1a (CMJN), traits de coupe 3 mm, fonds perdus 3 mm si photo bord à bord.
\end{enumerate}
\end{exoresolu}

\begin{experience}[Mini-TP : Créer une fiche de lecture sur une page A5]
\textbf{Objectif :} Maîtriser la création de cadres, l'application de styles et l'alignement sur une petite surface (A5).
\textbf{Logiciel :} Scribus (libre) ou Publisher / Canva.
\textbf{Consignes :}
\begin{enumerate}
    \item Nouveau document : A5 (148×210 mm), Portrait, Marges 12 mm tout autour.
    \item Créer 2 styles : \texttt{TitreFiche} (Montserrat Bold 20 pt, Bleu \#1A3C6E, Espace après 6 pt) et \texttt{CorpsFiche} (Merriweather 10 pt, Noir, Interl 13 pt, Justifié, Retrait 1ère ligne 4 mm).
    \item Page 1 : Cadre titre (haut, largeur page). Cadre image (sous titre, 50\% largeur, aligné gauche). Cadre texte principal (remplissant le reste, chainé si besoin).
    \item Insérer : Titre du livre, Auteur, Genre. Image couverture. Résumé (texte factice \texttt{lorem ipsum} ou vrai résumé).
    \item Pied de page : Cadre fin (hauteur 8 mm) : « Fiche réalisée par [Prénom Nom] — 3ème [Classe] — 2025 ». Police 8 pt, Gris, Centré.
    \item \textbf{Vérification :} Activer « Aperçu d'impression » (séparations CMJN). Vérifier que le bleu est bien une couleur spot (ou CMJN défini). Exporter PDF « Haute qualité ».
\end{enumerate}
\textbf{Observations attendues :} L'image ne déborde pas dans la marge. Le texte ne touche pas le cadre image (marge interne cadre 3 mm). Les lignes de texte de la colonne de droite s'alignent sur celles de gauche (grille de base). Pas de ligne isolée en haut/bas de cadre.
\end{experience}

\begin{aretenir}[Synthèse : Les 7 piliers d'une mise en page réussie]
\begin{enumerate}
    \item \textbf{Format \& Orientation} adaptés au support final et à l'imprimante.
    \item \textbf{Marges} généreuses (sécurité massicot + reliure).
    \item \textbf{Colonnes} régulières, gouttière visible, grille de base alignée.
    \item \textbf{En-tête/Pied de page} informatifs, cohérents sur tous les maîtres-pages.
    \item \textbf{Cadres (Blocs)} : tout est objet positionné, chaîné pour le flux texte.
    \item \textbf{Typographie} : 2 polices max (Serif/Sans), hiérarchie taille/gras claire, interligne 1,2–1,5.
    \item \textbf{Couleur} : Palette restreinte, contraste fort, anticipation CMJN / Coût impression.
\end{enumerate}
\end{aretenir}

\coursec{Démarche de conception d'un document de PAO}

Maintenant que nous connaissons \emph{les éléments d'une mise en page} (marges, colonnes, alignements, styles, gabarits), une question essentielle se pose : \emph{comment passer de l'idée à l'affiche finale sans perdre de temps ni commettre d'erreurs ?} Un document de PAO réussi ne s'improvise pas devant l'écran ; il suit une démarche structurée, itérative et rigoureuse. C'est cette méthode, adaptée à notre niveau de 3ème, que nous allons détailler. Elle vous servira pour le BEPC, mais aussi pour tout projet de communication visuelle : affiche de kermesse, bulletin de l'association des élèves, flyer pour la boutique du quartier, ou même votre futur CV.

\courssub{Les six étapes de la conception : une méthode universelle}

La conception d'un document de PAO suit un cycle en six étapes. Les trois premières se font \emph{avant} d'allumer l'ordinateur : c'est le travail de réflexion et de préparation. Les trois dernières sont la réalisation technique, le contrôle qualité et la diffusion. On représente souvent ce processus par un organigramme qui montre bien la boucle de retour entre la vérification et la réalisation.

\begin{popfigure}
\begin{tikzpicture}[
    node distance=1.2cm and 1.5cm,
    every node/.style={font=\small},
    etape/.style={rectangle, rounded corners, draw=popblue, thick, fill=popblue!20, text width=3.2cm, align=center, minimum height=1.8cm},
    fleche/.style={-Stealth, thick, popdark},
    retour/.style={-Stealth, thick, poporange, dashed}
]
% Nœuds
\node[etape, align=center] (e1){\textbf{1. Analyse du besoin}\\\textit{Document ? Public ? Message ?}};
\node[etape, right=of e1, align=center] (e2){\textbf{2. Collecte des éléments}\\\textit{Textes, photos, logos}};
\node[etape, right=of e2, align=center] (e3){\textbf{3. Maquette (brouillon)}\\\textit{Croquis main levée}};
\node[etape, below=of e2, align=center] (e4){\textbf{4. Réalisation sur ordinateur}\\\textit{Logiciel, blocs, insertion}};
\node[etape, left=of e4, align=center] (e5){\textbf{5. Vérification}\\\textit{Orthographe, lisibilité, cohérence}};
\node[etape, left=of e5, align=center] (e6){\textbf{6. Impression / Export PDF}\\\textit{Essai, correction, final}};
% Flèches principales
\draw[fleche] (e1) -- (e2);
\draw[fleche] (e2) -- (e3);
\draw[fleche] (e3) -- (e4);
\draw[fleche] (e4) -- (e5);
\draw[fleche] (e5) -- (e6);
% Boucle de retour Vérification -> Réalisation
\draw[retour] (e5.west) -- ++(-1,0) |- (e4.west) node[midway, left, align=center, text width=2cm, font=\tiny] {\textbf{Boucle de correction}\\\textit{Si erreur : retour à l'étape 4}};
% Flèche de départ (cycle)
\draw[fleche, popgreen] (e6.south) -- ++(0,-0.5) -| (e1.south) node[midway, below, align=center, text width=2cm, font=\tiny] {\textit{Nouveau projet}};
\end{tikzpicture}
\legende{Organigramme des 6 étapes de conception d'un document PAO. La flèche pointillée orange symbolise la correction itérative : on revient à la réalisation tant que la vérification n'est pas validée.}
\end{popfigure}

\begin{methode}[Les 6 étapes de conception d'un document de PAO — À connaître pour le BEPC]
\textbf{Étape 1 — Analyse du besoin} : Avant tout, on répond aux questions fondamentales :
\begin{itemize}
    \item \textbf{Quel document ?} (Affiche A3, flyer A5 recto-verso, carte de visite, brochure 4 pages…)
    \item \textbf{Pour quel public ?} (Élèves, parents, passants dans la rue, clients de la boutique…)
    \item \textbf{Quel message principal ?} (Une date, un prix, un slogan, une invitation…)
    \item \textbf{Quel budget d'impression ?} (Noir et blanc sur photocopieuse école ? Couleur chez l'imprimeur du quartier ? Nombre d'exemplaires ?)
\end{itemize}
Ces réponses dictent \emph{tous} les choix techniques ultérieurs (format, couleurs, résolution, logiciel).

\textbf{Étape 2 — Collecte des éléments} : On rassemble \emph{tout} le contenu \textbf{avant} de mettre en page.
\begin{itemize}
    \item Textes : rédigés, relus, corrigés, validés (pas de « je mettrai le texte plus tard »).
    \item Images : photos (résolution suffisante), logos (de préférence vectoriels), icônes.
    \item Informations pratiques : dates, heures, lieux, numéros de téléphone, adresses email, comptes réseaux sociaux, QR codes.
\end{itemize}
\textit{Astuce :} Créez un dossier sur le bureau nommé \texttt{Projet\_PAO\_NomProjet} avec des sous-dossiers \texttt{/textes}, \texttt{/images}, \texttt{/logos}.

\textbf{Étape 3 — Maquette (brouillon)} : C'est l'étape la plus négligée et la plus cruciale. On dessine \textbf{à la main}, sur une feuille blanche (format A4 plié en 2 ou 4 pour simuler le format final), la disposition approximative de chaque bloc : titre, image principale, colonnes de texte, logo, pied de page. On y note les dimensions relatives (ex: « titre = 1/3 hauteur page »), les alignements, les couleurs dominantes. Cela prend 10 minutes et économise des heures de « glissé-déplacé » hasardeux à l'écran.

\textbf{Étape 4 — Réalisation sur ordinateur} :
\begin{enumerate}
    \item Choix du logiciel adapté (Word/Writer pour simple, Publisher/Canva pour mise en page avancée).
    \item Création du document : bon format, bonnes marges, gabarit si document multipage.
    \item Création des \emph{blocs} (cadres) textes et images selon la maquette.
    \item Importation des contenus collectés (Étape 2).
    \item Application des styles (polices, tailles, couleurs, interlignes) définis lors de l'analyse.
\end{enumerate}

\textbf{Étape 5 — Vérification (Contrôle qualité)} : On ne valide pas tant que tout n'est pas coché :
\begin{itemize}
    \item \textbf{Orthographe/grammaire} : relecture à voix haute, correcteur activé, \emph{et} relecture par un camarade.
    \item \textbf{Lisibilité} : contraste suffisant (texte foncé sur fond clair ou inverse), police lisible à distance (pour une affiche), tailles hiérarchisées.
    \item \textbf{Qualité images} : pas de pixelisation visible (zoom 100\% à l'écran), pas de déformation (proportions conservées).
    \item \textbf{Cohérence informations} : la date correspond-elle au jour de la semaine ? Le lieu est-il exact ? Le numéro de téléphone a-t-il 9 chiffres (format Cameroun) ? Les liens/QR codes fonctionnent-ils ?
\end{itemize}
\textbf{Si la moindre erreur est trouvée : retour à l'Étape 4 (flèche orange sur l'organigramme).}

\textbf{Étape 6 — Impression ou Export} :
\begin{itemize}
    \item \textbf{Impression d'essai (épreuve)} : une seule copie sur l'imprimante finale (ou papier équivalent). Vérification couleurs, coupes, pliage.
    \item \textbf{Correction finale} si besoin (retour Étape 4).
    \item \textbf{Impression finale} (nombre d'exemplaires) \textbf{OU} \textbf{Export PDF} (pour diffusion numérique : WhatsApp, email, site web). Le PDF verrouille la mise en page.
\end{itemize}
\end{methode}

\courssub{Étape 1 : Analyse du besoin — Le cahier des charges simplifié}

Tout projet PAO commence par une discussion (avec soi-même ou le « client »). Imaginons : \emph{« Je dois faire l'affiche de la fête de fin d'année du collège. »} Sans analyse, on ouvre Word, on met une photo floue, un titre trop petit, et on imprime 50 exemplaires en noir et blanc sur du A4. Résultat : personne ne la lit.

\begin{definition}[Analyse du besoin (Cahier des charges)]
C'est l'étape d'identification des contraintes et objectifs du document. Elle produit une fiche récapitulative qui guide \emph{tous} les choix techniques. Pour un élève de 3ème, elle se résume à 4 questions clés :
\begin{enumerate}
    \item \textbf{Format final et support :} A4, A3, A5, carte de visite ? Portrait ou Paysage ? Papier ordinaire (80 g/m²) ou papier un peu plus épais (135--170 g/m²) ?
    \item \textbf{Public cible et distance de lecture :} Lu de près (flyer en main) ou de loin (affiche sur mur à 3 m) ? Public jeune (couleurs vives, police ludique) ou officiel (sobriété, police serif) ?
    \item \textbf{Message prioritaire (accroche) :} Quelle est l'\textbf{unique} information que le lecteur doit retenir en 3 secondes ? (Ex: \textbf{« FÊTE FIN D'ANNÉE — 15 JUIN — 14H »})
    \item \textbf{Contraintes techniques et budgétaires :} Impression noir/blanc ou couleur ? Recto seul ou recto-verso ? Imprimante école (marges non imprimantes larges) ou imprimeur pro ? Délai ?
\end{enumerate}
\end{definition}

\begin{exemplebox}[Analyse pour l'affiche « Fête de fin d'année — Collège Bilingue de Bafoussam »]
\begin{center}
\begin{tabularx}{\linewidth}{|l|X|}
\hline
\textbf{Critère} & \textbf{Réponse / Choix justifié} \\
\hline
Format / Support & \textbf{A3 Portrait (29,7 × 42 cm)}, papier couleur 135 g/m². \textit{Justification :} assez grand pour être vu à 5 m sur les panneaux du quartier, tient sur les panneaux d'affichage standard, papier un peu épais pour tenir dehors. \\
\hline
Public & Élèves, parents, passants du quartier. Lecture \textbf{loin (3--5 m)} et \textbf{proche (main propre)}. \\
\hline
Message prioritaire & \textbf{« GRANDE FÊTE — SAMEDI 15 JUIN — 14H00 »} (Gros, haut, couleur vive). Secondaire : lieu, programme, contact. \\
\hline
Budget / Impression & 50 exemplaires couleur chez \textit{Photocopie Express} (quartier). \textit{Contrainte :} marges techniques de l'imprimeur (souvent 5 mm). Pas de finition coûteuse (vernis). \\
\hline
Logiciel envisagé & \texttt{Canva} (modèles prêts, export PDF simple) ou \texttt{LibreOffice Writer} / \texttt{Word} (disponibles au collège). \\
\hline
\end{tabularx}
\end{center}
\end{exemplebox}

\courssub{Étape 2 : Collecte des éléments — « Mise en place » avant « Mise en page »}

\begin{attention}[Erreur classique : « Je mettrai le texte plus tard »]
\textbf{Commencer la mise en page avec du « Lorem Ipsum » ou des cases vides « Photo ici »} est la meilleure façon de casser sa mise en page trois fois. La longueur \emph{réelle} du texte définit la hauteur du bloc ; le ratio \emph{réel} de la photo définit l'espace qu'elle occupe. \textbf{Règle d'or :} On ne touche au logiciel de PAO (Étape 4) que quand \textbf{100\% des contenus sont dans le dossier projet, validés et nommés proprement}.
\end{attention}

\begin{methode}[Check-list de collecte (à faire avant l'ordinateur)]
\begin{enumerate}
    \item \textbf{Textes} : Fichier \texttt{.txt} ou \texttt{.odt} unique, sans mise en forme (pas de gras, pas de tabulations), juste les sauts de paragraphe. Relu \textbf{deux fois}.
    \item \textbf{Images} :
    \begin{itemize}
        \item Photos : \textbf{bonne qualité} (pas de photo floue ou trop petite prise sur WhatsApp). Format \texttt{.jpg} ou \texttt{.png}.
        \item Logos / Icônes / Dessins : \textbf{Vectoriels} (\texttt{.svg}, \texttt{.pdf} vectoriel) \emph{idéal} pour le logo de l'école (ne pixelise jamais). Si seulement \texttt{.png} dispo : prendre la plus grande résolution possible.
    \end{itemize}
    \item \textbf{Données variables} : Date (format JJ MOIS AAAA), Heure (24h), Lieu (nom + quartier + ville), Téléphone (format +237 6XX XX XX XX), Email, QR code (généré à l'avance, testé).
    \item \textbf{Chartes graphiques} : Codes couleur (ex: Bleu école \texttt{\#003366}, Orange énergie \texttt{\#FF6600}), polices officielles (si existantes), logo officiel (ne pas le déformer, ne pas le recolorier).
\end{enumerate}
\end{methode}

\courssub{Étape 3 : La maquette (brouillon) — Le plan de l'architecte}

\begin{definition}[Maquette (ou « Rough » / « Wireframe » papier)]
C'est un \textbf{croquis rapide à main levée} (crayon à papier, gomme, règle) sur une feuille représentant le format final (pliée si nécessaire). Elle positionne les \emph{blocs} (zones texte, zones image, filets, fonds de couleur) avec leurs proportions relatives. Elle ne contient \textbf{aucun détail graphique} (pas de polices jolies, pas d'ombrages), seulement des annotations : « Titre — Police grasse 48 pt — Rouge », « Photo groupe — 50\% largeur », « QR code — 3×3 cm — Bas droite ».
\end{definition}

\begin{popfigure}
\begin{tikzpicture}[
    scale=0.8,
    every node/.style={font=\footnotesize},
    bloc/.style={rectangle, draw=popdark, thick, fill=popcream, rounded corners=3pt, inner sep=4pt, align=center},
    trait/.style={draw=popmuted, dashed, thin},
    annotation/.style={popmuted, font=\tiny, align=left}
]
% Cadre A3 Portrait simulé (ratio ~1:1.414) -> largeur 14cm, hauteur ~19.8cm
% Coordonnées en cm (unité par défaut TikZ)
\draw[popline, thick] (0,0) rectangle (14,19.8);
% Marge visuelle
\draw[popmuted, thin] (1,1) rectangle (13,18.8);
% --- BLOCS MAQUETTE (Haut vers Bas) ---
% 1. Bandeau haut couleur
\node[bloc, minimum width=13, minimum height=1.2, anchor=north] (b1) at (7, 19.3) 
    {\textbf{BANDEAU HAUT}\\\textcolor{popblue}{Fond bleu école \#003366}};
% 2. TITRE PRINCIPAL
\node[bloc, minimum width=12, minimum height=2.5, anchor=north] (b2) at (7, 17.6) 
    {\textbf{TITRE PRINCIPAL}\\\texttt{GRANDE FÊTE}\\\textit{Police: Impact / 72 pt / Blanc / Centré}};
% 3. SOUS-TITRE DATE/LIEU
\node[bloc, minimum width=12, minimum height=1.5, anchor=north] (b3) at (7, 14.6) 
    {\textbf{SOUS-TITRE}\\\texttt{Samedi 15 Juin — 14h00}\\\textit{Police: Arial Bold / 36 pt / Orange \#FF6600}};
% 4. IMAGE CENTRALE (Photo élèves)
\node[bloc, minimum width=12, minimum height=5.5, anchor=north] (b4) at (7, 12.6) 
    {\textbf{IMAGE PRINCIPALE}\\\textit{Photo groupe élèves (paysage)}\\\texttt{Bonne résolution — Cadre arrondi}};
% 5. ZONE TEXTE PROGRAMME (2 colonnes)
\node[bloc, minimum width=5.5, minimum height=4, anchor=north west] (b5a) at (1.5, 6.6) 
    {\textbf{COLONNE G}\\\textit{Programme détaillé}\\\texttt{• 14h Accueil}\\\texttt{• 14h30 Spectacles}\\\texttt{• 16h Remise prix}};
\node[bloc, minimum width=5.5, minimum height=4, anchor=north east] (b5b) at (12.5, 6.6) 
    {\textbf{COLONNE D}\\\textit{Infos pratiques}\\\texttt{• Entrée libre}\\\texttt{• Buvette sur place}\\\texttt{• Parking dispo}};
% 6. BANDEAU BAS : LOGOS + CONTACT + QR
\node[bloc, minimum width=13, minimum height=2.2, anchor=south] (b6) at (7, 1) 
    {\textbf{PIED DE PAGE}\\\texttt{Logo École | Logo Mairie | Logo Sponsor}\\\texttt{Contact: 6XX XX XX XX | college@email.cm}\\\texttt{[ QR CODE 3×3 cm ]}};
% Annotations fléchées (exemples)
\draw[trait] (b1.north west) -- ++(-1,0.5) node[annotation, anchor=south east] {Marge haute 1,5 cm};
\draw[trait] (b2.south) -- ++(0,-0.5) -| ++(3,-1) node[annotation, anchor=north west] {Espace 0,5 cm sous bandeau};
\draw[trait] (b4.south) -- ++(0,-0.5) -| ++(-3,-1) node[annotation, anchor=north east] {Gouttière 1 cm entre img et texte};
\draw[trait] (b6.south west) -- ++(-1,-0.5) node[annotation, anchor=north east] {Marge bas 1,5 cm};
% Légende style
\node[annotation, anchor=south west] at (0.2, 0.2) {\textit{Croquis main levée — Crayon à papier — Annotations techniques}};
\end{tikzpicture}
\legende{Maquette (brouillon) schématisée d'une affiche A3 Portrait. Chaque rectangle représente un bloc (cadre texte ou image). Les annotations techniques (marges, gouttières, tailles de police, codes couleur) guideront la réalisation sur ordinateur.}
\end{popfigure}

\courssub{Étape 4 : Réalisation sur ordinateur — De la maquette au fichier source}

C'est l'étape technique où l'on transpose la maquette dans le logiciel choisi. Le choix du logiciel dépend de l'analyse (Étape 1) :

\begin{itemize}
    \item \textbf{Traitement de texte (Word, LibreOffice Writer) :} Suffisant pour documents simples (1 page, peu d'images, texte dominant). \textit{Attention :} placement précis des blocs parfois difficile.
    \item \textbf{Logiciel de mise en page / PAO (Publisher, Canva, Scribus) :} Indispensable pour : multipages, gabarits, placement libre des blocs (calques), export PDF de qualité.
\end{itemize}

\begin{methode}[Procédure pas-à-pas de réalisation (ex: avec Canva / Publisher / Word)]
\begin{enumerate}
    \item \textbf{Nouveau document} : Format exact (A3 297×420 mm), Marges (ex: 15 mm), Unités : millimètres.
    \item \textbf{Gabarit (Maître de page)} : Si document > 1 page, placer les éléments fixes (numéros de page, filets, logo filigrane) sur le gabarit.
    \item \textbf{Création des cadres (blocs) selon la maquette} :
    \begin{itemize}
        \item Outil « Cadre image » / « Cadre texte » : dessiner les zones photos/logos et textes.
        \item \textbf{Aligner et distribuer} : utiliser les guides, la grille, les outils d'alignement (haut, centres, espacement égal). \textit{Ne jamais placer « à l'œil »}.
    \end{itemize}
    \item \textbf{Importation contenus} :
    \begin{itemize}
        \item Texte : Copier-coller depuis le fichier source → cliquer dans le cadre texte → appliquer le \textbf{Style de paragraphe} correspondant (TitrePrincipal, SousTitre, CorpsTexte, PiedPage).
        \item Images : Insérer image dans le cadre → \texttt{Ajuster l'image au cadre} (proportions conservées).
    \end{itemize}
    \item \textbf{Mise en forme fine} : Interlignes, césure (coupe-mots), alinéas, couleurs (mode RVB pour impression scolaire/bureautique), effets d'ombre (avec modération).
    \item \textbf{Vérification visuelle à 100\%} : Zoom 100\%, parcourir toute la page. Vérifier les « veuves » (dernière ligne seule en haut) et « orphelines » (première ligne seule en bas).
\end{enumerate}
\end{methode}

\courssub{Étape 5 : Vérification — Le contrôle qualité impitoyable}

C'est l'étape qui sépare l'amateur du professionnel. On ne vérifie \textbf{pas} en diagonale. On utilise une check-list physique (papier) qu'on coche au stylo.

\begin{aretenir}[Check-list de vérification PAO (à imprimer et utiliser à chaque projet)]
\begin{enumerate}
    \item \textbf{Orthographe / Grammaire / Typographie} : \texttt{Ctrl+A} → Vérificateur orthographique. Relire \textbf{à voix haute} (force à voir chaque mot). Vérifier : majuscules aux accents (É, À), guillemets français « », espaces insécables avant : ; ? ! et après «, chiffres groupés par 3 (1 000), tirets demi-cadratin – pour énumérations.
    \item \textbf{Informations critiques (Zéro tolérance)} : \textbf{Date / Jour cohérents ?} (Samedi 15 Juin 2025 = VRAI Samedi). \textbf{Lieu exact ?} \textbf{Téléphone 9 chiffres ?} \textbf{Email valide ?} \textbf{QR code scanné et testé ?}
    \item \textbf{Images} : Aucune pixelisation au zoom 100\%. Aucune déformation (ratio L/H respecté). Mode couleur : \textbf{RVB} pour écran/imprimante bureau.
    \item \textbf{Mise en page / Lisibilité} : Alignements parfaits (guides magnétiques). Hiérarchie visuelle claire (Titre > Sous-titre > Corps). Contraste suffisant (test : imprimer en N\&B, est-ce encore lisible ?). Pas de « fleuves » (espaces alignés verticalement dans texte justifié).
    \item \textbf{Aspects techniques impression} : Marges respectées (rien ne touche le bord non imprimable) ? Pages en ordre (pour brochure) ? Police \textbf{intégrée} dans le PDF ?
\end{enumerate}
\end{aretenir}

\begin{attention}[La boucle de retour : Vérification $\rightarrow$ Réalisation]
Si vous trouvez \textbf{une seule} erreur à l'étape 5, vous \textbf{devez} revenir à l'étape 4, la corriger, puis \textbf{refaire l'étape 5 complète}. Ne corrigez pas « juste l'erreur » sans revérifier le reste (une correction peut en créer une autre : texte qui reflow, image qui bouge). C'est la flèche orange pointillée sur l'organigramme.
\end{attention}

\courssub{Étape 6 : Impression d'essai, impression finale, export PDF}

\begin{propriete}[Le format PDF : le standard universel de l'échange imprimable]
Le \textbf{PDF (Portable Document Format)} est le format \emph{final} par excellence pour la PAO.
\begin{itemize}
    \item Il \textbf{fige} la mise en page : polices, images, positions, couleurs restent identiques quel que soit l'ordinateur ou l'imprimante qui l'ouvre.
    \item Il supporte les \textbf{repères de coupe} (si activés à l'export) et la \textbf{transparence}.
    \item Pour l'imprimeur : on exporte en \textbf{PDF « Haute qualité »} ou \textbf{PDF/X} (norme professionnelle).
    \item Pour le web / WhatsApp / Email : on exporte en \textbf{PDF standard} (compression images 150 ppp) pour alléger le fichier.
\end{itemize}
\textbf{Règle :} On \textbf{ne donne jamais} le fichier source (\texttt{.docx}, \texttt{.pub}, \texttt{.canva}) à l'imprimeur ou au client, seulement le PDF validé.
\end{propriete}

\begin{methode}[Export et impression finale]
\begin{enumerate}
    \item \textbf{Impression d'essai (Épreuve)} : Imprimer \textbf{une seule page} (ou le recto-verso plié) sur le papier final (ou le plus proche possible). Vérifier : couleurs réelles (l'écran ment), découpe (les bords sont-ils nets ?), pliage (le texte tombe-t-il dans le pli ?), lisibilité réelle.
    \item \textbf{Correction ultime} : Si l'épreuve révèle un défaut (couleur trop foncée, marge trop juste), retour Étape 4 $\rightarrow$ 5 $\rightarrow$ 6.
    \item \textbf{Export PDF final} :
    \begin{itemize}
        \item Imprimeur : \texttt{Fichier > Exporter > Enregistrer sous PDF} $\rightarrow$ Qualité \texttt{Haute qualité / Impression} $\rightarrow$ Cocher \textbf{Repères de coupe} si dispo. $\rightarrow$ \texttt{Exporter}.
        \item Diffusion numérique : Qualité \texttt{Standard} ou \texttt{Taille de fichier minimale} (images 150 ppp). Vérifier que le fichier final fait < 5--10 Mo pour WhatsApp/Email.
    \end{itemize}
    \item \textbf{Archivage} : Sauvegarder le \textbf{fichier source} (projet logiciel) + le \textbf{PDF final} + le \textbf{dossier « Collecte »} (Étape 2) dans un dossier nommé \texttt{PAO\_NomProjet\_AAAAMMJJ\_FINAL}. C'est votre « dossier de fabrication » pour réimprimer ou modifier l'an prochain.
\end{enumerate}
\end{methode}

\courssub{Justification de la démarche : l'argumentaire écrit (Compétence BEPC)}

Au BEPC, on vous demandera souvent : \emph{« Justifiez les choix techniques effectués pour la réalisation de ce document. »} Ce n'est pas une question d'opinion (« c'est joli »), mais une question de \textbf{lien logique} entre l'analyse du besoin (Étape 1) et les choix techniques (Étapes 3--4--6). La structure type d'une justification est :

\begin{center}
\textbf{Choix technique} $\leftarrow$ \textbf{Parce que} \textbf{Contrainte / Objectif identifié à l'Étape 1}
\end{center}

\begin{exemplebox}[Rédaction guidée de la justification d'Armelle (Affiche Fête Fin d'Année)]
\textbf{Consigne :} \textit{« Armelle a réalisé l'affiche A3 de la fête du collège. Elle a utilisé Canva. Rédigez la justification de ses trois choix principaux. »}

\textbf{Choix 1 : Format et logiciel}
\begin{quote}
\textit{« J'ai choisi le \textbf{format A3 portrait (297×420 mm)} \textbf{parce que} l'analyse du besoin (Étape 1) indiquait une impression sur panneaux d'affichage extérieurs standards A3. J'ai utilisé \textbf{Canva} \textbf{parce que} c'est un outil en ligne gratuit qui gère nativement le format A3, propose des modèles d'affiche, et permet un export PDF haute qualité prêt pour l'imprimeur, sans installer de logiciel. »}
\end{quote}

\textbf{Choix 2 : Image vectorielle pour le logo}
\begin{quote}
\textit{« J'ai inséré le logo de l'école en \textbf{format vectoriel (.svg)} \textbf{parce que} le logo doit apparaître net aussi bien sur l'affiche A3 (grand format) que sur les flyers A5 (petit format) dérivés du même projet. Une image matricielle (JPG/PNG) aurait pixelisé lors de l'agrandissement A3. Le vectoriel garantit une netteté parfaite à toute taille. »}
\end{quote}

\textbf{Choix 3 : Couleurs et police du titre}
\begin{quote}
\textit{« J'ai choisi la \textbf{police Impact en 72 pt, blanche, sur fond bleu école (\#003366)} pour le titre principal \textbf{parce que} la lecture se fait à distance (3--5 m) par des passants (public cible). Impact est une police grasse, sans empattement, à chasse large, très lisible de loin. Le contraste blanc sur bleu foncé assure une lisibilité maximale. La couleur bleue reprend la charte graphique de l'école pour l'identité visuelle. »}
\end{quote}
\end{exemplebox}

\begin{savaistu}[Le PDF : une invention pour « imprimer partout pareil »]
En 1991, John Warnock (cofondateur d'Adobe) lance le projet « Camelot » : créer un format de fichier qui s'affiche et s'imprime \emph{identiquement} sur Mac, Windows, Unix, n'importe quelle imprimante, sans avoir les polices installées. Le PDF 1.0 sort en 1993. Aujourd'hui, c'est le standard mondial : factures, bulletins de notes, contrats, livres, affiches. Au Cameroun, l'administration (Ministères, Mairies, Banques) exige le PDF pour les dépôts en ligne. Un fichier Word \texttt{.docx} peut s'afficher différemment selon la version de Word ou les polices manquantes ; le PDF, lui, est \textbf{fidèle}. C'est pourquoi l'Étape 6 se termine presque toujours par « Export PDF ».
\end{savaistu}

\begin{exoresolu}[Exercice type BEPC : Démarche et Justification]
\textbf{Énoncé :} \textit{M. Kamga, gérant de la boutique « \textit{Informatique Express »} à Yaoundé (quartier Mvog-Ada), veut créer un flyer A5 recto-verso (148×210 mm) pour promouvoir la réparation de téléphones. Il a 5 000 FCFA de budget pour l'impression (noir et blanc sur papier 80 g, 500 exemplaires chez un photocopieur de quartier). Il a un logo vectoriel, une photo de téléphone (JPG 800×600 px), et le texte suivant : « Réparation Express — Écran, Batterie, Connectique — 30 min — 2 000 F à 15 000 F — 6 99 12 34 56 — Carrefour Mvog-Ada — Ouvert 8h-20h ».}

\begin{enumerate}
    \item \textbf{Analyse du besoin} : Complétez le tableau ci-dessous pour M. Kamga.
    \item \textbf{Maquette} : Dessinez (schéma) la maquette du \textbf{recto} en indiquant les blocs, tailles de police approximatives et alignements.
    \item \textbf{Justification} : M. Kamga hésite entre Word et Publisher. Justifiez le choix de Publisher en 2 arguments techniques liés au besoin.
    \item \textbf{Vérification} : La photo fait 800×600 px. Elle occupera 7 cm de large sur le flyer A5. Est-elle de qualité suffisante pour une impression noir/blanc correcte ? (Rappel : 1 pouce = 2.54cm ; seuil minimal 150 ppp).
\end{enumerate}

\tcblower
\textbf{Solution détaillée}

\textbf{1. Analyse du besoin (Tableau)}

\begin{center}
\begin{tabularx}{\linewidth}{|l|X|}
\hline
\textbf{Critère} & \textbf{Réponse justifiée} \\
\hline
Document / Format & \textbf{Flyer A5 Recto-Verso (148×210 mm)}. Format standard, économique, tient dans la main / poche. \\
\hline
Support / Impression & \textbf{Papier 80 g/m² (standard), Noir \& Blanc seul}. Budget 5 000 F pour 500 ex. = 10 F/ex. Imprimeur quartier : photocopieuse N\&B. \textit{Marges techniques min 5--10 mm.} \\
\hline
Public / Lecture & \textbf{Passants, clients quartier, lecture de près (main propre)}. Info rapide. \\
\hline
Message prioritaire & \textbf{« RÉPARATION EXPRESS — 30 MIN »} (Service rapide = argument n°1). \\
\hline
Contraintes techniques & Mode couleur : \textbf{Niveaux de gris / Noir seul}. Résolution images : correcte. Logiciel : doit gérer A5, recto-verso, gabarit, export PDF. \\
\hline
\end{tabularx}
\end{center}

\textbf{2. Maquette du Recto (Description textuelle / Schéma)}

\begin{center}
\begin{tikzpicture}[scale=0.7, every node/.style={font=\small}, bloc/.style={draw=popdark, fill=popcream, rounded corners=2pt, inner sep=3pt, align=center}]
% A5 = 148x210 mm. Dessin ~7x10 cm
\draw[thick] (0,0) rectangle (7,10);
\node[bloc, minimum width=6.5, minimum height=1.2, anchor=north, align=center] at (3.5, 9.8){\textbf{BANDEAU HAUT (Gris 20\%)}\\\texttt{INFORMATIQUE EXPRESS}\\\textit{Logo (Vecteur) | Police: Arial Black 18pt}};
\node[bloc, minimum width=6.5, minimum height=1.8, anchor=north, align=center] at (3.5, 8.2){\textbf{TITRE ACCROCHE}\\\texttt{RÉPARATION EXPRESS}\\\textit{Police: Impact 28pt | Gras | Centré | Noir 100\%}};
\node[bloc, minimum width=6.5, minimum height=1.2, anchor=north, align=center] at (3.5, 6.1){\textbf{SOUS-TITRE SERVICES}\\\texttt{Écran • Batterie • Connectique}\\\textit{Police: Arial 14pt | Centré | Gris 50\%}};
\node[bloc, minimum width=6, minimum height=3.5, anchor=north, align=center] at (3.5, 4.5){\textbf{IMAGE CENTRALE}\\\textit{Téléphone (Cadre arrondi)}\\\texttt{Photo 7cm larg. | Centrée}};
\node[bloc, minimum width=6.5, minimum height=1.5, anchor=south, align=center] at (3.5, 0.3){\textbf{PIED : CONTACT \& LOCALISATION}\\\texttt{30 min — 2 000 à 15 000 F}\\\texttt{6 99 12 34 56 — Carrefour Mvog-Ada}\\\texttt{8h-20h}\\\textit{Police: Arial Bold 11pt | Centré}};
\end{tikzpicture}
\end{center}
\textit{Annotations : Marges 5 mm. Gouttières 3 mm. Alignement centre pour tout (symétrie flyer). Logo vectoriel en haut à gauche dans bandeau.}

\textbf{3. Justification du choix Publisher (vs Word)}

\begin{quote}
\textit{« Je choisis \textbf{Publisher} plutôt que Word pour deux raisons techniques majeures liées au besoin :}
\begin{enumerate}
    \item \textbf{Gestion native des blocs (cadres) et du recto-verso :} Le flyer est recto-verso. Publisher permet de déplacer, redimensionner et superposer des blocs textes/images librement (comme en PAO) sans que le texte ne « saute » de page comme dans Word. Les gabarits « Recto » et « Verso » distincts facilitent l'imposition pour le pliage.
    \item \textbf{Export PDF optimisé pour photocopieuse N\&B :} Publisher exporte en PDF avec polices intégrées et conversion automatique en niveaux de gris pur pour l'impression noir/blanc, évitant les trames de couleur parasites que Word peut générer. Cela garantit un texte net et des images en nuances de gris propres sur les 500 exemplaires. »
\end{enumerate}
\end{quote}

\textbf{4. Calcul de la résolution effective de la photo}

\begin{itemize}
    \item Largeur image source : $L_{px} = 800$ pixels.
    \item Largeur impression souhaitée : $L_{cm} = 7$ cm.
    \item Conversion cm $\rightarrow$ pouces : $1 \text{ pouce} = 2,54 \text{ cm}$.
    \item $L_{pouces} = \frac{7}{2,54} \approx 2,76 \text{ pouces}$.
    \item Résolution effective (PPP/DPI) : $\text{PPP} = \frac{L_{px}}{L_{pouces}} = \frac{800}{2,76} \approx \mathbf{290 \text{ ppp}}$.
\end{itemize}

\textbf{Conclusion :} $290 \text{ ppp} > 150 \text{ ppp}$ (seuil minimal N\&B). \textbf{La photo est suffisante} pour une impression noir/blanc correcte sur photocopieuse.
\end{exoresolu}

\begin{exercice}[Entraînement BEPC : Maquette et Justification]
\textbf{Contexte :} L'association « Jeunes Leaders du Lycée » organise une conférence « Orientation Post-Bac » le \textbf{Samedi 12 Octobre 2024 à 9h} dans la salle des fêtes de la mairie d'arrondissement. Public : élèves de Terminale et parents. Budget : impression 100 affiches A3 couleur sur papier 135 g chez un imprimeur pro (marges techniques 5 mm demandées). Éléments : Logo asso (PNG 500×500 px seulement), photo salle (JPG 3000×2000 px), texte complet fourni.

\begin{enumerate}
    \item Remplissez l'analyse du besoin (Format, Support, Public, Message, Contraintes imprimeur).
    \item Proposez une maquette schématique du recto (blocs + annotations techniques).
    \item Le logo n'existe qu'en PNG 500×500 px. Il doit faire 4 cm de large sur l'affiche A3. Calculez sa résolution effective. Est-ce acceptable ? Quelle solution technique proposez-vous pour améliorer la netteté ?
    \item Rédigez la justification du choix : « Export final en PDF Haute Qualité avec repères de coupe. »
\end{enumerate}
\end{exercice}

\begin{corrige}[Exercice : Maquette et Justification]
\textbf{1. Analyse du besoin}
\begin{itemize}
    \item Format : \textbf{A3 Portrait (297×420 mm)}.
    \item Support : \textbf{Papier couché 135 g/m², Couleur}. Marges techniques \textbf{5 mm} (exigence imprimeur).
    \item Public : \textbf{Élèves Terminale + Parents}. Lecture : \textbf{Loin (panneau mairie 3--5 m) et Proche}.
    \item Message prioritaire : \textbf{« CONFÉRENCE ORIENTATION POST-BAC — 12 OCTOBRE — 9H »}.
    \item Contraintes : PDF Haute Qualité, marges 5 mm, repères de coupe si possible, polices intégrées, images 300 ppp RVB/CMJN. Délai : prêt 3 jours avant.
\end{itemize}

\textbf{2. Maquette Recto (Description)}
Bandeau haut (H=3 cm, Fond bleu marine) → Logo asso (G) + Titre Asso (D). Espace 1 cm.
TITRE PRINCIPAL (H=4 cm, Police grasse 60-70 pt, Blanc/Jaune vif, Centré) : « ORIENTATION POST-BAC ».
SOUS-TITRE (H=1.5cm, Police 30 pt, Blanc) : « Samedi 12 Octobre — 9h00 ».
IMAGE SALLE (H=10 cm, Pleine largeur moins marges, Cadre fin blanc) : Photo ambiance.
BANDEAU BAS (H=3 cm, Fond bleu marine) : Lieu (Mairie), Contact (Tel, FB, QR Code), Logos partenaires. Marge bas 5 mm.

\textbf{3. Logo PNG 500×500 px à 4 cm}
$L_{pouces} = 4 / 2,54 \approx 1,57 \text{ pouces}$.
$\text{PPP} = 500 / 1,57 \approx \mathbf{318 \text{ ppp}}$.
\textbf{Acceptable} ( > 300 ppp). \textit{Mais} : PNG = fond transparent ? Si fond blanc, gêne sur bandeau bleu. \textbf{Solution :} Ouvrir dans GIMP/Photopea → Détourer (sélection couleur / chemin) → Exporter en \textbf{PNG transparent 300 ppp} (redimensionné à 1200 px large pour marge) OU revectoriser (Inkscape « Vectoriser bitmap ») pour obtenir un \textbf{.svg} parfait (recommandé pour logo).

\textbf{4. Justification PDF Haute Qualité + Repères de coupe}
\textit{« L'export en \textbf{PDF Haute Qualité} est exigé par l'imprimeur professionnel car il garantit que le fichier contient \textbf{toutes les polices intégrées} et les images à leur résolution maximale (300 ppp), évitant le flou ou les substitutions de police. Les \textbf{repères de coupe} (traits de coupe) indiquent exactement où couper le papier (format fini 297×420 mm) par rapport aux marges techniques de 5 mm demandées. Sans eux, l'imprimeur ne sait pas où se trouve le bord utile de l'affiche, risquant de couper dans le texte ou de laisser un filet blanc. Cela assure une production industrielle sans défaut. »}
\end{corrige}

\coursec{Applications concrètes et rédaction de la démarche}

Dans la section précédente, nous avons appris la \cle{démarche de conception} d'un document de PAO : analyser le besoin, choisir le format et les éléments de mise en page, réaliser le document, puis le vérifier avant de le diffuser. Il est temps de mettre cette méthode en pratique ! Dans cette dernière section, nous allons étudier trois cas concrets de la vie courante — une affiche de tournoi de football, un bulletin d'information et une carte d'invitation — puis apprendre à \cle{rédiger et justifier} notre démarche, comme on te le demandera au BEPC. Enfin, nous verrons comment \cle{évaluer} la qualité d'un document de PAO.

\courssub{Étude de cas 1 : l'affiche d'annonce d'un tournoi de football}

\textbf{Situation.} Le club sportif de ton quartier à Yaoundé organise un tournoi de football pendant les vacances. Le président du club te demande de concevoir une \cle{affiche} qui sera imprimée en 50 exemplaires et collée dans les boutiques, les églises et les écoles du quartier.

\textbf{Analyse du besoin.} Une affiche a un rôle bien précis : être \emph{vue de loin} et \emph{comprise en quelques secondes}. Elle doit donc être grande, colorée, avec peu de texte mais des informations complètes.

\textbf{Les choix effectués et leur justification :}

\begin{itemize}
  \item \textbf{Format :} A4 (ou A3 si le budget le permet), en orientation \cle{portrait}. Justification : le A4 est le format standard des imprimantes et photocopieuses des imprimeurs ; il est assez grand pour être lu à deux mètres.
  \item \textbf{Couleurs :} les couleurs du club (par exemple vert et jaune), avec un titre en gros caractères gras. Justification : les couleurs vives attirent le regard ; le titre doit être la première chose vue.
  \item \textbf{Image :} une photo d'un ballon de football ou de l'équipe, de bonne qualité. Justification : une image parle immédiatement, même à quelqu'un qui lit vite.
  \item \textbf{Informations obligatoires :} on les retient avec la liste magique « \cle{quoi, qui, quand, où, combien, contact} » :
  \begin{itemize}
    \item \emph{Quoi ?} Tournoi de football inter-quartiers ;
    \item \emph{Qui ?} Organisé par le Club Sportif de Nkolbisson, ouvert aux 13--17 ans ;
    \item \emph{Quand ?} Du 10 au 20 août, matchs à 15 h ;
    \item \emph{Où ?} Terrain du lycée de Nkolbisson ;
    \item \emph{Combien ?} Inscription : 1\,000 FCFA par équipe ;
    \item \emph{Contact ?} Téléphone du président du club.
  \end{itemize}
\end{itemize}

\begin{attention}[L'erreur qui gâche tout]
L'erreur la plus fréquente sur une affiche est d'\textbf{oublier une information essentielle} : la date ou le lieu ! Imagine une magnifique affiche de tournoi... sans la date : personne ne viendra. Avant d'imprimer, vérifie \textbf{toujours} la liste complète : \cle{quoi, qui, quand, où, combien, contact}. Fais relire ton affiche par une autre personne : un œil neuf repère vite les oublis.
\end{attention}

\begin{exemplebox}[Un choix de couleurs a un prix !]
Le président du club hésite entre l'impression couleur et l'impression noir et blanc. Faisons le calcul :
\begin{itemize}
  \item 50 affiches A4 en \textbf{couleur} à 200 FCFA l'unité : $50 \times 200 = 10\,000$ FCFA ;
  \item 50 affiches A4 en \textbf{noir et blanc} à 50 FCFA l'unité : $50 \times 50 = 2\,500$ FCFA.
\end{itemize}
La couleur coûte donc $10\,000 - 2\,500 = 7\,500$ FCFA de plus, soit 4 fois plus cher. \textbf{Conclusion :} si le budget du club est de 5\,000 FCFA, on peut imprimer 20 affiches couleur pour les endroits très fréquentés et compléter par des affiches noir et blanc. Le choix des couleurs n'est pas seulement esthétique : il a un \cle{impact budgétaire} qu'il faut justifier.
\end{exemplebox}

\courssub{Étude de cas 2 : le bulletin d'information de deux pages}

\textbf{Situation.} Le conseil d'administration de ton collège te confie la création du \cle{bulletin d'information} de fin d'année : deux pages A4 distribuées aux parents lors de la remise des bulletins de notes.

\textbf{Analyse du besoin.} Contrairement à l'affiche, un bulletin est un document que l'on \emph{lit attentivement}, posé sur une table. Il contient plusieurs articles (le mot du principal, les résultats au BEPC, les activités de l'année, les dates de la rentrée). Le problème à résoudre est donc différent : \emph{organiser beaucoup de texte pour que la lecture reste agréable}.

\textbf{Les choix effectués et leur justification :}

\begin{itemize}
  \item \textbf{Format :} A4 portrait, deux pages. Justification : format standard, facile à photocopier et à agraffer.
  \item \textbf{Colonnes :} le texte des articles est disposé en \cle{deux colonnes}. Justification : des lignes trop longues fatiguent l'œil ; les colonnes, comme dans les journaux, rendent la lecture plus rapide et plus agréable.
  \item \textbf{En-tête :} sur chaque page, le nom du collège, le logo et le titre du bulletin (« Info-Collège -- n° 12 -- Juin »). Justification : le lecteur identifie immédiatement qui publie le document.
  \item \textbf{Pied de page :} le numéro de page (« page 1/2 ») et un contact. Justification : si les pages se détachent, on peut les remettre dans l'ordre.
  \item \textbf{Photos :} deux ou trois photos (la fête de la jeunesse, les lauréats) avec une \cle{légende} sous chacune. Justification : les photos illustrent les articles et aèrent la page ; la légende explique ce que l'on voit.
  \item \textbf{Titres des articles :} en gras, plus gros que le texte, avec la même police partout. Justification : le lecteur repère vite les sujets qui l'intéressent.
\end{itemize}

\begin{attention}[Trop de polices tue la lisibilité]
Ne mélange jamais plus de \textbf{deux ou trois polices} dans un même document : une pour les titres, une pour le texte. Un bulletin avec six polices différentes donne une impression de désordre et fatigue le lecteur. De même, évite les photos floues ou étirées : une image déformée (écrasée ou allongée) est le signe d'un travail non soigné.
\end{attention}

\courssub{Étude de cas 3 : la carte d'invitation à une cérémonie}

\textbf{Situation.} Ta famille organise la remise d'un diplôme de ta grande sœur et te demande de concevoir les \cle{cartes d'invitation} à envoyer aux invités.

\textbf{Analyse du besoin.} Une carte d'invitation est un petit document \emph{personnel} et \emph{élégant}. Contrairement à l'affiche qui crie, la carte chuchote : elle doit être sobre, raffinée, et contenir exactement les informations utiles, rien de plus.

\textbf{Les choix effectués et leur justification :}

\begin{itemize}
  \item \textbf{Format :} A6 (10,5 cm $\times$ 14,8 cm) ou A5 (14,8 cm $\times$ 21 cm). Justification : le format A6 tient dans une poche ou un petit sac ; on obtient d'ailleurs 4 cartes A6 en découpant une feuille A4, ce qui économise le papier.
  \item \textbf{Sobriété :} fond blanc ou crème, une seule couleur d'accent (doré ou bleu foncé), beaucoup d'espace vide. Justification : l'élégance vient de la simplicité ; une carte surchargée paraît désordonnée.
  \item \textbf{Police élégante :} une police fine et soignée pour les noms, une police simple et lisible pour les informations pratiques. Justification : la police participe à l'ambiance de la cérémonie.
  \item \textbf{Informations :} qui invite (la famille), qui est invité (le nom de l'invité), l'événement (remise de diplôme), la date et l'heure, le lieu exact, et éventuellement « RSVP » avec un numéro de téléphone pour confirmer sa présence.
\end{itemize}

\begin{popfigure}
\begin{center}
\begin{tikzpicture}[scale=0.9, every node/.style={transform shape}]
  % Affiche
  \draw[fill=popgreenL, draw=popdark, thick] (0,0) rectangle (2.6,3.6);
  \node[font=\bfseries\small, text=popdark] at (1.3,3.2) {TOURNOI};
  \draw[fill=white, draw=popline] (0.4,1.9) rectangle (2.2,2.9);
  \node[font=\scriptsize, text=popmuted] at (1.3,2.4) {photo ballon};
  \node[font=\tiny, text=popdark, align=center] at (1.3,1.3) {10--20 août\\Terrain du lycée\\Inscription : 1\,000 F};
  \node[font=\scriptsize\itshape, text=popmuted, align=center] at (1.3,-0.5) {Affiche A4\\colorée, informative};
  % Bulletin
  \draw[fill=popblueL, draw=popdark, thick] (4.2,0) rectangle (6.8,3.6);
  \draw[fill=white, draw=popline] (4.3,3.1) rectangle (6.7,3.5);
  \node[font=\tiny, text=popdark] at (5.5,3.3) {en-tête : logo + titre};
  \draw[line width=1.2pt, popline] (4.5,2.9) -- (5.4,2.9);
  \draw[line width=1.2pt, popline] (4.5,2.6) -- (5.4,2.6);
  \draw[line width=1.2pt, popline] (4.5,2.3) -- (5.4,2.3);
  \draw[line width=1.2pt, popline] (5.7,2.9) -- (6.6,2.9);
  \draw[line width=1.2pt, popline] (5.7,2.6) -- (6.6,2.6);
  \draw[line width=1.2pt, popline] (5.7,2.3) -- (6.6,2.3);
  \draw[fill=white, draw=popline] (4.5,1.1) rectangle (5.4,2.0);
  \node[font=\tiny, text=popmuted] at (4.95,1.55) {photo};
  \draw[line width=1.2pt, popline] (5.7,1.9) -- (6.6,1.9);
  \draw[line width=1.2pt, popline] (5.7,1.6) -- (6.6,1.6);
  \draw[line width=1.2pt, popline] (5.7,1.3) -- (6.6,1.3);
  \node[font=\tiny, text=popmuted] at (5.5,0.4) {page 1/2};
  \node[font=\scriptsize\itshape, text=popmuted, align=center] at (5.5,-0.5) {Bulletin A4\\2 colonnes, en-tête, photos};
  % Carte
  \draw[fill=popgoldL, draw=popdark, thick] (8.6,1.0) rectangle (11.2,3.2);
  \node[font=\scriptsize\itshape, text=popdark] at (9.9,2.7) {Invitation};
  \node[font=\tiny, text=popdark, align=center] at (9.9,1.9) {Remise de diplôme\\Samedi 28 juin, 17 h\\Salle des fêtes};
  \node[font=\scriptsize\itshape, text=popmuted, align=center] at (9.9,-0.5) {Carte A6\\sobre, élégante};
\end{tikzpicture}
\end{center}
\legende{Trois documents de PAO, trois objectifs différents : l'affiche attire de loin, le bulletin informe en détail, la carte invite avec élégance.}
\end{popfigure}

\courssub{Rédiger et justifier une démarche de création}

Au BEPC, on ne te demandera pas seulement de créer un document : on te demandera aussi de \textbf{rédiger ta démarche}, c'est-à-dire d'expliquer par écrit comment tu as procédé et \emph{pourquoi} tu as fait chaque choix. C'est un savoir-faire officiel du programme : « rédiger et justifier une démarche, une réponse ou une conclusion ».

\begin{methode}[Rédiger une démarche de création en 4 parties]
\begin{enumerate}
  \item \textbf{Le besoin :} décris la situation et le problème à résoudre. Qui a besoin du document ? Pour quoi faire ? Qui va le lire ? \emph{Exemple : « Le club sportif veut informer les habitants du quartier d'un tournoi de football. »}
  \item \textbf{Les choix effectués :} liste les décisions prises : format, orientation, couleurs, polices, images, informations incluses. \emph{Exemple : « J'ai choisi le format A4 portrait, les couleurs vert et jaune du club, une photo de ballon. »}
  \item \textbf{La justification de chaque choix :} c'est la partie la plus importante ! Pour chaque choix, explique \emph{pourquoi} il répond au besoin. Utilise des mots comme « parce que », « afin de », « ce qui permet de ». \emph{Exemple : « J'ai choisi le format A4 parce que c'est le format standard des imprimeurs, ce qui réduit le coût d'impression. »}
  \item \textbf{La conclusion :} vérifie que le document répond bien au besoin de départ et dis-le. \emph{Exemple : « L'affiche contient toutes les informations (quoi, qui, quand, où, combien, contact) et respecte le budget du club : elle répond donc au besoin. »}
\end{enumerate}
\end{methode}

\begin{exoresolu}[Rédiger la démarche de la carte d'invitation]
\textbf{Énoncé.} Rédige, en quatre parties, la démarche de création de la carte d'invitation à la cérémonie de remise de diplôme (étude de cas 3).
\tcblower
\textbf{Solution rédigée :}

\emph{1. Besoin.} Ma famille organise une cérémonie de remise de diplôme et doit inviter une cinquantaine de personnes. Il faut un document petit, élégant et complet, que chaque invité pourra garder sur lui.

\emph{2. Choix effectués.} J'ai choisi le format A6, un fond crème avec un liseré doré, une police fine et élégante pour les noms et une police simple pour les informations pratiques. La carte contient : le nom de la famille qui invite, le nom de l'invité, l'événement, la date, l'heure, le lieu et un numéro de téléphone pour confirmer.

\emph{3. Justifications.} Le format A6 a été choisi parce qu'il tient dans une poche et que quatre cartes se découpent dans une seule feuille A4, ce qui réduit le coût. Le fond crème et le liseré doré ont été choisis afin de donner une impression d'élégance, adaptée à une cérémonie. La police simple pour les informations pratiques permet de lire la date et le lieu sans effort. Le numéro de téléphone permet aux invités de confirmer leur présence, ce qui aide la famille à prévoir le nombre de chaises et de repas.

\emph{4. Conclusion.} La carte est sobre, complète et peu coûteuse à produire : elle répond exactement au besoin de la famille.
\end{exoresolu}

\courssub{Évaluer un document de PAO}

Comment savoir si un document est réussi ? On utilise une \cle{grille d'évaluation} avec des critères précis et des points. Cette grille sert aussi à s'auto-évaluer avant d'imprimer.

\begin{aretenir}[Les 5 critères de qualité d'un document de PAO]
\begin{itemize}
  \item \textbf{Lisibilité :} le texte se lit sans effort (taille des caractères, contraste, polices limitées) ;
  \item \textbf{Équilibre :} les éléments sont bien répartis sur la page, ni trop chargée ni trop vide ;
  \item \textbf{Pertinence des informations :} toutes les informations essentielles sont présentes (quoi, qui, quand, où, combien, contact) et il n'y a rien d'inutile ;
  \item \textbf{Qualité des images :} les images sont nettes, non déformées, et en rapport avec le sujet ;
  \item \textbf{Respect des règles de mise en page :} marges régulières, alignements soignés, titres hiérarchisés, en-tête et pied de page si nécessaire.
\end{itemize}
\end{aretenir}

\begin{popfigure}
\begin{center}
\begin{tikzpicture}
  \draw[fill=popblue, draw=popblue] (0,4.0) rectangle (10,4.7);
  \node[text=white, font=\bfseries\small, anchor=west] at (0.2,4.35) {Critère};
  \node[text=white, font=\bfseries\small] at (9.0,4.35) {Points};
  \foreach \y/\txt/\pts in {3.3/{Lisibilité (taille, contraste, polices)}/4, 2.6/{Équilibre de la page}/4, 1.9/{Pertinence des informations}/4, 1.2/{Qualité des images}/4, 0.5/{Respect des règles de mise en page}/4}{
    \draw[draw=popline] (0,\y-0.35) rectangle (10,\y+0.35);
    \node[anchor=west, font=\small] at (0.2,\y) {\txt};
    \node[font=\small\bfseries, text=poporange] at (9.0,\y) {/ \pts};
  }
  \draw[fill=popgoldL, draw=popline] (0,-0.55) rectangle (10,0.15);
  \node[anchor=west, font=\small\bfseries] at (0.2,-0.2) {TOTAL};
  \node[font=\small\bfseries, text=popdark] at (9.0,-0.2) {/ 20};
\end{tikzpicture}
\end{center}
\legende{Exemple de grille d'évaluation d'un document de PAO sur 20 points : chaque critère vaut 4 points.}
\end{popfigure}

\begin{experience}[Évaluer l'affiche d'un camarade]
\textbf{Objectif :} utiliser la grille d'évaluation pour juger un document et proposer des améliorations.

\textbf{Matériel :} l'affiche du tournoi réalisée par un camarade (ou la tienne), la grille d'évaluation ci-dessus.

\textbf{Manipulation :}
\begin{enumerate}
  \item Lis l'affiche et attribue une note de 0 à 4 pour chacun des cinq critères ;
  \item Pour chaque point perdu, écris une phrase expliquant le défaut (exemple : « la date est écrite trop petite, elle ne se lit pas à deux mètres ») ;
  \item Propose une amélioration concrète pour chaque défaut ;
  \item Calcule le total sur 20 et compare avec l'évaluation d'un autre groupe.
\end{enumerate}

\textbf{Observations attendues :} les écarts de notation entre groupes montrent que certains défauts (image floue, information manquante) sont objectifs, tandis que d'autres (goût des couleurs) sont plus subjectifs.

\textbf{Interprétation :} une grille d'évaluation rend le jugement plus juste et plus précis qu'un simple « c'est joli » ou « c'est moche » : elle permet d'\emph{argumenter}, ce qui est exactement la compétence « rédiger et justifier une conclusion ».
\end{experience}

\begin{savaistu}[Les imprimeurs camerounais et le PDF]
À Yaoundé et à Douala, de nombreux imprimeurs acceptent aujourd'hui les fichiers \cle{PDF} envoyés simplement par WhatsApp ou apportés sur une clé USB. Le format PDF est idéal car il conserve exactement la mise en page, les polices et les images, quel que soit l'ordinateur sur lequel on l'ouvre. Avant d'envoyer ton fichier, ouvre-le toi-même en PDF pour vérifier une dernière fois qu'aucune information ne manque !
\end{savaistu}

\begin{savaistu}[Complément (hors programme strict) : le droit d'auteur sur les images]
Les images trouvées sur Internet ne sont pas libres d'utilisation : elles appartiennent à leurs auteurs, protégés par le \cle{droit d'auteur}. Pour tes documents, utilise des \textbf{images libres de droits} (par exemple sur les sites qui les proposent gratuitement) ou tes propres photos, et prends l'habitude de \textbf{citer la source} de chaque image utilisée (nom de l'auteur ou du site). C'est une règle de respect et d'honnêteté numérique que tu retrouveras dans ta vie scolaire et professionnelle.
\end{savaistu}

\begin{exercice}[À toi de jouer]
Ton école organise une journée portes ouvertes.
\begin{enumerate}
  \item Choisis le type de document le plus adapté (affiche, bulletin ou carte) et justifie ton choix en deux phrases ;
  \item Liste les informations obligatoires que ce document doit contenir ;
  \item Rédige ta démarche complète en quatre parties (besoin, choix, justifications, conclusion) ;
  \item Évalue le document d'un camarade avec la grille sur 20 points.
\end{enumerate}
\end{exercice}

\begin{corrige}[Exercice]
\begin{enumerate}
  \item Le document le plus adapté est l'\textbf{affiche}, car il faut informer un large public (élèves, parents, habitants du quartier) rapidement et de loin. On peut la compléter par un bulletin pour les détails du programme.
  \item Informations obligatoires : \emph{quoi} (journée portes ouvertes), \emph{qui} (le collège, ouvert à tous), \emph{quand} (date et horaires), \emph{où} (adresse exacte du collège), \emph{combien} (entrée libre), \emph{contact} (téléphone du secrétariat).
  \item Exemple de démarche : \emph{Besoin} : le collège veut faire connaître sa journée portes ouvertes au public du quartier. \emph{Choix} : affiche A4 portrait, couleurs de l'établissement, photo de la façade, informations essentielles. \emph{Justifications} : le A4 car c'est le format standard des imprimeurs (coût réduit) ; les couleurs de l'établissement pour être reconnu ; la photo de la façade pour que les visiteurs repèrent facilement le bâtiment ; la liste « quoi, qui, quand, où, combien, contact » pour ne rien oublier. \emph{Conclusion} : l'affiche est lisible, complète et peu coûteuse ; elle répond au besoin.
  \item L'évaluation utilise la grille : lisibilité, équilibre, pertinence, qualité des images, respect des règles de mise en page, chacun noté sur 4, pour un total sur 20, chaque point perdu étant expliqué par une phrase.
\end{enumerate}
\end{corrige}

\newpage

\coursec{Activité d'intégration}

\coursec{Initiation à la publication assistée par ordinateur}

\courssub{Objectifs de la leçon}

À la fin de cette leçon, tu seras capable de :

\begin{itemize}
\item définir l'infographie et la publication assistée par ordinateur (\cle{PAO}) ;
\item distinguer une \cle{image matricielle} (bitmap) d'une \cle{image vectorielle} et expliquer l'impact de l'agrandissement sur leur qualité ;
\item citer des logiciels d'infographie (matriciels et vectoriels) et des logiciels de PAO ;
\item énoncer et appliquer les règles fondamentales de mise en page (marges, hiérarchie, polices, couleurs, équilibre) ;
\item suivre une démarche de conception rigoureuse : analyse du besoin $\rightarrow$ maquette $\rightarrow$ réalisation $\rightarrow$ vérification $\rightarrow$ justification ;
\item rédiger une justification argumentée de tes choix techniques et esthétiques.
\end{itemize}

\courssub{Situation déclenchante : Le tournoi de Nkol-Eton}

L'association des jeunes du quartier Nkol-Eton, à Yaoundé, organise son grand tournoi annuel de football. Le président, M. Etoundi, te sollicite, toi et tes camarades de 3\textsuperscript{ème}, pour concevoir l'affiche de l'événement. Voici sa lettre :

\begin{exemplebox}[La demande de M. Etoundi]
« Chers élèves, notre tournoi de football se déroulera du \textbf{samedi 12 juillet au samedi 26 juillet 2025}, tous les samedis et dimanches à partir de \textbf{15 heures}, au \textbf{stade municipal de Nkol-Eton}. Huit équipes de quartiers voisins participent. Les frais d'inscription sont de \textbf{5\,000 FCFA par équipe}, à verser avant le \textbf{5 juillet 2025} auprès du secrétaire de l'association, M\textsuperscript{me} Abena, joignable au \textbf{6\,90\,12\,34\,56}. L'entrée est \textbf{gratuite} pour les spectateurs. Le vainqueur recevra un trophée et une enveloppe de \textbf{50\,000 FCFA}. Nous avons besoin d'une \textbf{affiche au format A4 en orientation portrait} pour annoncer l'événement : elle sera imprimée en 50 exemplaires et collée dans les boutiques, les églises, les écoles et au carrefour principal. Nous disposons du logo de l'association (un dessin simple) et d'une photographie de l'équipe championne de l'année dernière. Merci de nous aider ! »
\end{exemplebox}

\textbf{Moyens mis à disposition :} salle multimédia avec ordinateurs équipés de \textbf{Microsoft Publisher} et \textbf{Scribus}, imprimante, appareil photo numérique.
\textbf{Ressources graphiques fournies :}
\begin{itemize}
\item Photographie de l'équipe : fichier de 300 ko, dimensions $800 \times 600$ pixels, légèrement floue.
\item Logo de l'association : dessin vectoriel (redessiné par un ancien élève).
\end{itemize}

\textbf{Problème :} Comment concevoir, en groupe, une affiche attractive, lisible, conforme aux règles de la PAO, et justifier chacun de vos choix ?

\courssub{Activités d'apprentissage (Travail en groupes de 4 élèves)}

Chaque groupe remettra : la maquette papier, l'affiche imprimée et le texte de justification.

\begin{enumerate}
\item \textbf{Analyse du besoin.} Relis la lettre et dresse la liste complète des informations \emph{obligatoires}. Classe-les en \emph{informations principales} (à mettre en évidence) et \emph{informations secondaires}.
\item \textbf{Choix techniques.} Pour la photo et le logo, précise le type d'image (matricielle ou vectorielle). Laquelle peut être agrandie sans perte de qualité ? Quel risque en agrandissant fortement la photo de $800 \times 600$ pixels ?
\item \textbf{Maquette papier.} Dessine à main levée, sur feuille A4 portrait, la structure de l'affiche : délimite les zones (titre, image, infos pratiques, contact) et note l'emplacement de chaque élément.
\item \textbf{Réalisation sur ordinateur.} Crée l'affiche avec Publisher ou Scribus en respectant : \textbf{deux polices maximum}, image nette et bien cadrée, marges régulières, titre très visible, deux ou trois couleurs harmonieuses.
\item \textbf{Rédaction de la démarche.} Rédige un texte (environ une demi-page) justifiant : format/orientation, logiciel, type d'images, couleurs, polices, hiérarchie.
\item \textbf{Présentation et évaluation.} Présente l'affiche et la démarche (5 min). La classe évalue avec la grille de critères (voir Bilan).
\end{enumerate}

\courssub{Apports théoriques (Ressources à mobiliser)}

\begin{definition}[Infographie et PAO]
L'\textbf{infographie} désigne l'ensemble des techniques de création, de modification et de traitement d'images numériques par ordinateur.
La \textbf{publication assistée par ordinateur (PAO)} est l'utilisation de l'ordinateur pour concevoir des documents destinés à l'impression (affiches, journaux, brochures, livres, cartes de visite) en combinant texte, images et mise en page.
\end{definition}

\begin{definition}[Images matricielles et vectorielles]
Une \textbf{image matricielle} (ou \emph{bitmap}) est une grille de points appelés \cle{pixels}, chacun portant une couleur. Sa qualité dépend de sa \cle{résolution} (nombre de pixels par pouce, ppp/dpi). Si on l'agrandit au-delà de sa résolution native, les pixels deviennent visibles : c'est la \cle{pixellisation} (effet « mosaïque » ou flou). Formats courants : JPEG, PNG, BMP, GIF.

Une \textbf{image vectorielle} est décrite par des formules mathématiques (coordonnées de points, courbes de Bézier, formes géométriques). Elle est \textbf{indépendante de la résolution} : on peut l'agrandir ou la réduire indéfiniment sans aucune perte de netteté. Formats courants : SVG, AI, EPS, PDF (vectoriel).
\end{definition}

\begin{propriete}[Logiciels d'infographie et de PAO]
\begin{center}
\begin{tabularx}{\linewidth}{|l|X|X|}
\hline
\rowcolor{popblueL} \textbf{Catégorie} & \textbf{Exemples (Matriciel / Bitmap)} & \textbf{Exemples (Vectoriel)} \\
\hline
Infographie (création/retouche images) & Paint, GIMP, Photoshop & Inkscape, Illustrator, CorelDRAW \\
\hline
PAO (mise en page documents) & \multicolumn{2}{c|}{Microsoft Publisher, Scribus (libre/gratuit), Canva (en ligne), InDesign} \\
\hline
\end{tabularx}
\end{center}
\emph{Note :} Un logiciel de PAO importe les deux types d'images mais ne les crée pas (sauf formes vectorielles simples).
\end{propriete}

\begin{propriete}[Règles fondamentales de mise en page]
Pour un document lisible, équilibré et professionnel :
\begin{enumerate}
\item \textbf{Marges régulières :} zone blanche d'au moins 1,5 à 2 cm sur les 4 côtés (zone de sécurité pour l'impression).
\item \textbf{Hiérarchie visuelle :} l'information la plus importante (titre, accroche) est la plus visible (taille, gras, couleur, position haute).
\item \textbf{Polices :} \textbf{maximum 2 familles de polices} (ex. : une avec empattements \emph{serif} pour le titre, une sans empattements \emph{sans-serif} pour le corps). Éviter les polices fantaisistes illisibles.
\item \textbf{Couleurs :} palette limitée (2 à 3 couleurs principales + noir/blanc), harmonieuses (cercle chromatique), contraste suffisant fond/texte.
\item \textbf{Équilibre texte/images :} ni trop dense, ni trop vide. Alignements soignés (gauche, centré, justifié). Respect des proportions des images (pas d'étirement).
\end{enumerate}
\end{propriete}

\begin{methode}[Démarche de conception d'un document PAO]
\begin{enumerate}
\item \textbf{Analyser le besoin :} public cible, message clé, support (format, orientation), contraintes (budget, délai), ressources disponibles (textes, images, logos).
\item \textbf{Maquetter (papier/crayon) :} zoning (emplacement des blocs), hiérarchie, flux de lecture (Z ou F), choix provisoires polices/couleurs.
\item \textbf{Réaliser (logiciel PAO) :} créer le document (format, marges), importer/placer images, saisir/mettre en forme textes, appliquer charte graphique.
\item \textbf{Vérifier et corriger :} relecture orthographique, cohérence des informations, qualité images (pas de pixellisation), respect marges, aperçu avant impression.
\item \textbf{Exporter/Imprimer et Justifier :} PDF pour l'impression (haute définition, traits de coupe si pro), impression test. Rédiger la note de justification des choix.
\end{enumerate}
\end{methode}

\courssub{Démarche de résolution guidée et corrigé commenté}

\begin{exoresolu}[Conception de l'affiche du tournoi de Nkol-Eton]
Énoncé : En appliquant la démarche de conception, produire l'affiche du tournoi et rédiger la justification des choix.
\tcblower
\textbf{Étape 1 : Analyse du besoin et liste des informations.}

Lecture attentive de la lettre. Extraction et classification des données :

\begin{center}
\begin{tabularx}{\linewidth}{|l|X|}
\hline
\rowcolor{popblueL} \textbf{Catégorie} & \textbf{Informations obligatoires} \\
\hline
Principales (visibilité max) & Nom : « Tournoi de football de Nkol-Eton » ; Dates : 12 au 26 juillet 2025 ; Lieu : Stade municipal de Nkol-Eton ; Entrée : Gratuite. \\
\hline
Secondaires (lisibles de près) & Horaires : Samedis \& dimanches, 15 h ; Participants : 8 équipes ; Inscription : 5\,000 FCFA/équipe avant le 5/07/2025 ; Contact : M\textsuperscript{me} Abena -- 6\,90\,12\,34\,56 ; Récompense : Trophée + 50\,000 FCFA. \\
\hline
\end{tabularx}
\end{center}

\emph{Justification :} Une affiche se lit en 3 secondes (lecture diagonale). Le \textbf{Quoi, Quand, Où, Combien (entrée)} doivent sauter aux yeux. Le reste (contact, détails inscription, prime) s'adresse aux personnes déjà intéressées qui s'approchent.

\textbf{Étape 2 : Choix et traitement des images.}

\begin{itemize}
\item \textbf{Photo de l'équipe (300 ko, $800 \times 600$ px, floue) :} Image \cle{matricielle}. Résolution faible pour de l'impression A4 (300 ppp requis $\rightarrow$ $2480 \times 3508$ px). Un agrandissement important provoquera une \textbf{pixellisation} visible (flou, blocs carrés). \textbf{Décision :} L'utiliser en \emph{petite taille} (ex. 1/4 de page, $\approx 100 \times 75$ mm $\rightarrow$ $\approx 120$ ppp, acceptable), bien cadrée (recadrage sur les visages/action), sans déformation.
\item \textbf{Logo de l'association :} Image \cle{vectorielle}. Redimensionnable à l'infini sans perte. \textbf{Décision :} Le placer en \emph{grand format} (haut de page, à gauche ou centré) pour l'identité visuelle forte.
\end{itemize}

\textbf{Étape 3 : Maquette papier (Zoning).}

Sur feuille A4 portrait ($210 \times 297$ mm), marges 1,5 cm. Organisation en 4 zones verticales (flux de lecture en Z) :

\begin{popfigure}
\begin{tikzpicture}[scale=0.65, every node/.style={font=\small}]
% Page frame
\draw[popdark, thick] (0,0) rectangle (21,29.7);
% Margins
\draw[popmuted, dashed] (1.5,1.5) rectangle (19.5,28.2);
% Zones
% Zone 1: Header (Logo + Title) - 20% height ~ 6cm
\draw[popblue, thick] (1.5,23.7) rectangle (19.5,28.2);
\node[popblue, align=center] at (10.5,26) {ZONE 1 : EN-TÊTE (≈ 20\%) \newline Logo (G) + TITRE Principal (D/C)};
% Zone 2: Image - 40% height ~ 12cm
\draw[popgreen, thick] (1.5,11.7) rectangle (19.5,23.7);
\node[popgreen, align=center] at (10.5,17.7) {ZONE 2 : IMAGE CENTRALE (≈ 40\%) \newline Photo équipe championne + « Entrée Gratuite »};
% Zone 3: Practical Info - 25% height ~ 7.5cm
\draw[poporange, thick] (1.5,4.2) rectangle (19.5,11.7);
\node[poporange, align=center] at (10.5,7.95) {ZONE 3 : INFOS PRATIQUES (≈ 25\%) \newline Dates, Horaires, Lieu, Inscription};
% Zone 4: Footer (Contact + Rewards) - 15% height ~ 4.5cm
\draw[poppurple, thick] (1.5,1.5) rectangle (19.5,4.2);
\node[poppurple, align=center] at (10.5,2.85) {ZONE 4 : PIED DE PAGE (≈ 15\%) \newline Contact, Récompenses};
\end{tikzpicture}
\legende{Maquette de zoning pour l'affiche A4 Portrait (marges 1,5 cm).}
\end{popfigure}

\textbf{Étape 4 : Réalisation sur ordinateur (Publisher / Scribus).}

\begin{enumerate}
\item Nouveau document : A4, Portrait, Marges 1,5 cm.
\item \textbf{Zones de texte / Images :} Créer 4 cadres (text frames / picture frames) selon le zoning.
\item \textbf{Import images :} Insérer logo (vectoriel $\rightarrow$ agrandir sans peur), photo (matricielle $\rightarrow$ taille réduite, recadrage).
\item \textbf{Typographie :} Police 1 (Titre) : \texttt{Impact} ou \texttt{Arial Black} (gras, condensé, très lisible de loin), Taille 50--60 pt, Couleur Vert foncé. Police 2 (Corps) : \texttt{Arial} ou \texttt{Calibri} (simple, lisible), Taille 14--18 pt, Noir ou Gris foncé. \textbf{Règle des 2 polices respectée.}
\item \textbf{Couleurs :} Vert (pelouse, \#006400) et Jaune (soleil, \#FFD700) sur fond Blanc. Titre Vert, encadrés/points Jaune, Texte Noir.
\item \textbf{Alignements :} Titre centré ou aligné gauche sur logo. Infos pratiques : liste à puces alignée gauche. Contact centré en bas.
\item \textbf{Vérification :} Aperçu avant impression $\rightarrow$ Imprimer 1 exemplaire test.
\end{enumerate}

\textbf{Étape 5 : Rédaction de la justification (Modèle attendu).}

\begin{exemplebox}[Exemple de justification rédigée]
« Notre groupe a choisi le format A4 portrait car c'est le standard le plus économique pour une impression de 50 exemplaires chez un imprimeur local ou au cybercafé du quartier, et il s'adapte aux panneaux d'affichage verticaux (murs, vitrines).

Nous avons utilisé \textbf{Microsoft Publisher} car c'est un logiciel de PAO installé sur les postes du collège, offrant une gestion aisée des cadres de texte et d'images (repérage magnétique, gabarits), adaptée à notre niveau.

Le \textbf{logo vectoriel} a été agrandi pour occuper la largeur haute de l'affiche, garantissant une netteté parfaite. La \textbf{photo matricielle} ($800 \times 600$ px), de faible résolution et floue d'origine, a été recadrée sur l'essentiel (les joueurs) et placée en taille réduite (largeur 10 cm) pour masquer les défauts et éviter la pixellisation à l'impression.

La \textbf{charte graphique} se limite à deux polices (\texttt{Impact} pour l'impact visuel du titre, \texttt{Arial} pour la lecture aisée des détails) et deux couleurs symboliques (Vert pour le football/espérance, Jaune pour la joie/soleil) sur fond blanc. Cela assure un contraste fort, une lisibilité optimale et une identité visuelle cohérente. Les marges de 1,5 cm protègent le contenu lors de la coupe papier. »
\end{exemplebox}

\textbf{Étape 6 : Présentation et évaluation par les pairs.}

Chaque groupe présente (5 min). Évaluation sur 20 points selon la grille :

\begin{center}
\begin{tabularx}{\linewidth}{|l|X|c|}
\hline
\rowcolor{popblueL} \textbf{Critère} & \textbf{Indicateurs de réussite} & \textbf{Points} \\
\hline
Respect des infos obligatoires & Toutes les infos principales et secondaires présentes, exactes, orthographe soignée. & 4 \\
\hline
Lisibilité \& Hiérarchie & Titre visible de loin (3 m), lecture fluide (Z/F), contraste fond/texte, infos principales distinctes. & 4 \\
\hline
Qualité des images & Logo net (vectoriel), photo non pixellisée (taille adaptée), recadrage pertinent, proportions respectées. & 4 \\
\hline
Règles de mise en page & 2 polices max, 2--3 couleurs harmonieuses, marges régulières, alignements soignés, équilibre texte/image. & 4 \\
\hline
Qualité de la justification & Texte structé, vocabulaire technique précis (matriciel/vectoriel, PAO, zoning, police, marge), arguments logiques. & 4 \\
\hline
\end{tabularx}
\end{center}
\end{exoresolu}

\courssub{Erreurs fréquentes à éviter}

\begin{attention}[Pièges classiques en PAO (niveau 3ème)]
\begin{itemize}
\item \textbf{Oubli d'informations critiques :} Date, lieu, contact manquants $\rightarrow$ affiche inutile. \emph{Check-list obligatoire avant impression.}
\item \textbf{Agrandissement abusif d'image matricielle :} Photo $800 \times 600$ px étirée sur toute la page $\rightarrow$ flou disgracieux, aspect « amateur ». \emph{Toujours vérifier la résolution effective après placement.}
\item \textbf{Festival de polices et de couleurs :} 4--5 polices différentes, arc-en-ciel de couleurs $\rightarrow$ illisible, visuel agressif, pas de hiérarchie. \emph{Règle d'or : 2 polices, 3 couleurs max.}
\item \textbf{Absence de marges (texte au bord) :} Risque de coupe à l'impression (perte d'infos), sensation d'étouffement visuel. \emph{Marge minimale 1,5 cm.}
\item \textbf{Déformation des images :} Étirement horizontal/vertical pour « remplir le cadre » $\rightarrow$ cercles $\rightarrow$ ellipses, visages déformés. \emph{Verrouiller les proportions (Shift + coin dans la plupart des logiciels).}
\item \textbf{Justification approximative :} « C'est joli », « J'aime bien » $\rightarrow$ note faible. \emph{Exiger : vocabulaire technique + lien avec le besoin client + contrainte technique.}
\end{itemize}
\end{attention}

\courssub{Bilan et compétences visées}

Cette leçon t'a permis de mobiliser l'ensemble des savoirs de l'unité « Infographie et PAO » dans une \textbf{situation authentique} de communication locale.

\begin{aretenir}[Ce qu'il faut retenir pour le BEPC]
\begin{itemize}
\item \textbf{Infographie} = traitement images ; \textbf{PAO} = mise en page documents imprimés.
\item \textbf{Matriciel (pixels)} : qualité fixe, pixellisation si agrandi (JPEG, PNG). \textbf{Vectoriel (maths)} : qualité infinie, redimensionnement libre (SVG, AI).
\item \textbf{Logiciels PAO :} Publisher, Scribus, Canva. \textbf{Logiciels Infographie :} GIMP (matriciel), Inkscape (vectoriel).
\item \textbf{Mise en page pro :} Marges, Hiérarchie, 2 polices max, 3 couleurs max, Alignements, Proportions images.
\item \textbf{Démarche :} Besoin $\rightarrow$ Maquette $\rightarrow$ Réalisation $\rightarrow$ Vérification $\rightarrow$ Justification.
\item \textbf{Compétence clé BEPC :} \textbf{Justifier sa démarche} à l'écrit et à l'oral avec un vocabulaire technique précis. Ce n'est pas le résultat seul qui compte, mais la preuve que tu as fait des \emph{choix raisonnés}.
\end{itemize}
\end{aretenir}

\begin{savaistu}[Le numérique au Cameroun : La PAO au service de la vie locale]
Au Cameroun, la PAO est un levier d'employabilité pour les jeunes. De la conception d'affiches pour les événements culturels (Ngondo, Fête de la Musique, tournois inter-quartiers comme Nkol-Eton) à la création de cartes de visite pour les artisans (menuisiers, couturières, mécaniciens) en passant par les bulletins paroissiaux ou les journaux scolaires, la maîtrise de Scribus (gratuit) ou Publisher permet de lancer une micro-activité de \emph{design graphique} avec un simple ordinateur portable. De nombreux cybercafés de Yaoundé, Douala, Bafoussam ou Garoua proposent ce service. La rigueur technique (résolution images, marges, PDF/X pour l'imprimeur) distingue le professionnel de l'amateur.
\end{savaistu}

\newpage

\coursec{Méthodes et exercices résolus}

\coursec{Leçon 07 : Initiation à la publication assistée par ordinateur}

\begin{definition}[Infographie]
L'\cle{infographie} (informatique graphique) désigne l'ensemble des techniques de création, de modification et de traitement des \cle{images numériques} par ordinateur. Elle concerne l'image prise isolément : dessin, retouche photo, création de logo, schéma.
\end{definition}

\begin{definition}[Publication assistée par ordinateur (PAO)]
La \cle{PAO} (Publication Assistée par Ordinateur) est l'ensemble des techniques de \cle{mise en page} de documents destinés à l'impression ou à la diffusion numérique. Elle combine du \cle{texte}, des \cle{images} et des \cle{objets graphiques} (filets, cadres, fonds de couleur) dans une composition maîtrisée (journal, flyer, affiche, brochure, dépliant, bulletin).
\end{definition}

\begin{propriete}[Distinction fondamentale]
\begin{itemize}
\item \textbf{Infographie} $\rightarrow$ Objet produit : une \emph{image} (fichier image : .jpg, .png, .svg). Outils : éditeurs d'images (GIMP, Paint, Inkscape, Photoshop).
\item \textbf{PAO} $\rightarrow$ Objet produit : un \emph{document complet} (fichier de mise en page : .pub, .sla, .indd, ou .pdf final). Outils : logiciels de PAO (Publisher, Scribus, Canva, InDesign).
\end{itemize}
\textbf{Règle pratique :} on fait souvent de l'infographie \emph{avant} la PAO (préparer les images), puis on importe ces images dans le logiciel de PAO pour composer le document.
\end{propriete}

\courssub{Savoir-faire 1 : Distinguer infographie et PAO dans une situation concrète}

\begin{methode}[Identifier le domaine concerné : infographie ou PAO]
Pour appliquer les notions de l'unité à une situation de la vie courante, il faut d'abord reconnaître de quelle activité il s'agit.
\begin{enumerate}
\item \textbf{Repérer l'objet à produire} : s'agit-il de créer ou de retoucher une \cle{image} (logo, photo, affiche visuelle) ou de composer un \cle{document complet} (journal, flyer, bulletin, brochure) ?
\item \textbf{Appliquer la définition} :
\begin{itemize}
\item création/traitement d'images par ordinateur $\Rightarrow$ \cle{infographie} ;
\item mise en page de documents combinant texte et images en vue d'une impression ou d'une diffusion $\Rightarrow$ \cle{PAO}.
\end{itemize}
\item \textbf{Identifier le type d'image} si besoin : photo ou dessin complexe $\Rightarrow$ image \cle{matricielle} (bitmap, faite de pixels) ; logo ou schéma redimensionnable sans perte $\Rightarrow$ image \cle{vectorielle} (faite d'objets géométriques).
\item \textbf{Choisir le logiciel adapté} : logiciel de dessin/retouche pour l'infographie (Paint, GIMP, Inkscape), logiciel de PAO pour la mise en page (Microsoft Publisher, Scribus, Canva).
\item \textbf{Justifier la réponse} en une phrase complète qui cite la définition et relie au cas concret.
\end{enumerate}
\end{methode}

\begin{exoresolu}[Le projet du club jeunesse de Yaoundé]
Le club jeunesse d'un quartier de Yaoundé veut : (a) retoucher la photo prise lors de sa dernière journée sportive ; (b) composer un dépliant de deux pages annonçant sa journée portes ouvertes, avec textes, photos et logo. Pour chaque tâche, indique s'il s'agit d'infographie ou de PAO, et propose un logiciel adapté. Justifie tes réponses.
\tcblower
\textbf{Raisonnement.} On applique les définitions : l'infographie concerne la création et le traitement des images ; la PAO concerne la mise en page d'un document complet combinant texte et images.

\textbf{Tâche (a) : retoucher la photo.}
\begin{itemize}
\item Il s'agit de modifier une \emph{image} (corriger la luminosité, recadrer, effacer un élément gênant). C'est donc de l'\cle{infographie}, plus précisément le traitement d'une image \cle{matricielle} (une photo est composée de pixels).
\item Logiciel adapté : un logiciel de retouche d'images comme \textbf{GIMP} (gratuit, multiplateforme) ou \textbf{Paint} (simple, livré avec Windows).
\end{itemize}

\textbf{Tâche (b) : composer le dépliant.}
\begin{itemize}
\item Il s'agit d'organiser sur deux pages du \emph{texte}, des \emph{photos} et un \emph{logo} en vue d'une impression : c'est de la mise en page, donc de la \cle{PAO}.
\item Logiciel adapté : un logiciel de PAO comme \textbf{Microsoft Publisher}, \textbf{Scribus} (gratuit, open source) ou \textbf{Canva} (en ligne, version gratuite).
\end{itemize}

\textbf{Conclusion.} La retouche de la photo relève de l'infographie (logiciel GIMP ou Paint), car on traite une image matricielle ; la composition du dépliant relève de la PAO (logiciel Publisher, Scribus ou Canva), car on met en page un document complet mêlant texte et images.
\end{exoresolu}

\courssub{Savoir-faire 2 : Choisir entre image matricielle et image vectorielle}

\begin{definition}[Image matricielle (Bitmap)]
Une \cle{image matricielle} est une grille rectangulaire de points appelés \cle{pixels} (picture elements). Chaque pixel porte une couleur codée (niveau de gris, RVB, CMJN). La taille de l'image est fixe : largeur $\times$ hauteur en pixels. \textbf{Exemples :} photographie numérique, capture d'écran, dessin réaliste avec dégradés. \textbf{Formats courants :} JPEG (compression avec perte, photo), PNG (compression sans perte, transparence), BMP, GIF (animation).
\end{definition}

\begin{definition}[Image vectorielle]
Une \cle{image vectorielle} est décrite par des \cle{objets géométriques} (points, segments, courbes de Bézier, cercles, polygones) définis par des équations mathématiques (coordonnées, équations de droites/courbes, couleurs de remplissage et de contour). Elle n'a pas de résolution fixe. \textbf{Exemples :} logo, pictogramme, schéma technique, police de caractères, carte. \textbf{Formats courants :} SVG (standard web), WMF/EMF (Windows), EPS (imprimerie), PDF (peut contenir du vectoriel).
\end{definition}

\begin{propriete}[Comportement à l'agrandissement]
\begin{itemize}
\item \textbf{Image matricielle :} agrandissement $\Rightarrow$ \cle{pixellisation} (les pixels deviennent visibles, l'image perd en netteté, effet « escalier » sur les contours). La qualité dépend de la \cle{résolution} initiale (ppp / dpi : points par pouce). Pour l'impression : 300 ppp minimum ; pour l'écran : 72--96 ppp.
\item \textbf{Image vectorielle :} agrandissement $\Rightarrow$ \textbf{aucune perte de qualité}. L'ordinateur recalcule les objets géométriques à la nouvelle taille. Idéal pour les logos (carte de visite $\leftrightarrow$ banderole 4 m).
\end{itemize}
\end{propriete}

\begin{methode}[Choisir le bon type d'image]
\begin{enumerate}
\item \textbf{Analyser la nature de l'image} : photo, dessin réaliste avec dégradés $\Rightarrow$ matricielle ; logo, schéma, texte stylisé, formes géométriques $\Rightarrow$ vectorielle.
\item \textbf{Analyser l'usage prévu} : l'image devra-t-elle être fortement agrandie (banderole, affiche) ? Si oui, préférer le \cle{vectoriel}, qui ne perd pas en qualité à l'agrandissement.
\item \textbf{Connaître les formats} : matriciel $\Rightarrow$ JPEG, PNG, BMP, GIF ; vectoriel $\Rightarrow$ SVG, WMF, EPS, PDF (vectoriel).
\item \textbf{Anticiper le défaut} : une image matricielle agrandie devient \cle{pixellisée} ; une image vectorielle reste nette car elle est recalculée à partir de formules géométriques.
\item \textbf{Rédiger la justification} : citer la nature de l'image, l'usage et la conséquence sur la qualité.
\end{enumerate}
\end{methode}

\begin{popfigure}
\begin{tikzpicture}[
    scale=0.8,
    every node/.style={font=\small},
    pixel/.style={draw=popline, fill=popblue!20, minimum size=0.6cm},
    vector/.style={draw=poporange, thick, fill=poporange!10},
    curve/.style={draw=poporange, very thick, rounded corners=2pt}
]
% Matricielle
\begin{scope}[xshift=0cm]
\node[align=center, text width=4cm] at (2.5,5.5) {\textbf{Image Matricielle (Pixels)}};
% Grille originale
\foreach \x in {0,...,4} \foreach \y in {0,...,3} {
    \pgfmathsetmacro{\c}{random(0,100)}
    \node[pixel, fill=popblue!\c] at (\x*0.65, -\y*0.65) {};
}
\node[align=center, text width=4cm] at (2.5,-3.5) {Zoom 400\% : pixels visibles};
% Zoom simulé
\begin{scope}[xshift=7cm, yshift=-1cm, scale=2.5]
\foreach \x in {1,...,2} \foreach \y in {1,...,2} {
    \pgfmathsetmacro{\c}{random(0,100)}
    \node[pixel, fill=popblue!\c, minimum size=0.65cm] at (\x*0.65, -\y*0.65) {};
}
\end{scope}
\draw[->, thick, popdark] (5.5, -1.5) -- (7, -1.5) node[midway, above] {Agrandissement};
\end{scope}
% Vectorielle
\begin{scope}[xshift=16cm]
\node[align=center, text width=4cm] at (2.5,5.5) {\textbf{Image Vectorielle (Objets)}};
% Formes vectorielles
\draw[vector] (0.5, 0) rectangle (4.5, 3);
\draw[curve] (1, 1.5) .. controls (2, 3) and (3, 3) .. (4, 1.5);
\draw[vector] (2.5, 0.5) circle (0.8);
\node[align=center, text width=4cm] at (2.5,-1.5) {Zoom 400\% : contours nets};
% Zoom simulé (recalculé)
\begin{scope}[xshift=7cm, yshift=-1cm, scale=2.5]
\draw[vector] (0.2, 0) rectangle (1.8, 1.2);
\draw[curve] (0.4, 0.6) .. controls (0.8, 1.2) and (1.2, 1.2) .. (1.6, 0.6);
\draw[vector] (1, 0.2) circle (0.32);
\end{scope}
\draw[->, thick, popdark] (5.5, -0.5) -- (7, -0.5) node[midway, above] {Agrandissement};
\end{scope}
\end{tikzpicture}
\legende{Comparaison visuelle : l'image matricielle (pixels) se dégrade à l'agrandissement (pixellisation), l'image vectorielle (objets géométriques) reste parfaitement nette.}
\end{popfigure}

\begin{exoresolu}[Le logo de la boutique « Chez Manga »]
Monsieur Manga, commerçant à Douala, veut créer le logo de sa boutique. Ce logo figurera sur les cartes de visite (petite taille) mais aussi sur une grande banderole de 4 mètres. Son neveu lui propose de le dessiner comme une photo au format JPEG. Ce choix est-il judicieux ? Justifie et propose une solution.
\tcblower
\textbf{Analyse de la situation.}
\begin{itemize}
\item Un logo est composé de formes simples (courbes, lettres, aplats de couleur) : il se prête bien au mode \cle{vectoriel}.
\item L'usage impose un \emph{fort agrandissement} : passer d'une carte de visite ($\approx$ 5 cm) à une banderole de 4 mètres (facteur 80).
\end{itemize}

\textbf{Conséquence du choix JPEG.} Le JPEG est un format d'image \cle{matricielle} : l'image est une grille de pixels de taille fixe. En l'agrandissant pour la banderole, chaque pixel deviendra un carré visible : le logo sera \emph{pixellisé} et flou, ce qui donnera une image de mauvaise qualité pour la boutique.

\textbf{Solution proposée.} Créer le logo avec un logiciel de dessin \emph{vectoriel} (par exemple \textbf{Inkscape}, gratuit et open source) et l'enregistrer au format \textbf{SVG} (Scalable Vector Graphics). Une image vectorielle est décrite par des objets géométriques (points, droites, courbes de Bézier) : l'ordinateur la recalcule à chaque taille, donc elle reste parfaitement nette, aussi bien sur la carte de visite que sur la banderole.

\textbf{Conclusion.} Le choix du JPEG n'est pas judicieux : un logo destiné à être fortement agrandi doit être réalisé en image vectorielle (format SVG, logiciel Inkscape), car le vectoriel se redimensionne sans aucune perte de qualité, contrairement au matriciel qui se pixellise.
\end{exoresolu}

\begin{aretenir}[Rappel rapide : Matriciel vs Vectoriel]
\begin{center}
\begin{tabularx}{\linewidth}{|l|X|X|}\hline
\rowcolor{popblueL}
\textbf{Critère} & \textbf{Image Matricielle} & \textbf{Image Vectorielle} \\\hline
\textbf{Structure} & Grille de pixels & Objets géométriques (formules) \\\hline
\textbf{Adaptée à} & Photos, images complexes, dégradés & Logos, schémas, textes, pictogrammes \\\hline
\textbf{Agrandissement} & Perte de qualité (pixellisation) & \textbf{Aucune perte} (recalcul mathématique) \\\hline
\textbf{Résolution} & Fixe (ppp / dpi) & Indépendante de la résolution \\\hline
\textbf{Formats} & JPEG, PNG, GIF, BMP & SVG, WMF, EPS, PDF (vectoriel) \\\hline
\textbf{Logiciels} & GIMP, Paint, Photoshop & Inkscape, Illustrator, Draw (LibreOffice) \\\hline
\end{tabularx}
\end{center}
\end{aretenir}

\courssub{Savoir-faire 3 : Identifier les éléments d'une mise en page}

\begin{definition}[Éléments constitutifs d'une page de PAO]
Une page de PAO structurée comporte les éléments suivants :
\begin{itemize}
\item \textbf{Les \cle{marges} :} espaces blancs de réserve entre le bord physique du papier et la zone utile (zone de texte/images). Elles assurent la lisibilité et permettent la reliure/le massicotage.
\item \textbf{Les \cle{colonnes} :} découpage vertical de la zone utile en plusieurs bandes étroites (facilite la lecture, structure l'information : journaux, dépliants).
\item \textbf{Le \cle{titre} et les \cle{sous-titres} :} textes hiérarchisés par la taille, la police, la graisse (gras) et la couleur pour guider l'œil.
\item \textbf{L'\cle{en-tête} (header) :} zone répétée en haut de chaque page (nom du document, logo, date, numéro de version).
\item \textbf{Le \cle{pied de page} (footer) :} zone répétée en bas de chaque page (numérotation, mentions légales, contact, date).
\item \textbf{La \cle{numérotation des pages} :} indispensable pour les documents multipages (souvent dans le pied de page).
\item \textbf{Les \cle{images} et \cle{objets graphiques} :} photos (matricielles), logos/schémas (vectoriels), filets (lignes), cadres, fonds de couleur, formes.
\item \textbf{L'\cle{habillage} (wrap) :} mode de circulation du texte autour d'une image (carré, serré, au-dessus/au-dessous, traversant).
\end{itemize}
\end{definition}

\begin{methode}[Reconnaître et nommer les éléments d'une page de PAO]
Face à un document (journal, affiche, dépliant), savoir nommer ses constituants :
\begin{enumerate}
\item \textbf{Observer les bords} : identifier les \cle{marges} (espace vide périphérique).
\item \textbf{Repérer les zones répétées} : haut de page $\rightarrow$ \cle{en-tête} ; bas de page $\rightarrow$ \cle{pied de page} (souvent avec \cle{numérotation}).
\item \textbf{Analyser la structure du corps} : nombre de \cle{colonnes}, emplacement des \cle{titres} et \cle{sous-titres}.
\item \textbf{Localiser les éléments graphiques} : \cle{images}, \cle{logos}, \cle{filets}, \cle{cadres}, \cle{habillage} du texte.
\item \textbf{Décrire de l'extérieur vers l'intérieur} : marges $\rightarrow$ en-tête/pied $\rightarrow$ colonnes $\rightarrow$ blocs contenu.
\end{enumerate}
\end{methode}

\begin{popfigure}
\begin{tikzpicture}[
    scale=0.7,
    every node/.style={font=\small},
    pageborder/.style={draw=popdark, thick, rectangle, minimum width=14cm, minimum height=10cm},
    margin/.style={draw=popblue, dashed, fill=popblue!10},
    zone/.style={draw=poporange, fill=poporange!15, rounded corners=3pt},
    textline/.style={draw=popmuted, densely dashed, minimum height=0.3cm},
    img/.style={draw=popgreen, fill=popgreen!20, pattern=north east lines, pattern color=popgreen!50}
]
% Page physique
\node[pageborder] (page) at (0,0) {};
% Marges
\node[margin, minimum width=13cm, minimum height=9cm] at (0,0) {};
\node[align=center, text width=1.5cm] at (-6.5, 0) {\textbf{Marges}\\(blancs de réserve)};
% En-tête
\node[zone, minimum width=12cm, minimum height=1.2cm] (header) at (0, 3.9) {};
\node[align=center] at (0, 3.9) {\textbf{En-tête} : Nom journal, Date, Logo};
% Pied de page
\node[zone, minimum width=12cm, minimum height=0.8cm] (footer) at (0, -3.9) {};
\node[align=center] at (0, -3.9) {\textbf{Pied de page} : N° page \quad Mentions légales};
% Colonnes (2)
\node[zone, minimum width=5.5cm, minimum height=5.5cm] (col1) at (-3.2, 0.2) {};
\node[zone, minimum width=5.5cm, minimum height=5.5cm] (col2) at (3.2, 0.2) {};
% Titre dans col1
\node[textline, minimum width=4.5cm, fill=poppink!30] at (-3.2, 2.8) {\textbf{Titre principal (Gras, Grande taille)}};
% Texte lignes
\foreach \i in {1,...,6} {
    \node[textline, minimum width=4.5cm] at (-3.2, 2.8-\i*0.6) {};
    \node[textline, minimum width=4.5cm] at (3.2, 2.8-\i*0.6) {};
}
% Image dans col2 (habillage)
\node[img, minimum width=4.5cm, minimum height=2.5cm] at (3.2, 2.2) {};
\node at (3.2, 2.2) {\textbf{Image / Photo}};
\node[align=center, text width=4.5cm] at (3.2, 0.8) {Texte habillé\\autour de l'image};
% Flèches légendes
\draw[->, popblue, thick] (-7.5, 2) -- (-6.2, 1);
\draw[->, poporange, thick] (0, 4.6) -- (0, 4.1);
\draw[->, poporange, thick] (0, -4.6) -- (0, -4.1);
\draw[->, popgreen, thick] (7, 2.2) -- (5.5, 2.2);
\end{tikzpicture}
\legende{Anatomie d'une page de PAO (ex. journal scolaire) : marges, en-tête, pied de page, colonnes, titre, image avec habillage.}
\end{popfigure}

\begin{exoresolu}[Analyse du journal du collège]
Le journal du collège Biyem-Assi est composé ainsi : une bande blanche de 2 cm tout autour de chaque page ; le nom du journal et la date en haut de chaque page ; le texte des articles réparti sur deux colonnes ; un grand titre en gras au début de chaque article ; des photos insérées entre les articles ; un numéro en bas de chaque page. Identifie et nomme chaque élément de mise en page décrit.
\tcblower
On décrit le document de l'extérieur vers l'intérieur :
\begin{enumerate}
\item « Bande blanche de 2 cm tout autour de la page » $\Rightarrow$ ce sont les \cle{marges}.
\item « Nom du journal et date en haut de chaque page » $\Rightarrow$ c'est l'\cle{en-tête} (zone répétée sur toutes les pages).
\item « Texte réparti sur deux colonnes » $\Rightarrow$ mise en page en \cle{colonnes}.
\item « Grand titre en gras au début de chaque article » $\Rightarrow$ le \cle{titre} (avec mise en forme des caractères : graisse et taille).
\item « Photos insérées entre les articles » $\Rightarrow$ les \cle{images} (objets graphiques incorporés, ici matricielles).
\item « Numéro en bas de chaque page » $\Rightarrow$ la \cle{numérotation des pages}, placée dans le \cle{pied de page}.
\end{enumerate}
\textbf{Conclusion.} Le journal utilise les six éléments classiques d'une mise en page de PAO : marges, en-tête, colonnes, titres mis en forme, images et pied de page avec numérotation.
\end{exoresolu}

\courssub{Savoir-faire 4 : Rédiger et justifier une démarche de conception d'un document de PAO}

\begin{methode}[Concevoir un document de PAO en 5 étapes justifiées]
Pour rédiger une démarche complète et la justifier (compétence « Rédiger et justifier une démarche ») :
\begin{enumerate}
\item \textbf{Définir l'objectif et le public cible} : à quoi sert le document (informer, inviter, sensibiliser, vendre) ? À qui s'adresse-t-il (âge, niveau de lecture, contexte) ? \emph{Justification :} « On choisit une affiche visuelle car le public (élèves pressés) ne lira pas un long texte. »
\item \textbf{Choisir le support et le format} : format papier (A4, A5, A3), orientation (portrait/paysage), nombre de pages, type de pliage (dépliant 2 volets, 3 volets roulés/accordéon). \emph{Justification :} « Format A4 portrait : standard imprimante collège, adapté à l'affichage vertical. »
\item \textbf{Rassembler et préparer le contenu (Pré-PAO)} : rédiger les textes (titres, corps, légendes), sélectionner les images, \emph{retoucher les images en infographie} (recadrage, luminosité, conversion vectoriel/logo) \textbf{avant} importation. \emph{Justification :} « On retouche la photo sous GIMP pour corriger la surexposition ; on vectorise le logo pour l'agrandissement. »
\item \textbf{Réaliser la mise en page (PAO)} : créer le gabarit (marges, colonnes, repères), insérer blocs de texte et images, appliquer les styles (polices : max 2 familles, tailles hiérarchisées, couleurs cohérentes), gérer l'habillage, placer en-tête/pied de page/numérotation. \emph{Justification :} « On limite à 2 polices (sans serif pour titres, serif pour corps) pour la lisibilité ; couleurs vert/blanc pour thème environnement. »
\item \textbf{Vérifier, valider et diffuser} : relecture orthographique/typographique, vérification alignements/débordements (aperçu avant impression), impression test, export PDF (qualité impression ou web), archivage des sources. \emph{Justification :} « Aperçu avant impression évite le gaspillage papier ; PDF assure la fidélité sur tout poste. »
\end{enumerate}
\textbf{Réflexe rédactionnel :} à chaque étape, écrire une phrase « On fait X \emph{parce que} Y (lien avec objectif/public/contraintes techniques) ».
\end{methode}

\begin{exoresolu}[L'affiche de la campagne de propreté]
Le club environnement de ton collège veut produire une affiche A4 pour sensibiliser les élèves au tri des déchets dans la cour. Rédige la démarche complète de conception de cette affiche en justifiant chaque choix.
\tcblower
\textbf{Étape 1 — Objectif et public.} L'affiche doit \emph{sensibiliser} les élèves du collège au tri des déchets. Le public étant des élèves pressés qui passent dans la cour, le message doit être court, percutant et très visuel : cela justifie le choix d'une \emph{affiche} (format unique, lecture rapide) plutôt qu'un long texte ou un dépliant.

\textbf{Étape 2 — Support et format.} On choisit le format \textbf{A4} (21 $\times$ 29,7 cm) en orientation \textbf{portrait}, car c'est le format natif des imprimantes du collège, économique en papier, et il se prête bien à l'affichage vertical sur les panneaux d'affichage des couloirs.

\textbf{Étape 3 — Préparation du contenu.} On rédige un slogan court (ex. « Trie tes déchets, protège ton école ! »), on sélectionne une photo des poubelles de tri du collège et on la retouche avec un logiciel d'infographie (\textbf{GIMP}) pour améliorer sa luminosité et la recadrer ; on prépare aussi le logo du club en image \textbf{vectorielle} (Inkscape, format SVG) pour qu'il reste net quel que soit son emplacement.

\textbf{Étape 4 — Mise en page.} Dans un logiciel de PAO (\textbf{Scribus}, \textbf{Publisher} ou \textbf{Canva}) :
\begin{itemize}
\item Gabarit : marges 1,5 cm, une seule colonne (affiche = 1 bloc visuel).
\item Placement : slogan en grand titre (police \texttt{Impact} ou \texttt{Arial Black}, 48 pt, couleur \texttt{vert foncé}) en haut ; photo centrée, large ; 3 consignes courtes en dessous (puces, police \texttt{Arial} 18 pt) ; logo club en bas à droite.
\item Harmonie : palette limitée (vert, blanc, noir), 2 polices max, alignements respectés (grille de repères).
\end{itemize}

\textbf{Étape 5 — Vérification et diffusion.} On relit l'orthographe (slogan, consignes), on utilise l'\textbf{aperçu avant impression} pour vérifier qu'aucun élément ne dépasse dans les marges (zone de coupe), on imprime un exemplaire test (épreuve), puis on lance l'impression des 50 exemplaires ou on exporte en \textbf{PDF haute qualité (300 ppp)} pour diffusion numérique sur l'ENT du collège.

\textbf{Conclusion.} La démarche de conception suit cinq étapes justifiées : définir l'objectif/le public, choisir le format, préparer le contenu (infographie amont), réaliser la mise en page (PAO), puis vérifier/diffuser. Chaque choix (A4 portrait, slogan court, photo retouchée, logo vectoriel, couleurs harmonieuses) découle de l'objectif de sensibilisation rapide et des contraintes techniques (imprimante A4, affichage mural).
\end{exoresolu}

\courssub{Savoir-faire 5 : Choisir le logiciel adapté à une tâche}

\begin{propriete}[Panorama des logiciels courants (niveau 3ème)]
\begin{center}
\begin{tabularx}{\linewidth}{|l|X|l|}\hline
\rowcolor{popblueL}
\textbf{Catégorie} & \textbf{Exemples (Gratuits / Payants)} & \textbf{Usage type} \\\hline
\textbf{Retouche matricielle} & \textbf{GIMP} (gratuit, référence), Paint (basique), Photoshop (pro) & Photos : recadrage, couleur, filtres, suppression d'éléments \\\hline
\textbf{Dessin vectoriel} & \textbf{Inkscape} (gratuit, référence), Illustrator (pro), Draw (LibreOffice) & Logos, schémas, plans, pictogrammes, typographie \\\hline
\textbf{PAO (Mise en page)} & \textbf{Scribus} (gratuit, pro), \textbf{Canva} (en ligne, freemium), \textbf{Publisher} (Microsoft, payant), InDesign (pro) & Journaux, flyers, affiches, dépliants, brochures, livres \\\hline
\textbf{Traitement de texte} & Word, Writer (LibreOffice), Google Docs & Lettres, rapports, CV, documents \emph{principalement textuels} \\\hline
\end{tabularx}
\end{center}
\textbf{Attention :} Un traitement de texte (Word) \emph{peut} faire de la mise en page simple, mais la PAO offre : blocs déplaçables librement (calques), gestion native des colonnes complexes, habillage précis des images, gabarits de pages (maîtres), repères d'alignement, export PDF/X pour imprimerie.
\end{propriete}

\begin{methode}[Associer tâche et logiciel]
\begin{enumerate}
\item \textbf{Classer la tâche} :
\begin{itemize}
\item Retouche photo, correction couleurs, suppression fond $\Rightarrow$ \textbf{Éditeur matriciel} (GIMP, Paint).
\item Création logo, schéma, pictogramme, texte courbé $\Rightarrow$ \textbf{Éditeur vectoriel} (Inkscape).
\item Composition document multipages, colonnes, pliage, blocs libres $\Rightarrow$ \textbf{Logiciel PAO} (Scribus, Publisher, Canva).
\item Rédaction lettre, rapport, dissertation $\Rightarrow$ \textbf{Traitement de texte} (Word, Writer).
\end{itemize}
\item \textbf{Ne pas confondre} PAO et Traitement de texte : la PAO gère la \emph{géométrie de la page} (blocs objets), le traitement de texte gère le \emph{flux de texte}.
\item \textbf{Tenir compte des contraintes} : gratuit/payant, installé/en ligne, OS (Windows/Linux/Mac), niveau de l'utilisateur.
\item \textbf{Justifier} le choix : « Je choisis \texttt{Logiciel X} \emph{car} c'est un \texttt{Type} adapté à \texttt{Tâche} (fonction Y). »
\end{enumerate}
\end{methode}

\begin{exoresolu}[Équiper le cybercafé « Le Bon Plan »]
La gérante du cybercafé « Le Bon Plan » à Bafoussam veut équiper ses postes pour trois services : (1) retoucher les photos d'identité des clients ; (2) créer des logos et cartes de visite ; (3) composer des menus et des affiches publicitaires pour les restaurants du quartier. Son budget est limité. Propose un logiciel adapté à chaque service en justifiant, en privilégiant si possible des solutions gratuites.
\tcblower
\textbf{Service 1 — Retouche de photos d'identité.} Une photo d'identité est une image \cle{matricielle} ; il faut un éditeur capable de recadrer aux normes (35$\times$45 mm), corriger l'exposition, le fond. Choix : \textbf{GIMP}, gratuit, complet, multiplateforme (Windows/Linux), gère les calques et la résolution d'impression (300 ppp). \textit{Alternative basique : Paint (recadrage seulement).}

\textbf{Service 2 — Logos et cartes de visite.} Un logo doit pouvoir être agrandi (carte 5 cm $\rightarrow$ enseigne 1 m) sans perte : il faut un éditeur \cle{vectoriel}. Choix : \textbf{Inkscape}, gratuit, standard SVG, gère le texte courbé, les dégradés, l'export PDF vectoriel pour l'imprimeur.

\textbf{Service 3 — Menus et affiches.} Composer un document mêlant textes, photos, mise en page soignée, pliage (menu 2 volets) relève de la \cle{PAO}. Choix : \textbf{Scribus} (gratuit, pro, gestion CMJN, PDF/X pour imprimerie) ou \textbf{Canva} (en ligne, modèles prêts à l'emploi, collaboration simple, version gratuite suffisante). \textit{Alternative payante : Microsoft Publisher (souvent déjà sur poste Windows).}

\textbf{Conclusion.} Avec la suite \textbf{GIMP + Inkscape + Scribus (ou Canva)}, le cybercafé couvre ses trois services avec des logiciels \textbf{gratuits}, standards, adaptés chacun à la nature technique de la tâche : traitement d'image matricielle, création vectorielle, mise en page professionnelle.
\end{exoresolu}

\courssub{Savoir-faire 6 : Rédiger une réponse argumentée et conclure (Méthodologie BEPC)}

\begin{methode}[Structurer une réponse justifiée au BEPC]
Pour satisfaire à la compétence « Rédiger et justifier une démarche, une réponse ou une conclusion », applique la structure en 4 temps :
\begin{enumerate}
\item \textbf{Annonce claire de la réponse} (thèse) dès la première phrase. \textit{Ex : « Jean a raison : le logiciel de PAO est le choix pertinent. »}
\item \textbf{Règle / Définition du cours} qui fonde l'argument. \textit{Ex : « Or, la PAO est spécialisée dans la mise en page de documents combinant texte et images avec gestion des blocs, colonnes et habillage, contrairement au traitement de texte orienté flux linéaire. »}
\item \textbf{Application au cas concret} de l'énoncé (liaison règle $\rightarrow$ situation). \textit{Ex : « Ici, le dépliant nécessite un pliage 3 volets, des colonnes, des images habillées : fonctions natives de Publisher, difficiles et instables sous Word. »}
\item \textbf{Conclusion formelle} qui reprend la question. \textit{Ex : « Par conséquent, Jean a raison : le dépliant doit être réalisé avec un logiciel de PAO. »}
\end{enumerate}
\textbf{Sanction :} Une réponse sans justification (définition + application) ne vaut que la moitié des points. Toujours écrire des phrases complètes, vocabulaire technique précis.
\end{methode}

\begin{exoresolu}[Débat : Word ou Publisher pour le dépliant ?]
Deux élèves de 3ème doivent réaliser le dépliant de la fête de fin d'année. Awa propose d'utiliser le traitement de texte Word ; Jean propose le logiciel de PAO Publisher. Qui a raison ? Rédige une réponse argumentée et conclus.
\tcblower
\textbf{Réponse annoncée.} Jean a raison : pour composer un dépliant, le logiciel de PAO (Publisher) est le mieux adapté.

\textbf{Règle du cours.} Le traitement de texte (Word, Writer) est conçu pour rédiger des documents essentiellement textuels à flux continu (lettres, rapports). La PAO (Publisher, Scribus, Canva) est spécialisée dans la \emph{mise en page} : elle permet de placer librement des blocs de texte et d'images (objets flottants), de gérer facilement les colonnes, les gabarits de pliage (dépliant 2/3 volets) et l'habillage précis des images.

\textbf{Application au cas.} Le dépliant de la fête combine un titre accrocheur, des photos, un programme en colonnes et un pliage en trois volets : ce sont précisément les fonctions où la PAO est supérieure au traitement de texte. Avec Word, le placement des images et la gestion des volets de pliage seraient difficiles, instables à la réimpression et imprécis pour l'imprimeur.

\textbf{Conclusion.} Jean a raison : le dépliant, document de mise en page complexe mêlant texte et images avec pliage, doit être réalisé avec un logiciel de PAO comme Publisher (ou Scribus/Canva), le traitement de texte restant réservé aux documents principalement rédactionnels.
\end{exoresolu}

\begin{savaistu}[Histoire \& Contexte Cameroun : De la composition à la PAO]
\begin{itemize}
\item \textbf{Années 80-90 (Cameroun) :} Les journaux (\textit{Cameroon Tribune}, \textit{La Nouvelle Expression}) et les imprimeries (SOPECAM, imprimeries privées à Douala/Yaoundé) utilisaient la \textbf{photocomposition} : texte tapé sur machine, tirage sur papier photosensible, découpage au cutter, collage manuel sur maquette (paste-up), flashage sur film, plaque offset. Lente, coûteuse, peu flexible.
\item \textbf{1985 :} Aldus \textbf{PageMaker} (Mac) + imprimante LaserWriter (PostScript) = naissance de la PAO moderne (« Desktop Publishing »). Révolution : « What You See Is What You Get » (WYSIWYG).
\item \textbf{Aujourd'hui au Cameroun :} La chaîne graphique est entièrement numérique (CtP : Computer-to-Plate). Les cybercafés, studios graphiques et services communication des administrations/entreprises utilisent la PAO (Publisher, CorelDraw, InDesign, Canva, Scribus). La maîtrise de la PAO est une compétence employable immédiate (création cartes de visite, affiches, faire-part, journaux scolaires, bulletins de notes).
\item \textbf{Standard d'échange :} Le \textbf{PDF (Portable Document Format)}, inventé par Adobe (1993), est devenu le format universel de diffusion et d'impression (PDF/X pour l'imprimerie professionnelle). Un document PAO se termine presque toujours par un \texttt{Exporter en PDF}.
\end{itemize}
\end{savaistu}

\begin{experience}[TP Guidé : Créer une affiche « Journée Portes Ouvertes » avec Scribus (ou Canva)]
\textbf{Objectif :} Mettre en œuvre la démarche de conception (5 étapes) pour produire une affiche A4 portrait.
\textbf{Matériel/Logiciels :} Postes Windows/Linux, \textbf{Scribus} (installé) \textbf{ou} navigateur web + compte \textbf{Canva} (gratuit). Dossier ressources : \texttt{logo\_ecole.svg}, \texttt{photo\_batiment.jpg}, \texttt{texte\_jpo.txt}.
\textbf{Durée :} 1 séance (50 min).

\begin{enumerate}
\item \textbf{Préparation (Infographie) -- 10 min :}
\begin{itemize}
\item Ouvrir \texttt{photo\_batiment.jpg} dans \textbf{GIMP}. Recadrer format 4:3, augmenter légèrement le contraste (\textit{Couleurs $\rightarrow$ Luminosité-Contraste}), exporter en \texttt{photo\_batiment\_ok.jpg} (JPEG qualité 90\%).
\item Ouvrir \texttt{logo\_ecole.svg} dans \textbf{Inkscape}. Vérifier qu'il est bien vectoriel (pas d'image incrustée). Enregistrer sous \texttt{logo\_ecole\_cmyk.svg} (pour impression).
\end{itemize}
\item \textbf{Nouveau document PAO -- 5 min :}
\begin{itemize}
\item Scribus : \textit{Fichier $\rightarrow$ Nouveau} $\rightarrow$ Document simple, Format A4, Orientation Portrait, Marges 15 mm, Unité mm. Canva : \textit{Créer un design $\rightarrow$ Affiche (A4)}.
\end{itemize}
\item \textbf{Gabarit \& Repères -- 5 min :}
\begin{itemize}
\item Scribus : \textit{Page $\rightarrow$ Gérer les gabarits} $\rightarrow$ Ajouter repères verticaux (marges + centre), horizontaux (tiers). Canva : \textit{Fichier $\rightarrow$ Afficher les marges / Repères}.
\end{itemize}
\item \textbf{Composition -- 20 min :}
\begin{enumerate}
\item \textbf{Bloc titre haut :} Insérer cadre texte (largeur page), saisir « JOURNÉES PORTES OUVERTES », Police \texttt{Anton} ou \texttt{Montserrat ExtraBold}, Taille 48 pt, Couleur \texttt{Bleu marine} (\#002B5C), Alignement centré.
\item \textbf{Image centrale :} Insérer cadre image (\textit{Fichier $\rightarrow$ Importer}), choisir \texttt{photo\_batiment\_ok.jpg}, ajuster \textit{Échelle au cadre} (proportions conservées).
\item \textbf{Bloc infos (Date, Lieu, Heure) :} Cadre texte sous l'image, police \texttt{Montserrat Regular} 18 pt, interligne 1.4, puces \texttt{$\bullet$}. Couleur \texttt{Gris foncé} (\#333333).
\item \textbf{Logo + Contact bas :} Insérer \texttt{logo\_ecole\_cmyk.svg} (cadre image vectoriel), à gauche. Cadre texte à droite : « Lycée de Biyem-Assi $\mid$ BP 123 Yaoundé $\mid$ www.lycee-biyem.cm », police 14 pt.
\item \textbf{Habillage :} Vérifier que le texte ne chevauche pas l'image (Scribus : \textit{Propriétés $\rightarrow$ Forme $\rightarrow$ Type de contour $\rightarrow$ Cadre de l'image}).
\end{enumerate}
\item \textbf{Vérification \& Export -- 10 min :}
\begin{itemize}
\item \textit{Fichier $\rightarrow$ Aperçu} (Scribus) ou \textit{Partager $\rightarrow$ Aperçu} (Canva) : vérifier marges, lisibilité, alignements (grille).
\item \textit{Fichier $\rightarrow$ Exporter $\rightarrow$ Enregistrer sous PDF} : Préréglage \textbf{Imprimeur (PDF/X-3 ou Haute qualité 300 ppp)}, \textbf{Repères de coupe} activés. Nommer \texttt{Affiche\_JPO\_2025.pdf}.
\item Ouvrir le PDF pour contrôle final.
\end{itemize}
\end{enumerate}
\textbf{Observations attendues :} L'image vectorielle (logo) reste nette au zoom 400\% dans le PDF ; l'image matricielle (photo) montre des pixels si on zoome au-delà de 100\% (résolution 300 ppp respectée). Le pliage n'est pas géré ici (affiche 1 page), mais les repères de coupe sont présents pour le massicotage.
\textbf{Compétences validées :} Distinction matriciel/vectoriel, chaîne Infographie $\rightarrow$ PAO, gabarit/marges/colonnes, styles texte, export PDF pro.
\end{experience}

\begin{attention}[Les 5 erreurs qui coûtent des points au BEPC]
\begin{enumerate}
\item \textbf{Confondre infographie et PAO.} L'infographie crée et retouche les \emph{images} ; la PAO compose la \emph{mise en page} d'un document complet. Dire « on fait de la PAO pour retoucher une photo » est faux.
\item \textbf{Confondre image matricielle et vectorielle.} Retenir : matricielle = pixels = photo = pixellisation à l'agrandissement ; vectorielle = objets géométriques = logo = agrandissement sans perte. Ne jamais proposer un JPEG pour un logo de banderole.
\item \textbf{Citer un logiciel sans le relier à la tâche.} « On utilise Publisher » ne suffit pas : il faut écrire « on utilise Publisher \emph{parce que} c'est un logiciel de PAO adapté à la mise en page ».
\item \textbf{Oublier des étapes dans la démarche de conception.} La démarche comporte 5 étapes : objectif/public, format, contenu, mise en page, vérification/diffusion. Oublier l'étape 1 (objectif et public) ou l'étape 5 (relecture et aperçu avant impression) fait perdre des points.
\item \textbf{Répondre sans justifier ni conclure.} Une réponse du type « oui » ou « c'est matriciel » sans définition, sans application au cas concret et sans phrase de conclusion ne rapporte qu'une fraction des points : toujours structurer en réponse $\rightarrow$ règle $\rightarrow$ application $\rightarrow$ conclusion.
\end{enumerate}
\end{attention}

\begin{exercice}[Entraînement BEPC -- Sujet type]
\textbf{Contexte :} L'association des anciens élèves du Lycée de Ngaoundéré organise ses retrouvailles. Le bureau doit produire : (1) une photo de groupe retouchée (supprimer un intrus, corriger la lumière) ; (2) le logo de l'association (création originale) ; (3) un programme détaillé sur 4 pages A5 pliées en deux (livret) ; (4) une grande banderole 3$\times$1 m pour l'entrée de la salle.
\textbf{Questions :}
\begin{enumerate}
\item Pour chacune des 4 productions, indique s'il s'agit d'une tâche d'infographie ou de PAO.
\item Pour chaque tâche, propose \textbf{un} logiciel gratuit adapté et justifie ton choix en une phrase.
\item Le logo sera sur la banderole (3 m) et sur le livret (3 cm). Quel type d'image (matriciel/vectoriel) et quel format de fichier choisis-tu ? Justifie techniquement.
\item Rédige la démarche complète (5 étapes justifiées) pour la réalisation du livret (production 3).
\end{enumerate}
\end{exercice}

\begin{corrige}[Exercice Entraînement BEPC]
\begin{enumerate}
\item \textbf{Production 1 (Photo retouchée) :} \cle{Infographie} (traitement image matricielle). \\
\textbf{Production 2 (Logo création) :} \cle{Infographie} (dessin vectoriel). \\
\textbf{Production 3 (Livret 4 pages) :} \cle{PAO} (mise en page document multipages). \\
\textbf{Production 4 (Banderole) :} \cle{PAO} (mise en page grand format, bien que souvent confiée à un imprimeur, la maquette se fait en PAO). \textit{Note : si la banderole ne contient que le logo agrandi, c'est de l'infographie vectorielle ; mais ici « banderole pour l'entrée » implique texte + logo + déco $\Rightarrow$ PAO.}
\item \textbf{Logiciels gratuits proposés :}
\begin{itemize}
\item Tâche 1 : \textbf{GIMP} — éditeur matriciel complet pour retouche photo (suppression, correction lumière).
\item Tâche 2 : \textbf{Inkscape} — éditeur vectoriel pour création logo original (courbes, texte, formes).
\item Tâche 3 : \textbf{Scribus} — logiciel PAO pro pour livret multipages, gestion gabarits, pliage, numérotation, export PDF/X.
\item Tâche 4 : \textbf{Scribus} (ou Canva) — PAO pour grand format, gestion repères de coupe/fonds perdus, export PDF haute résolution.
\end{itemize}
\item \textbf{Logo : Type \cle{Vectoriel}, Format \cle{SVG} (ou PDF vectoriel).} \\
\textbf{Justification :} Le logo doit passer de 3 cm (livret) à 300 cm (banderole), facteur 100. Une image matricielle (JPEG/PNG) de 3 cm à 300 ppp fait $\approx$ 350 px de large ; agrandie 100$\times$, elle ferait 35 000 px $\rightarrow$ fichier énorme ou pixellisation si résolution initiale faible. L'image vectorielle (SVG) est décrite par des équations géométriques : le RIP (Raster Image Processor) de l'imprimeur la rastérise \emph{à la résolution finale de l'imprimante grand format} (ex. 150 ppp), garantissant une netteté parfaite à toute taille. C'est le standard professionnel pour les logos.
\item \textbf{Démarche Livret (5 étapes) :}
\begin{enumerate}
\item \textbf{Objectif/Public :} Informer les anciens élèves du programme détaillé (conférences, repas, soirées). Public : adultes (30--60 ans), lecture assise $\Rightarrow$ police lisible (serif 11 pt), structure claire, numérotation impérative.
\item \textbf{Format/Support :} Livret 4 pages = 1 feuille A4 pliée en 2 (format fini A5, 14,8 $\times$ 21 cm). Impression recto-verso sur papier 135 g/m². Gabarit : marges 12 mm (intérieur 15 mm pour pliure), 2 colonnes par page.
\item \textbf{Contenu :} Récupération textes (Word $\rightarrow$ copier-coller sans format), photos retouchées sous GIMP (JPEG 300 ppp), logo vectoriel (SVG). Vérification orthographe \emph{avant} import.
\item \textbf{Mise en page (Scribus) :} Gabarit « Page intérieure » (en-tête « Retrouvailles 2025 », pied « p. N / 4 », numérotation auto). Flux de texte dans colonnes liées. Titres style « Titre 1 » (Montserrat Bold 14 pt, bleu marine). Images légendées, habillage carré. Vérification veuves/orphelines.
\item \textbf{Vérification/Diffusion :} Aperçu « Pages face à face » pour vérifier ordre des pages après pliage (imposition 2-up). Export PDF/X-1a (repères coupe, fonds perdus 3 mm). Envoi à l'imprimeur (clé USB/mail) + archivage dossier source (.sla + images + polices).
\end{enumerate}
\end{enumerate}
\end{corrige}

\newpage

\coursec{Exercices}

\courssub{Je vérifie mes connaissances}

\begin{exercice}[QCM : les bons réflexes de l'infographie]
Pour chaque question, entoure la (ou les) bonne(s) réponse(s).
\begin{enumerate}
\item Un \cle{pixel} est :
\begin{itemize}
\item[a)] une unité de mesure de la vitesse d'impression ;
\item[b)] le plus petit élément d'une image matricielle ;
\item[c)] un logiciel de retouche photo ;
\item[d)] un format de fichier vectoriel.
\end{itemize}
\item Parmi les logiciels suivants, lequel est un logiciel de \cle{PAO} ?
\begin{itemize}
\item[a)] Microsoft Word ; b) Microsoft Publisher ; c) Paint ; d) VLC.
\end{itemize}
\item Une image \cle{vectorielle} :
\begin{itemize}
\item[a)] se déforme (devient floue) quand on l'agrandit ;
\item[b)] est décrite par des formules mathématiques et peut être agrandie sans perte de qualité ;
\item[c)] est toujours une photographie ;
\item[d)] occupe toujours plus de mémoire qu'une image matricielle.
\end{itemize}
\item Le format de fichier typiquement \cle{matriciel} est :
\begin{itemize}
\item[a)] JPEG ; b) SVG ; c) AI ; d) EPS.
\end{itemize}
\end{enumerate}
\end{exercice}

\begin{exercice}[Vrai ou faux, avec justification]
Réponds par Vrai ou Faux, puis justifie ta réponse en une ou deux phrases.
\begin{enumerate}
\item La PAO (publication assistée par ordinateur) consiste à utiliser l'ordinateur pour concevoir des documents destinés à être imprimés ou diffusés : affiches, journaux, cartes, brochures.
\item Une photo prise avec un téléphone portable est une image vectorielle.
\item Dans une mise en page soignée, il est conseillé d'utiliser le plus grand nombre possible de polices de caractères et de couleurs pour attirer le lecteur.
\item La résolution d'une image, exprimée en dpi (points par pouce), influe sur la qualité de son impression.
\end{enumerate}
\end{exercice}

\begin{exercice}[Définitions à compléter]
Recopie et complète chaque définition avec les mots ou expressions suivants : \emph{infographie ; pixel ; image vectorielle ; mise en page ; PAO}.
\begin{enumerate}
\item La \ldots\ldots\ldots est l'ensemble des techniques de création et de traitement d'images et de documents à l'aide de l'ordinateur.
\item Une \ldots\ldots\ldots est une image décrite par des objets géométriques (lignes, courbes, formes) définis mathématiquement.
\item Le \ldots\ldots\ldots est le plus petit point coloré constituant une image matricielle.
\item La \ldots\ldots\ldots désigne l'organisation des éléments (textes, images, titres, marges) sur une page.
\item La \ldots\ldots\ldots est l'utilisation de logiciels spécialisés pour réaliser des documents à imprimer (journaux, affiches, cartes de visite).
\end{enumerate}
\end{exercice}

\begin{exercice}[Calcul direct : dimensions d'impression]
Une photo numérique a pour dimensions $600 \times 900$ pixels. On souhaite l'imprimer à la résolution de $300$ dpi (300 pixels par pouce). On rappelle que $1$ pouce $= 2,54$ cm.
\begin{enumerate}
\item Calcule la largeur et la hauteur de l'image imprimée en pouces.
\item Convertis ces dimensions en centimètres.
\end{enumerate}
\end{exercice}

\courssub{Je m'entraîne}

\begin{exercice}[Chaque besoin, son outil]
Pour chacun des besoins suivants, indique le \textbf{type d'image} (matricielle ou vectorielle) et le \textbf{logiciel} le plus adaptés, en justifiant ton choix en une phrase. Logiciels disponibles : Paint, Inkscape (dessin vectoriel), Microsoft Publisher, Photoshop (retouche photo).
\begin{enumerate}
\item Créer le logo du club informatique du collège, qui devra être imprimé en grand format sur une banderole.
\item Retoucher une photo de famille prise lors de la fête de l'école (luminosité, cadrage).
\item Mettre en page le journal scolaire de 8 pages (articles, photos, titres).
\item Dessiner un schéma simple de l'ordinateur pour une affiche de la salle multimédia.
\end{enumerate}
\end{exercice}

\begin{exercice}[Corriger une affiche ratée]
Le président du club des jeunes de Nkongsamba a conçu une affiche pour une journée portes ouvertes. Voici sa description :
\begin{itemize}
\item le texte utilise 6 polices de caractères différentes et 8 couleurs ;
\item le titre est écrit en petits caractères au bas de la page ;
\item une photo, étirée pour remplir la page, apparaît déformée et floue ;
\item aucune marge : le texte touche les bords du papier ;
\item la date et le lieu de l'événement sont absents.
\end{itemize}
\begin{enumerate}
\item Relève cinq erreurs de conception commises sur cette affiche.
\item Pour chaque erreur, propose une correction et justifie-la à l'aide des règles de mise en page vues en cours.
\end{enumerate}
\end{exercice}

\begin{exercice}[Maquette d'un bulletin scolaire]
On te fournit les éléments suivants pour réaliser le bulletin d'information du collège : un texte brut de 400 mots, le nom du collège et sa devise, deux photographies (une du proviseur, une de la nouvelle salle informatique) et la date de parution.
\begin{enumerate}
\item Schématise sur papier (ou à l'aide de formes dans un logiciel) la maquette de la première page : en-tête avec le nom du collège, titre, corps du texte en deux colonnes, emplacement des deux images, pied de page avec la date.
\item Justifie chacun de tes choix de placement (position du titre, taille des images, colonnes) en t'appuyant sur les règles d'équilibre et de lisibilité d'une mise en page.
\end{enumerate}
\end{exercice}

\begin{exercice}[Résolution et qualité d'impression]
Une photo de $2400 \times 1800$ pixels doit être imprimée à $300$ dpi. On rappelle : $1$ pouce $= 2,54$ cm.
\begin{enumerate}
\item Calcule les dimensions maximales d'impression en pouces, puis en centimètres.
\item Le club photo veut imprimer cette image sur une affiche de $60 \times 45$ cm à $300$ dpi. Est-ce possible sans perte de qualité ? Justifie ta réponse par un calcul.
\item Que se passera-t-il si l'on force l'impression en grand format ? Quel type d'image aurait fallu utiliser pour un logo à imprimer en très grand format ?
\end{enumerate}
\end{exercice}

\courssub{Je me prépare au BEPC}

\begin{exercice}[La carte d'invitation de la fête de la jeunesse]
Le collège bilingue de Bafoussam organise la fête de la jeunesse. Le préfet des études te charge de concevoir les cartes d'invitation destinées aux parents et aux autorités.
\begin{enumerate}
\item Définis les termes \cle{infographie} et \cle{PAO}, puis explique en quoi la conception de cette carte relève des deux.
\item Décris la démarche complète de conception de la carte d'invitation, en rédigeant et en justifiant chaque étape : analyse du besoin (informations indispensables à faire figurer), choix du logiciel (parmi Word, Publisher, Paint) avec deux arguments, choix du format et de l'orientation du papier, choix des couleurs et des polices (nombre maximal conseillé).
\item La carte doit contenir le blason du collège. Ce blason sera aussi imprimé en grand format sur une banderole. Quel type d'image (matricielle ou vectorielle) choisis-tu pour le blason ? Justifie.
\item Cite deux erreurs de mise en page à éviter absolument sur ce document officiel.
\end{enumerate}
\end{exercice}

\begin{exercice}[Le cybercafé « Eto'o Print » à Douala]
M. Kamga tient un cybercafé à Akwa. Une cliente lui apporte une photo de $3000 \times 2000$ pixels et demande une affiche publicitaire de $30 \times 20$ cm, imprimée à $300$ dpi ($1$ pouce $= 2,54$ cm).
\begin{enumerate}
\item Rappelle ce qu'est un pixel et ce qu'est la résolution d'une image.
\item Calcule les dimensions maximales d'impression de la photo à $300$ dpi, en pouces puis en centimètres (arrondis au dixième).
\item La photo convient-elle pour l'affiche demandée sans perte de qualité ? Justifie par la comparaison des dimensions.
\item La cliente veut aussi ajouter le logo de son entreprise sur l'affiche. M. Kamga hésite entre une image matricielle et une image vectorielle. Conseille-le en donnant deux avantages du type d'image choisi.
\item Propose à M. Kamga un logiciel de PAO adapté à la réalisation de l'affiche et un logiciel de retouche pour la photo, en justifiant chaque choix.
\end{enumerate}
\end{exercice}

\begin{exercice}[Le journal scolaire du collège de Maroua]
Le club presse du collège de Maroua veut publier son journal trimestriel de 4 pages, imprimé au format A4. Le journal contiendra : un éditorial, trois articles, cinq photos, les résultats sportifs sous forme de tableau et une banderole publicitaire pour la kermesse.
\begin{enumerate}
\item Définis la PAO et cite deux exemples de documents qu'elle permet de réaliser.
\item Parmi les logiciels suivants, lesquels sont adaptés à la mise en page du journal : Microsoft Publisher, Scribus, Paint, Microsoft Excel ? Justifie ta réponse.
\item Les cinq photos proviennent d'un téléphone portable : de quel type d'images s'agit-il ? Que risque-t-on si on les agrandit exagérément dans le journal ?
\item Pour la banderole publicitaire de la kermesse, qui sera imprimée en très grand format, quel type d'image faut-il privilégier pour le dessin du logo ? Justifie.
\item Rédige la démarche complète de conception du journal, de l'analyse du besoin jusqu'à l'impression, en justifiant au moins trois choix (logiciel, format, organisation de la page, nombre de polices et de couleurs).
\end{enumerate}
\end{exercice}

\newpage

\coursec{Corrigés des exercices}

\begin{corrige}[Exercice 1 — QCM : les bons réflexes de l'infographie]
\begin{enumerate}
\item \textbf{b)} Un \cle{pixel} (contraction de \emph{picture element}) est le plus petit élément constitutif d'une image \cle{matricielle} (ou bitmap). C'est un petit carré coloré ; l'ensemble de ces carrés forme l'image.
\item \textbf{b)} \cle{Microsoft Publisher} est un logiciel de \cle{PAO} (Publication Assistée par Ordinateur) dédié à la mise en page de documents complexes (affiches, journaux, cartes). Word est un traitement de texte, Paint un éditeur d'images matricielles basique, VLC un lecteur multimédia.
\item \textbf{b)} Une image \cle{vectorielle} est décrite par des objets géométriques (points, segments, courbes de Bézier) définis par des formules mathématiques. Elle est \emph{indépendante de la résolution} : on peut l'agrandir indéfiniment sans perte de qualité (pas de pixelisation).
\item \textbf{a)} \cle{JPEG} (ou JPG) est le format standard pour les images \cle{matricielles} (photos). SVG, AI (Illustrator) et EPS sont des formats principalement \cle{vectoriels}.
\end{enumerate}
\end{corrige}

\begin{corrige}[Exercice 2 — Vrai ou faux, avec justification]
\begin{enumerate}
\item \textbf{Vrai.} C'est la définition exacte de la \cle{PAO} : concevoir des documents (affiches, journaux, brochures, cartes de visite, etc.) destinés à l'impression professionnelle ou à la diffusion numérique (PDF), en gérant finement la mise en page, la typographie et les images.
\item \textbf{Faux.} Une photo prise avec un téléphone portable (ou tout appareil photo numérique) est une image \cle{matricielle} (matrice de pixels). Elle dépend de la résolution ; si on l'agrandit trop, elle devient floue (pixelisation). Une image vectorielle est un dessin géométrique (logo, schéma, police de caractères).
\item \textbf{Faux.} En mise en page, la règle d'or est la \emph{sobriété} et la \emph{cohérence}. On limite généralement à \textbf{2 ou 3 polices maximum} (ex : une pour les titres, une pour le corps de texte) et à une \textbf{palette de 3 à 4 couleurs harmonieuses}. Trop de polices/couleurs nuisent à la lisibilité et donnent un aspect amateur.
\item \textbf{Vrai.} La résolution (exprimée en \cle{dpi} — \emph{dots per inch} — ou ppp — \emph{points par pouce}) détermine le nombre de pixels imprimés par pouce. Une résolution de 300 dpi est le standard pour une impression de qualité (offset, jet d'encre photo). En dessous (ex: 72 dpi), l'impression sera pixelisée.
\end{enumerate}
\end{corrige}

\begin{corrige}[Exercice 3 — Définitions à compléter]
\begin{enumerate}
\item La \cle{infographie} est l'ensemble des techniques de création et de traitement d'images et de documents à l'aide de l'ordinateur.
\item Une \cle{image vectorielle} est une image décrite par des objets géométriques (lignes, courbes, formes) définis mathématiquement.
\item Le \cle{pixel} est le plus petit point coloré constituant une image matricielle.
\item La \cle{mise en page} désigne l'organisation des éléments (textes, images, titres, marges) sur une page.
\item La \cle{PAO} est l'utilisation de logiciels spécialisés pour réaliser des documents à imprimer (journaux, affiches, cartes de visite).
\end{enumerate}
\end{corrige}

\begin{corrige}[Exercice 4 — Calcul direct : dimensions d'impression]
Données : Largeur $L_{px} = 600$ px, Hauteur $H_{px} = 900$ px, Résolution $R = 300$ dpi, $1$ pouce $= 2,54$ cm.

\begin{enumerate}
\item \textbf{Dimensions en pouces :}
\[
L_{\text{in}} = \frac{L_{px}}{R} = \frac{600}{300} = 2 \text{ pouces}
\]
\[
H_{\text{in}} = \frac{H_{px}}{R} = \frac{900}{300} = 3 \text{ pouces}
\]
\item \textbf{Dimensions en centimètres :}
\[
L_{\text{cm}} = 2 \times 2,54 = 5,08 \text{ cm}
\]
\[
H_{\text{cm}} = 3 \times 2,54 = 7,62 \text{ cm}
\]
\end{enumerate}
L'image imprimée à 300 dpi fera \textbf{5,08 cm $\times$ 7,62 cm}.
\end{corrige}

\begin{corrige}[Exercice 5 — Chaque besoin, son outil]
\begin{enumerate}
\item \textbf{Logo du club (banderole grand format)} \\
\textbf{Type :} \cle{Vectorielle}. \\
\textbf{Logiciel :} \cle{Inkscape} (dessin vectoriel libre et gratuit). \\
\textbf{Justification :} Un logo doit pouvoir être agrandi à n'importe quelle taille (ici une banderole de plusieurs mètres) sans perte de qualité ni pixelisation ; seul le format vectoriel le garantit.
\item \textbf{Retouche photo de famille} \\
\textbf{Type :} \cle{Matricielle}. \\
\textbf{Logiciel :} \cle{Photoshop} (ou GIMP, alternative libre). \\
\textbf{Justification :} Une photo est par nature une image matricielle (pixels) ; la retouche (luminosité, cadrage, correction couleurs) agit sur les pixels, ce que font les logiciels de retouche photo.
\item \textbf{Journal scolaire 8 pages} \\
\textbf{Type :} Document \cle{PAO} (mélange textes/images). \\
\textbf{Logiciel :} \cle{Microsoft Publisher} (ou Scribus). \\
\textbf{Justification :} La PAO gère les documents multipages, les gabarits, les colonnes, les repères de pliage et le flux de texte entre cadres, ce que ne font pas Word (traitement de texte linéaire) ni Paint.
\item \textbf{Schéma ordinateur pour affiche} \\
\textbf{Type :} \cle{Vectorielle}. \\
\textbf{Logiciel :} \cle{Inkscape}. \\
\textbf{Justification :} Un schéma est composé de formes géométriques simples (rectangles, flèches, texte) ; le vectoriel assure des contours nets à toute taille et un fichier léger, idéal pour l'impression affiche.
\end{enumerate}
\end{corrige}

\begin{corrige}[Exercice 6 — Corriger une affiche ratée]
\begin{enumerate}
\item \textbf{Cinq erreurs relevées :}
\begin{enumerate}
\item Utilisation excessive de polices (6) et de couleurs (8).
\item Titre illisible : trop petit et mal placé (bas de page).
\item Photo déformée (étirée sans respecter les proportions) et floue (résolution insuffisante pour le format).
\item Absence de marges (texte collé aux bords).
\item Informations essentielles manquantes (date, lieu).
\end{enumerate}
\item \textbf{Corrections et justifications :}
\begin{enumerate}
\item \textbf{Correction :} Limiter à 2 polices (1 titre, 1 corps) et 3 couleurs max (charte graphique). \textbf{Justification :} Respect de l'harmonie visuelle et de la lisibilité ; trop de styles fatiguent l'œil et brouillent le message.
\item \textbf{Correction :} Placer le titre en haut de page, en très gros caractères (hiérarchie visuelle). \textbf{Justification :} Le titre est l'accroche ; il doit être vu en premier (règle de la lecture en Z ou en F, zone forte en haut à gauche/centre).
\item \textbf{Correction :} Redimensionner l'image en conservant ses proportions (verrouiller le rapport L/H) ; utiliser une image de résolution suffisante (300 dpi pour le format final). \textbf{Justification :} Éviter la distorsion (déformation) et le flou (pixelisation) ; respecter le droit à l'image et l'esthétique.
\item \textbf{Correction :} Appliquer des marges de sécurité (min 1 cm, idéalement 1,5–2 cm). \textbf{Justification :} Espace de respiration visuelle, zone de non-imprimabilité (marge technique de l'imprimante), confort de lecture.
\item \textbf{Correction :} Ajouter clairement : Date, Heure, Lieu, Contact. \textbf{Justification :} Une affiche d'événement est un message informatif ; sans ces données, elle perd sa fonction première.
\end{enumerate}
\end{enumerate}
\end{corrige}

\begin{corrige}[Exercice 7 — Maquette d'un bulletin scolaire]
\begin{enumerate}
\item \textbf{Description de la maquette (représentation schématique) :}
\begin{center}
\begin{tabular}{|p{0.94\linewidth}|}
\hline
\textbf{\large EN-TÊTE} \\
\hline
\textbf{NOM DU COLLÈGE} (Centré, Police Serif, Grande taille, Couleur principale) \\
\textit{Devise du collège} (Centré, Italique, Plus petit) \\
\hline
\hline
\textbf{\large TITRE DU BULLETIN} (ex: « Bulletin d'Information n°1 ») \\
\hline
\hline
\begin{tabular}{@{}p{0.48\linewidth}@{\hspace{0.04\linewidth}}p{0.48\linewidth}@{}}
\textbf{COLONNE GAUCHE} & \textbf{COLONNE DROITE} \\
\hline
\textbf{Éditorial / Article principal} & \textbf{Photo 1 : Le Proviseur} \\
(Texte justifié, police 10-11 pt, interligne 1,2) & (Légende dessous) \\
& \\
\textbf{Photo 2 : Salle Informatique} & \textbf{Autres brèves / Actualités} \\
(Légende dessous) & (Texte justifié) \\
\end{tabular} \\
\hline
\hline
\textbf{\large PIED DE PAGE} \\
\hline
Date de parution \hfill Numéro de page \hfill Contact / Site web \\
\hline
\end{tabular}
\end{center}
\item \textbf{Justifications des choix :}
\begin{itemize}
\item \textbf{En-tête centré avec nom + devise :} Zone de repérage immédiate (identité visuelle), ancrage en haut (zone forte).
\item \textbf{Titre distinctif :} Hiérarchie de l'information : Nom du collège $>$ Titre du bulletin $>$ Articles.
\item \textbf{Deux colonnes :} Standard en presse écrite ; améliore la lisibilité (longueur de ligne idéale $\approx$ 60-70 caractères), dynamise la page, permet d'alterner texte/images.
\item \textbf{Placement des images :} Une photo par colonne (équilibre visuel). Photo du proviseur près de l'éditorial (lien sémantique). Photo salle info dans l'autre colonne (équilibre masse noire/blanche). Légendes obligatoires.
\item \textbf{Pied de page :} Informations de service (date, pagination) discrètes mais accessibles, sans perturber la lecture.
\item \textbf{Marges et gouttière :} Marges extérieures confortables (1,5-2 cm), gouttière (espace entre colonnes) $\ge$ 0,5 cm pour éviter la confusion de lecture.
\end{itemize}
\end{enumerate}
\end{corrige}

\begin{corrige}[Exercice 8 — Résolution et qualité d'impression]
Données : $L_{px} = 2400$, $H_{px} = 1800$, $R = 300$ dpi, $1 \text{ pouce} = 2,54 \text{ cm}$.

\begin{enumerate}
\item \textbf{Dimensions maximales à 300 dpi :}
\[
L_{\text{in}} = \frac{2400}{300} = 8 \text{ pouces} \quad\Rightarrow\quad L_{\text{cm}} = 8 \times 2,54 = 20,32 \text{ cm}
\]
\[
H_{\text{in}} = \frac{1800}{300} = 6 \text{ pouces} \quad\Rightarrow\quad H_{\text{cm}} = 6 \times 2,54 = 15,24 \text{ cm}
\]
\textbf{Format max : 20,3 cm $\times$ 15,2 cm (environ format A5).}

\item \textbf{Projet affiche 60 $\times$ 45 cm à 300 dpi :}
Calcul des pixels nécessaires :
\[
L_{\text{req}} = \frac{60}{2,54} \times 300 \approx 23,62 \times 300 = 7086 \text{ px}
\]
\[
H_{\text{req}} = \frac{45}{2,54} \times 300 \approx 17,72 \times 300 = 5315 \text{ px}
\]
Pixels disponibles : $2400 \times 1800$.
\textbf{Non, ce n'est pas possible.} L'image ne contient que $\approx 34\%$ des pixels requis en largeur et hauteur. L'agrandissement imposerait une interpolation massive $\rightarrow$ flou, artefacts, perte de netteté inacceptable.

\item \textbf{Conséquences et alternative :}
\begin{itemize}
\item \textbf{Si on force l'impression :} L'image sera \textbf{pixelisée} (effet « escalier » sur les contours), \textbf{floue} (perte de détails), de qualité inacceptable pour une affiche vue de près.
\item \textbf{Type d'image pour un logo grand format :} \cle{Image vectorielle}. Elle est définie par des équations, donc \textbf{redimensionnable à l'infini sans aucune perte de qualité}, contours nets, fichier léger. C'est le standard professionnel pour logos, pictogrammes, textes.
\end{itemize}
\end{enumerate}
\end{corrige}

\begin{corrige}[Exercice 9 — La carte d'invitation de la fête de la jeunesse (Prépa BEPC)]
\textbf{Barème indicatif : 10 points} — \textit{Points de vigilance : définitions précises, démarche structurée et justifiée, choix technique argumenté (vectoriel pour logo), erreurs de mise en page citées.}

\begin{enumerate}
\item \textbf{Définitions et lien :}
\begin{itemize}
\item \cle{Infographie} : Ensemble des techniques de création, de manipulation et de traitement d'images numériques (dessins, photos, graphiques) à l'aide d'un ordinateur.
\item \cle{PAO} (Publication Assistée par Ordinateur) : Utilisation de logiciels spécialisés pour composer des documents destinés à l'impression (mise en page, typographie, gestion des couleurs, export PDF HD).
\item \textbf{Lien :} La carte d'invitation nécessite de \textbf{créer/retravailler les éléments graphiques} (blason, fonds, mise en forme du texte) $\rightarrow$ \emph{Infographie}, puis de les \textbf{assembler en page} (positionnement, gabarit, repères de coupe, gestion quadrichromie) pour l'impression $\rightarrow$ \emph{PAO}. Les deux sont indissociables dans le flux de production.
\end{itemize}

\item \textbf{Démarche complète de conception :}
\begin{enumerate}
\item \textbf{Analyse du besoin (Informations indispensables) :}
\begin{itemize}
\item Nom de l'événement : « Fête de la Jeunesse ».
\item Organisateur : « Collège Bilingue de Bafoussam ».
\item Date, Heure de début/fin, Lieu précis (enceinte du collège / salle des fêtes).
\item Public cible : Parents, Autorités (Préfet, Maire, DRESEB...).
\item Mentions : RSVP (Répondez s'il vous plaît), Contact téléphone/email, Blason officiel.
\item Ton : Solennel, festif, institutionnel.
\end{itemize}
\item \textbf{Choix du logiciel : \cle{Microsoft Publisher} (ou Scribus).}
\textbf{Arguments :}
\begin{enumerate}
\item \textbf{Gabarits « Carte de visite / Invitation » intégrés} avec repères de pliage (rainage), marges de sécurité, fonds perdus — gain de temps et conformité imprimeur.
\item \textbf{Gestion avancée des blocs} (texte, images, formes) indépendants, alignement magnétique, calques, export PDF/X-1a prêt pour l'imprimerie (quadrichromie, 300 dpi, fonds perdus).
\end{enumerate}
\emph{Pourquoi pas Word ?} Mise en page libre mais fastidieuse pour du pliage précis, gestion des fonds perdus approximative. \emph{Pourquoi pas Paint ?} Uniquement matriciel, pas de gestion texte vectoriel, pas de multipage, pas d'export print.
\item \textbf{Choix du format et de l'orientation :}
\begin{itemize}
\item \textbf{Format :} A6 (10,5 $\times$ 14,8 cm) ou 12 $\times$ 17 cm (standard carte invitation, économique à l'impression, tient dans enveloppe C6/DL).
\item \textbf{Orientation :} \textbf{Paysage} (souvent plus esthétique pour invitation, permet grande photo fond + texte superposé) ou \textbf{Portrait} (classique, formel). Choix justifié par la maquette.
\end{itemize}
\item \textbf{Choix des couleurs et polices :}
\begin{itemize}
\item \textbf{Couleurs :} 2 à 3 max (ex : Bleu marine/Or du blason + Blanc + Gris foncé). Respect de la charte graphique du collège. Contraste élevé texte/fond pour lisibilité.
\item \textbf{Polices :} \textbf{2 maximum}. Ex : \texttt{Georgia} ou \texttt{Garamond} (serif, solennel) pour le corps ; \texttt{Montserrat} ou \texttt{Playfair Display} (titre/Blason). Interdiction : Comic Sans, polices fantaisies illisibles.
\end{itemize}
\end{enumerate}

\item \textbf{Blason du collège : Type \cle{Vectoriel}.}
\textbf{Justification :} Le blason sera imprimé \textbf{petit sur la carte} (quelques cm) et \textbf{très grand sur la banderole} (plusieurs mètres). Seul le vectoriel garantit une netteté parfaite à ces deux échelles extrêmes sans gérer deux fichiers sources distincts. Un fichier matriciel devrait faire des milliers de pixels pour la banderole $\rightarrow$ fichier énorme, lourd à manipuler.

\item \textbf{Deux erreurs de mise en page à éviter absolument :}
\begin{enumerate}
\item \textbf{Fautes d'orthographe/grammaire/syntaxes} sur un document officiel $\rightarrow$ discrédite l'institution. (Rellecture croisée obligatoire).
\item \textbf{Lisibilité compromise :} Contraste insuffisant (texte clair sur fond clair), police trop petite ($< 9$ pt pour corps), texte sur zone d'image trop chargée, marges trop justes (risque rognage infos lors de la coupe).
\end{enumerate}
\end{enumerate}
\end{corrige}

\begin{corrige}[Exercice 10 — Le cybercafé « Eto'o Print » à Douala]
\textbf{Barème indicatif : 12 points} — \textit{Points de vigilance : définitions précises (pixel, résolution), calculs justifiés avec unités, comparaison dimensions réelles vs demandées, argumentaire vectoriel vs matriciel, choix logiciels adaptés.}

\begin{enumerate}
\item \textbf{Définitions :}
\begin{itemize}
\item \cle{Pixel} : Plus petit élément adressable d'une image numérique matricielle ; petit carré (ou rectangle) de couleur uniforme. L'image est une grille (matrice) de pixels.
\item \cle{Résolution} (d'un fichier image pour l'impression) : Nombre de pixels par unité de longueur (généralement \cle{dpi} — \emph{dots per inch} — ou ppp). Elle relie la taille en pixels du fichier à la taille physique du tirage. Ex : 300 dpi = 300 pixels sur 1 pouce (2,54 cm).
\end{itemize}

\item \textbf{Calcul dimensions maximales (photo 3000 $\times$ 2000 px à 300 dpi) :}
\[
L_{\text{in}} = \frac{3000}{300} = 10 \text{ pouces} \quad\Rightarrow\quad L_{\text{cm}} = 10 \times 2,54 = \textbf{25,4 cm}
\]
\[
H_{\text{in}} = \frac{2000}{300} = 6,\overline{6} \text{ pouces} \quad\Rightarrow\quad H_{\text{cm}} = 6,\overline{6} \times 2,54 \approx \textbf{16,9 cm} \text{ (arrondi au dixième)}
\]
\textbf{Format maximal sans perte : 25,4 cm $\times$ 16,9 cm.}

\item \textbf{Adéquation pour l'affiche 30 $\times$ 20 cm :}
\begin{itemize}
\item Demandé : 30 cm (L) $\times$ 20 cm (H).
\item Possible max : 25,4 cm (L) $\times$ 16,9 cm (H).
\item \textbf{Non, la photo ne convient pas.} Les dimensions demandées dépassent les dimensions maximales de $\approx$ 4,6 cm en largeur et 3,1 cm en hauteur. Pour atteindre 30 $\times$ 20 cm à 300 dpi, il faudrait un fichier d'au moins :
\[
\frac{30}{2,54} \times 300 \approx 3543 \text{ px (L)} \quad\text{et}\quad \frac{20}{2,54} \times 300 \approx 2362 \text{ px (H)}
\]
Le fichier fourni (3000 $\times$ 2000) est insuffisant $\rightarrow$ interpolation forcée $\rightarrow$ perte de qualité visible.
\end{itemize}

\item \textbf{Conseil pour le logo : Image \cle{Vectorielle}.}
\textbf{Deux avantages décisifs :}
\begin{enumerate}
\item \textbf{Redimensionnement sans perte infinie :} Le logo restera net sur l'affiche 30$\times$20 cm comme sur une banderole 3$\times$2 m ou un favicon 16$\times$16 px. Un fichier matriciel unique ne peut pas couvrir cette plage sans perte.
\item \textbf{Contours nets et édition aisée :} Les courbes de Bézier restent lisses ; on peut changer les couleurs (Pantone/CMJN) facilement pour l'imprimeur ; fichier source léger.
\end{enumerate}

\item \textbf{Logiciels recommandés à M. Kamga :}
\begin{itemize}
\item \textbf{PAO pour l'affiche :} \cle{Microsoft Publisher} (ou \cle{Scribus} gratuit).
\textbf{Justification :} Gestion professionnelle des blocs texte/image, repères de coupe/fonds perdus, gabarits, export PDF haute définition (PDF/X) standard imprimeur, travail en CMJN.
\item \textbf{Retouche pour la photo :} \cle{Adobe Photoshop} (ou \cle{GIMP} gratuit).
\textbf{Justification :} Outils experts de correction (niveaux, courbes, balance des blancs, netteté, suppression défauts), travail par calques non destructif, gestion profils ICC, redimensionnement par interpolation bicubique de haute qualité si léger agrandissement nécessaire.
\end{itemize}
\end{enumerate}
\end{corrige}

\begin{corrige}[Exercice 11 — Le journal scolaire du collège de Maroua]
\textbf{Barème indicatif : 14 points} — \textit{Points de vigilance : définition PAO + exemples, tri logiciels justifié, nature photos téléphone + risque agrandissement, choix vectoriel banderole, démarche complète structurée (analyse $\rightarrow$ impression) avec 3 justifications minimales.}

\begin{enumerate}
\item \textbf{PAO et exemples :}
\cle{PAO} (Publication Assistée par Ordinateur) : Utilisation de logiciels dédiés (Publisher, Scribus, InDesign) pour composer, mettre en page et préparer des documents destinés à l'impression professionnelle ou à la diffusion numérique (PDF), en gérant la typographie, les images, les couleurs (CMJN), les gabarits et les repères techniques.
\textbf{Exemples de documents :} Journal/Newsletter, Affiche, Brochure dépliant, Carte de visite, Bulletin scolaire, Livre, Magazine.

\item \textbf{Logiciels adaptés pour la mise en page du journal :}
\begin{itemize}
\item \textbf{Microsoft Publisher} : \textbf{OUI}. Logiciel PAO grand public, gère multipages, colonnes, gabarits, fonds perdus, export PDF print.
\item \textbf{Scribus} : \textbf{OUI}. Logiciel PAO libre/gratuit (Open Source), standard dans l'éducation, fonctions équivalentes Publisher/InDesign, multiplateforme.
\item \textbf{Paint} : \textbf{NON}. Éditeur d'images \cle{matricielles} (raster) monopage. Pas de gestion de texte fluide, pas de colonnes, pas de gabarits, pas de multipage, export bitmap uniquement.
\item \textbf{Microsoft Excel} : \textbf{NON}. Tableur (calcul, tableaux). Peut imprimer des grilles mais pas mettre en page un journal (flux texte, colonnes liées, habillage images, typographie fine).
\end{itemize}

\item \textbf{Photos du téléphone :}
\begin{itemize}
\item \textbf{Type :} Images \cle{matricielles} (raster / bitmap). Formats typiques : JPEG, HEIC, PNG. Composées de pixels.
\item \textbf{Risque d'agrandissement excessif :} \textbf{Pixelisation} (effet « blocs » visibles), \textbf{flou} (perte de détails), \textbf{artefacts de compression} (si JPEG très compressé). La résolution native (ex: 4000 px) fixe la taille max d'impression nette (ex: $\approx$ 34 cm à 300 dpi). Au-delà, l'interpolation dégrade l'image.
\end{itemize}

\item \textbf{Banderole kermesse — Logo : Type \cle{Vectoriel}.}
\textbf{Justification :} Impression « très grand format » (plusieurs mètres). Une image matricielle nécessiterait une résolution native énorme (dizaines de milliers de pixels) $\rightarrow$ fichier ingérable, lourd, lent. L'image vectorielle (SVG, AI, EPS, PDF vectoriel) est \textbf{indépendante de la résolution} : un seul fichier source, contours mathématiquement parfaits à n'importe quelle taille, couleurs quadrichromie (CMJN) gérées proprement pour l'imprimeur.

\item \textbf{Démarche complète de conception du journal (4 pages A4) :}
\begin{enumerate}
\item \textbf{Analyse du besoin \& Chemin de fer :} Lister contenus (Édito, 3 articles, 5 photos, Tableau sport, Pub kermesse), définir pagination (P1: Une + Édito ; P2-3: Articles + Photos ; P4: Sport + Pub), public (élèves, parents, profs), ton (dynamique, lisible).
\item \textbf{Choix du logiciel : \cle{Scribus} (ou Publisher).}
\textbf{Justification 1 (Logiciel) :} Gratuit/Libre (budget club), standard éducatif Cameroun, gère nativement multipages, colonnes, styles de paragraphes, couleurs CMJN, export PDF/X-1a prêt pour l'imprimeur local.
\item \textbf{Création du gabarit (Master Pages) :} Marges (1,5 cm), gouttières, numérotation auto, en-tête/pied de page (nom journal, date, n° page), grille de mise en page (ex: 3 colonnes).
\textbf{Justification 2 (Organisation) :} Grille 3 colonnes = modularité (texte 1 col, photo 2 cols, titre 3 cols), régularité visuelle entre pages, gain de temps.
\item \textbf{Préparation des contenus :} Relecture/correction textes (Word $\rightarrow$ import), retouche photos (GIMP : recadrage, niveaux, 300 dpi, CMJN), création tableau sport (Scribus/Calc $\rightarrow$ import), \textbf{logo kermesse en vectoriel (Inkscape)}.
\item \textbf{Mise en page (Composition) :} Application styles (Titre, Intertitre, Corps, Légende), habillage images, alignement à la grille, respect hiérarchie visuelle.
\textbf{Justification 3 (Polices/Couleurs) :} \textbf{Max 2 familles de polices} (ex: \texttt{Roboto Slab} titres, \texttt{Roboto} corps) + \textbf{Charte 3 couleurs} (identité collège) = lisibilité, unité graphique, professionnalisme.
\item \textbf{Épreuvage (Preflight) :} Vérification : images 300 dpi ? CMJN ? Polices vectorisées/embarquées ? Fonds perdus (3 mm) ? Repères de coupe ? Orthographe ? Liens images valides ?
\item \textbf{Export \& Impression :} Export \textbf{PDF/X-1a:2001} (standard imprimerie), envoi à l'imprimeur (clé USB/Email/WeTransfer), contrôle du BAT (Bon À Tirer) avant tirage final.
\end{enumerate}
\end{enumerate}
\end{corrige}

\newpage

\coursec{Fiche bilan}

\courssub{Carte mentale de la leçon}

\begin{popfigure}
\begin{center}\tcbox[colback=popdark,colframe=popdark,coltext=white,arc=9pt,boxsep=4pt,left=6pt,right=6pt]{\bfseries\small Infographie et PAO}\end{center}
\vspace{-2pt}
\begin{tcbraster}[raster columns=3,raster equal height=rows,raster column skip=6pt,raster row skip=6pt,enhanced,blankest]
\begin{tcolorbox}[enhanced,colback=white,colframe=popblue,boxrule=0.9pt,arc=3pt,title={Définitions},fonttitle=\bfseries\scriptsize,coltitle=white,colbacktitle=popblue,left=4pt,right=4pt,top=2pt,bottom=2pt,boxsep=1pt,toptitle=1.5pt,bottomtitle=1.5pt,halign=left]
\begin{itemize}[nosep,leftmargin=*,itemsep=1pt,font=\scriptsize]
\item Infographie
\item PAO
\end{itemize}
\end{tcolorbox}
\begin{tcolorbox}[enhanced,colback=white,colframe=poporange,boxrule=0.9pt,arc=3pt,title={Images numériques},fonttitle=\bfseries\scriptsize,coltitle=white,colbacktitle=poporange,left=4pt,right=4pt,top=2pt,bottom=2pt,boxsep=1pt,toptitle=1.5pt,bottomtitle=1.5pt,halign=left]
\begin{itemize}[nosep,leftmargin=*,itemsep=1pt,font=\scriptsize]
\item Matricielles
\item Vectorielles
\end{itemize}
\end{tcolorbox}
\begin{tcolorbox}[enhanced,colback=white,colframe=popgreen,boxrule=0.9pt,arc=3pt,title={Logiciels},fonttitle=\bfseries\scriptsize,coltitle=white,colbacktitle=popgreen,left=4pt,right=4pt,top=2pt,bottom=2pt,boxsep=1pt,toptitle=1.5pt,bottomtitle=1.5pt,halign=left]
\begin{itemize}[nosep,leftmargin=*,itemsep=1pt,font=\scriptsize]
\item Paint, GIMP
\item Publisher, Scribus
\end{itemize}
\end{tcolorbox}
\begin{tcolorbox}[enhanced,colback=white,colframe=poppurple,boxrule=0.9pt,arc=3pt,title={Mise en page},fonttitle=\bfseries\scriptsize,coltitle=white,colbacktitle=poppurple,left=4pt,right=4pt,top=2pt,bottom=2pt,boxsep=1pt,toptitle=1.5pt,bottomtitle=1.5pt,halign=left]
\begin{itemize}[nosep,leftmargin=*,itemsep=1pt,font=\scriptsize]
\item Texte, images
\item Marges, colonnes
\end{itemize}
\end{tcolorbox}
\begin{tcolorbox}[enhanced,colback=white,colframe=poppink,boxrule=0.9pt,arc=3pt,title={Démarche de conception},fonttitle=\bfseries\scriptsize,coltitle=white,colbacktitle=poppink,left=4pt,right=4pt,top=2pt,bottom=2pt,boxsep=1pt,toptitle=1.5pt,bottomtitle=1.5pt,halign=left]
\begin{itemize}[nosep,leftmargin=*,itemsep=1pt,font=\scriptsize]
\item Analyser le besoin
\item Réaliser, vérifier
\end{itemize}
\end{tcolorbox}
\begin{tcolorbox}[enhanced,colback=white,colframe=popgold,boxrule=0.9pt,arc=3pt,title={Applications},fonttitle=\bfseries\scriptsize,coltitle=white,colbacktitle=popgold,left=4pt,right=4pt,top=2pt,bottom=2pt,boxsep=1pt,toptitle=1.5pt,bottomtitle=1.5pt,halign=left]
\begin{itemize}[nosep,leftmargin=*,itemsep=1pt,font=\scriptsize]
\item Affiches, journaux
\item Cartes, dépliants
\end{itemize}
\end{tcolorbox}
\end{tcbraster}

\legende{Carte mentale récapitulative : l'initiation à la publication assistée par ordinateur.}
\end{popfigure}

\begin{bilan}
\textbf{Définitions clés}
\begin{itemize}
  \item \cle{Infographie} : ensemble des techniques de création et de traitement d'images à l'aide de l'ordinateur.
  \item \cle{PAO} (Publication Assistée par Ordinateur) : utilisation de l'ordinateur et de logiciels spécialisés pour concevoir des documents destinés à être imprimés ou diffusés (affiches, journaux, dépliants, cartes de visite).
  \item \cle{Image matricielle} (ou bitmap) : image formée de petits points appelés \cle{pixels} ; elle perd en qualité quand on l'agrandit.
  \item \cle{Image vectorielle} : image décrite par des formes géométriques (lignes, courbes) ; elle peut être agrandie sans perte de qualité.
  \item \cle{Pixel} : plus petit élément d'une image matricielle, caractérisé par une couleur.
  \item \cle{Résolution} : nombre de pixels par unité de longueur, exprimée en ppp (points par pouce).
\end{itemize}

\textbf{Comparatif à connaître par c\oe ur}
\begin{center}
\begin{tabularx}{\linewidth}{|l|X|X|}
\hline
\rowcolor{popblueL} \textbf{Critère} & \textbf{Image matricielle} & \textbf{Image vectorielle} \\
\hline
Composition & Pixels (points colorés) & Objets géométriques \\
\hline
Agrandissement & Perte de qualité (flou) & Aucune perte de qualité \\
\hline
Usage typique & Photos, images scannées & Logos, schémas, plans \\
\hline
Exemples de formats & JPEG, PNG, BMP & SVG, WMF \\
\hline
\end{tabularx}
\end{center}

\textbf{Logiciels et éléments de mise en page}
\begin{itemize}
  \item Logiciels d'infographie : \texttt{Paint}, \texttt{GIMP} (images matricielles) ; \texttt{Inkscape}, module Dessin (images vectorielles).
  \item Logiciels de PAO : \texttt{Microsoft Publisher}, \texttt{Scribus}.
  \item Éléments d'une mise en page : les \cle{marges}, le \cle{format} du papier, l'\cle{orientation} (portrait ou paysage), les \cle{colonnes}, le \cle{texte} (police, taille), les \cle{images} et les \cle{titres}.
\end{itemize}

\textbf{Méthode express : concevoir un document de PAO}
\begin{enumerate}
  \item Analyser le besoin : quel document ? pour quel public ? quel message ?
  \item Rassembler les éléments : textes, images, logo.
  \item Choisir le format, l'orientation et les marges.
  \item Réaliser la mise en page avec le logiciel de PAO.
  \item Vérifier (orthographe, alignement, qualité des images) puis imprimer ou exporter.
\end{enumerate}
\end{bilan}

\begin{aretenir}[Checklist avant le BEPC]
\begin{itemize}
  \item $\square$ Je sais définir l'infographie et la PAO avec mes propres mots.
  \item $\square$ Je sais distinguer une image matricielle d'une image vectorielle.
  \item $\square$ Je connais la signification des mots \cle{pixel} et \cle{résolution}.
  \item $\square$ Je sais expliquer pourquoi une photo agrandie devient floue.
  \item $\square$ Je sais citer au moins deux logiciels d'infographie et deux logiciels de PAO.
  \item $\square$ Je connais les principaux éléments d'une mise en page (marges, orientation, colonnes, texte, images).
  \item $\square$ Je sais différencier les orientations portrait et paysage.
  \item $\square$ Je sais décrire les étapes de conception d'un document de PAO.
  \item $\square$ Je sais choisir le bon logiciel selon le document à réaliser (affiche, carte de visite, journal scolaire).
  \item $\square$ Je sais rédiger et justifier ma démarche de conception face à une situation concrète de la vie courante.
\end{itemize}
\end{aretenir}

\findecours

\end{document}
